%=============================================================================
%  ESSAI — MAÎTRISE EN INFORMATIQUE — UQAC
%  Gabarit LaTeX reproduisant la mise en forme de l'essai de référence
%
%  Sujet   : Intégration du référentiel MITRE ATT&CK dans un écosystème MCP
%  Auteur  : Kamil Saleh
%
%  Arborescence :
%     main.tex
%     chapitres/chapitre1.tex, chapitre2.tex, ...
%     figures/  (images et figures TikZ)
%     references.bib, revue_litterature.bib
%     logo-ia-assistance-limitee.png
%
%  Compilation :  latexmk -pdf main.tex
%                 (ou pdflatex → bibtex → pdflatex → pdflatex)
%=============================================================================
\documentclass[12pt,letterpaper,twoside]{report}

%--------------------------- Encodage et langue ------------------------------
\usepackage[utf8]{inputenc}
\usepackage[T1]{fontenc}
\usepackage[numbers,sort&compress]{natbib}   % à charger avant babel
\usepackage[french]{babel}
\usepackage{mathptmx}            % police Times (comme le modèle UQAC)

%------------------------------ Mise en page ---------------------------------
\usepackage[letterpaper,inner=3.5cm,outer=2.5cm,top=2.5cm,bottom=2.5cm]{geometry}
\usepackage{setspace}
\onehalfspacing                  % remplacer par \doublespacing si exigé
\usepackage{fancyhdr}
\usepackage{titlesec}
\usepackage{graphicx}
\usepackage{xcolor}
\usepackage{caption}

%--------------------------- Tableaux et figures -----------------------------
\usepackage{array}
\usepackage{calc}
\usepackage{multirow}
\usepackage{longtable}
\usepackage{booktabs}
\usepackage{float}
\usepackage{placeins}
\usepackage{amsmath,amssymb}
\usepackage{tikz}
\usetikzlibrary{positioning,arrows.meta,calc,fit,backgrounds}
\usepackage[most]{tcolorbox}
\usepackage{listings}
\usepackage{etoolbox}
\AtBeginEnvironment{verbatim}{\singlespacing}   % blocs de code compacts

\usepackage[hidelinks]{hyperref}  % toujours chargé en dernier

\graphicspath{{./}{images/}{figures/}}

% Styles communs des figures TikZ (couleurs, boîtes, flèches)
% =====================================================================
%  Styles partagés par toutes les figures TikZ de l'essai.
%  Chargé une seule fois depuis le préambule de main.tex.
%  Modifier une couleur ici la modifie dans toutes les figures.
% =====================================================================

\definecolor{figbleu}{RGB}{31,95,168}
\definecolor{figvert}{RGB}{26,122,76}
\definecolor{figor}{RGB}{154,123,23}
\definecolor{figrouge}{RGB}{192,57,43}
\definecolor{figgris}{RGB}{120,120,120}
\definecolor{figbordeaux}{RGB}{155,27,74}
\definecolor{figpeche}{RGB}{252,229,205}
\definecolor{figjaune}{RGB}{252,229,110}

\tikzset{
  % --- boîtes de base -------------------------------------------------
  fbase/.style   = {rounded corners=3pt, line width=0.7pt, align=center,
                    inner sep=5pt, font=\footnotesize},
  fbleu/.style   = {fbase, draw=figbleu,      fill=figbleu!6},
  fvert/.style   = {fbase, draw=figvert,      fill=figvert!6},
  for/.style     = {fbase, draw=figor,        fill=figor!12},
  frouge/.style  = {fbase, draw=figrouge,     fill=figrouge!5},
  fgris/.style   = {fbase, draw=figgris,      fill=black!4},
  fbord/.style   = {fbase, draw=figbordeaux,  fill=figbordeaux!4},
  fblanc/.style  = {fbase, draw=black!35,     fill=white},
  fpointille/.style = {fbase, draw=black!45, fill=black!4, dashed},
  % --- variantes alignées à gauche ------------------------------------
  fgauche/.style = {align=left, text width=3.3cm},
  % --- flèches --------------------------------------------------------
  fleche/.style  = {-{Stealth[length=5pt,width=4pt]}, line width=0.8pt,
                    draw=black!55},
  fleche2/.style = {{Stealth[length=5pt,width=4pt]}-{Stealth[length=5pt,width=4pt]},
                    line width=0.8pt, draw=black!55},
  fgrosse/.style = {-{Triangle[length=6pt,width=6pt]}, line width=1pt,
                    draw=black!60},
  fpale/.style   = {-{Stealth[length=5pt,width=4pt]}, line width=0.8pt,
                    draw=black!30},
  % --- étiquettes de flèche -------------------------------------------
  etiq/.style    = {font=\scriptsize, inner sep=2pt, fill=white,
                    text=black!65},
  etiqi/.style   = {etiq, font=\scriptsize\itshape},
  etiqt/.style   = {etiq, font=\scriptsize\ttfamily},
}

% Styles propres à l'annexe (schémas d'exécution et blocs de code)
% =====================================================================
%  Styles propres à l'annexe (schémas d'exécution et blocs de code).
%  Les styles de flèches y sont préfixés « a » (afleche, aflechep, aetiq)
%  pour ne pas écraser ceux de tikz-commun.tex utilisés par les chapitres.
% =====================================================================

% ---------- Palette ----------
\definecolor{cAtk}{HTML}{C0392B}\definecolor{cTar}{HTML}{2C3E50}
\definecolor{cZeek}{HTML}{2980B9}\definecolor{cGT}{HTML}{27AE60}
\definecolor{cScript}{HTML}{8E44AD}\definecolor{cOut}{HTML}{E67E22}
\definecolor{cBg}{HTML}{F4F6F7}\definecolor{cCode}{HTML}{ECF0F1}
\definecolor{cComment}{HTML}{7F8C8D}\definecolor{cStr}{HTML}{C0392B}

% ---------- Styles TikZ ----------
\tikzset{
  boite/.style={rounded corners=2pt, draw, thick, align=center,
                inner sep=3.5pt, font=\footnotesize, minimum height=0.75cm,
                fill=white},
  atk/.style   ={boite, draw=cAtk,   text=cAtk,   fill=cAtk!6},
  tar/.style   ={boite, draw=cTar,   text=cTar,   fill=cTar!6},
  zeek/.style  ={boite, draw=cZeek,  text=cZeek,  fill=cZeek!6},
  gt/.style    ={boite, draw=cGT,    text=cGT!60!black, fill=cGT!8},
  script/.style={boite, draw=cScript,text=cScript,fill=cScript!7},
  sortie/.style={boite, draw=cOut,   text=cOut!70!black, fill=cOut!9},
  afleche/.style ={-{Latex[length=2.3mm]}, thick},
  aflechep/.style={-{Latex[length=2.0mm]}, thick, dashed},
  aetiq/.style   ={font=\scriptsize\itshape, text=cComment},
}

% ---------- Blocs de code de l'annexe ----------
\lstset{
  basicstyle=\ttfamily\scriptsize, backgroundcolor=\color{cCode},
  frame=single, framerule=0pt, rulecolor=\color{cCode},
  xleftmargin=6pt, xrightmargin=3pt, framexleftmargin=5pt,
  breaklines=true, columns=fullflexible, keepspaces=true,
  showstringspaces=false, tabsize=2, aboveskip=3pt, belowskip=3pt,
  commentstyle=\color{cComment}\itshape, stringstyle=\color{cStr},
  literate=%
    {é}{{\'e}}1 {è}{{\`e}}1 {ê}{{\^e}}1 {à}{{\`a}}1 {â}{{\^a}}1
    {î}{{\^\i}}1 {ï}{{\"\i}}1 {ô}{{\^o}}1 {û}{{\^u}}1 {ù}{{\`u}}1
    {ç}{{\c c}}1 {É}{{\'E}}1 {«}{{\guillemotleft}}1 {»}{{\guillemotright}}1
    {’}{{'}}1 {–}{{-}}1 {…}{{...}}1,
}
\lstdefinestyle{bash}{language=bash}
\lstdefinestyle{json}{language=,moredelim=[s][\color{cTar}\bfseries]{"}{":}}
\lstdefinestyle{jsonfull}{language=,basicstyle=\ttfamily\tiny,
  moredelim=[s][\color{cTar}\bfseries]{"}{":}, aboveskip=3pt, belowskip=3pt}


%--------------------- Paramètres de l'essai (à modifier) --------------------
\newcommand{\titreessai}{Intégration du référentiel MITRE ATT\&CK dans un écosystème MCP}
\newcommand{\auteuressai}{Kamil Saleh}
\newcommand{\moisdepot}{Septembre 2026}
\newcommand{\pictogrammeia}{logo-ia-assistance-limitee.png}

\definecolor{uqacvert}{RGB}{86,125,42}

%--------------------------- Profondeur de titraille -------------------------
\setcounter{secnumdepth}{3}   % numérote jusqu'à 1.2.1.1
\setcounter{tocdepth}{2}      % table des matières : chapitres, sections, sous-sections
                              % (mettre 3 pour y inclure aussi les sous-sous-sections)

%------------------------- Titres francisés (babel) --------------------------
\addto\captionsfrench{%
  \renewcommand{\contentsname}{Table des matières}%
  \renewcommand{\listfigurename}{Table des figures}%
  \renewcommand{\listtablename}{Liste des tableaux}%
  \renewcommand{\tablename}{Tableau}%
  \renewcommand{\figurename}{Figure}%
  \renewcommand{\bibname}{Bibliographie}%
  \renewcommand{\chaptername}{Chapitre}%
}

%--------------------------- Format des titres -------------------------------
% CHAPITRE n  /  TITRE DU CHAPITRE   (centrés, en capitales)
\makeatletter
\titleformat{\chapter}[display]
  {\normalfont\normalsize\centering}
  {\MakeUppercase{\@chapapp}\ \thechapter}   % « CHAPITRE n », puis « ANNEXE A »
  {1.2em}
  {\MakeUppercase}
\makeatother
\titlespacing*{\chapter}{0pt}{0pt}{40pt}

% 1.1 TITRE DE SECTION (gras, capitales, aligné à gauche)
\titleformat{\section}{\normalfont\normalsize\bfseries}{\thesection}{1em}{\MakeUppercase}
\titlespacing*{\section}{0pt}{3ex plus 1ex minus .2ex}{2ex plus .2ex}

\titleformat{\subsection}{\normalfont\normalsize\bfseries}{\thesubsection}{1em}{}
\titlespacing*{\subsection}{0pt}{2.5ex plus 1ex minus .2ex}{1.5ex plus .2ex}

\titleformat{\subsubsection}{\normalfont\normalsize\itshape}{\thesubsubsection}{1em}{}
\titlespacing*{\subsubsection}{0pt}{2.2ex plus 1ex minus .2ex}{1.2ex plus .2ex}

%------------------------------ En-têtes -------------------------------------
\setlength{\headheight}{15pt}
\pagestyle{fancy}
\fancyhf{}
\fancyhead[LE,RO]{\thepage}      % numéro en haut, côté extérieur
\renewcommand{\headrulewidth}{0pt}
\fancypagestyle{plain}{\fancyhf{}\renewcommand{\headrulewidth}{0pt}}

% Les versos laissés vides par \cleardoublepage ne portent pas de numéro
\makeatletter
\def\cleardoublepage{\clearpage\if@twoside\ifodd\c@page\else
  \hbox{}\thispagestyle{empty}\newpage\if@twocolumn\hbox{}\newpage\fi\fi\fi}
\makeatother

%------------------------------- Légendes ------------------------------------
% Rendu : FIGURE 1.1 – Titre de la figure   /   TABLEAU 1.1 – Titre du tableau
% (ajouter textfont=bf pour mettre tout le texte de la légende en gras)
\captionsetup{font=small,labelfont={sc,bf},labelsep=endash,
              justification=justified,skip=6pt,belowskip=4pt}

%-------------------- Outils repris du chapitre 1 ----------------------------
% Chapitre non numéroté, mais inscrit dans la table des matières
% (résumé, introduction générale, conclusion générale)
\newcommand{\chapterwithtoc}[1]{%
  \chapter*{#1}\addcontentsline{toc}{chapter}{#1}}

% Rubrique non numérotée à l'intérieur d'une sous-section
\newcommand{\soustitre}[1]{%
  \par\addvspace{1.1em}%
  {\normalsize\bfseries #1\par}%
  \addvspace{0.25em}\nopagebreak}

% Encadré « Conclusion » en fin de rubrique
\definecolor{grisencadre}{gray}{0.96}
\definecolor{filetencadre}{gray}{0.55}
\newtcolorbox{encadre}{
  enhanced, breakable,
  colback=grisencadre, colframe=filetencadre,
  boxrule=0pt, leftrule=2.5pt,
  arc=0pt, outer arc=0pt,
  left=8pt, right=8pt, top=6pt, bottom=6pt,
  before skip=1em, after skip=1em
}

% Légende placée SOUS un longtable.
% \begin{longtable} incrémente déjà le compteur « table » : on le corrige
% pour que chaque tableau ne consomme qu'un seul numéro.
\newcommand{\legendetableau}[2]{%
  \addtocounter{table}{-1}\captionof{table}{#1}\label{#2}}

% Nécessaires aux tableaux convertis depuis pandoc
\providecommand{\tightlist}{\setlength{\itemsep}{0pt}\setlength{\parskip}{0pt}}
\makeatletter
\@ifundefined{real}{\newcommand{\real}[1]{#1}}{}
\makeatother

% Figures dessinées en TikZ : le fichier est mis à l'échelle de la justification.
\newcommand{\figuretikz}[1]{\resizebox{\linewidth}{!}{\input{figures/#1.tex}}}

% Insère l'image si elle est présente, sinon un cadre de rappel.
% Permet de compiler avant d'avoir déposé toutes les figures dans figures/.
\newcommand{\figurefichier}[1]{%
  \IfFileExists{#1.pdf}{\includegraphics[width=0.98\linewidth,keepaspectratio]{#1}}{%
  \IfFileExists{#1.png}{\includegraphics[width=0.98\linewidth,keepaspectratio]{#1}}{%
  \IfFileExists{#1.jpg}{\includegraphics[width=0.98\linewidth,keepaspectratio]{#1}}{%
  \fbox{\begin{minipage}[c][3.5cm][c]{0.9\linewidth}\centering\itshape\small
    Image absente : déposer le fichier \texttt{\detokenize{#1}.png} (ou .pdf)
  \end{minipage}}}}}}

% Flottants : ne pas laisser les figures s'échapper de leur section
\renewcommand{\topfraction}{0.90}
\renewcommand{\bottomfraction}{0.70}
\renewcommand{\textfraction}{0.10}
\renewcommand{\floatpagefraction}{0.75}
\setcounter{topnumber}{3}
\setcounter{bottomnumber}{2}
\setcounter{totalnumber}{4}

% Justification : évite les débordements en marge
\tolerance=2000
\emergencystretch=2em
\hbadness=4000

%=============================================================================
\begin{document}

%--------------------------- PAGE DE GARDE -----------------------------------
% Page de garde centrée sur la page (marges symétriques)
\newgeometry{letterpaper,left=2.5cm,right=2.5cm,top=2.5cm,bottom=2.5cm}
\begin{titlepage}
\singlespacing
\centering

\vspace*{\stretch{0.4}}

% Logo UQAC — remplacer par le logo officiel si vous en disposez :
% \includegraphics[width=4cm]{images/logo-uqac.png}
{\fontsize{36}{40}\selectfont\color{uqacvert}\textsf{UQAC}}\\[4pt]
{\small\textsf{Université du Québec}}\\
{\small\textsf{à Chicoutimi}}

\vspace{\stretch{2}}


\begin{minipage}{0.62\textwidth}
  \centering\bfseries\MakeUppercase{\titreessai}
\end{minipage}


\vspace{\stretch{2}}

\textbf{ESSAI}\\[1.4\baselineskip]

RÉALISÉ AU SEIN DU LABORATOIRE CYNERGIA

\vspace{\stretch{1.6}}

\textbf{PAR}\\[1.4\baselineskip]
\MakeUppercase{\auteuressai}

\vspace{\stretch{1.6}}

% Réduire la largeur ci-dessous pour changer l'endroit de la coupure du titre,
% ou insérer un \\ directement dans \titreessai pour forcer le retour à la ligne.


\vspace{\stretch{1.6}}

\MakeUppercase{\moisdepot}

\vspace{\stretch{1.4}}

% Déclaration d'utilisation de l'IA générative (pictogramme UQAC)
\includegraphics[width=0.36\textwidth]{\pictogrammeia}

\vspace*{\stretch{0.2}}
\end{titlepage}
\restoregeometry

%--------------------------- PAGES LIMINAIRES --------------------------------
\pagenumbering{roman}
\setcounter{page}{1}


\chapterwithtoc{Résumé}
\label{chap:resume}

Dans un centre d'opérations de sécurité, le goulot d'étranglement n'est plus la
détection mais la qualification de ce qui est détecté. Rattacher une activité
observée au référentiel MITRE ATT\&CK reste un travail manuel, lent et
tributaire d'une expertise rare, et la difficulté redouble chez un exploitant
d'infrastructure énergétique, dont les incidents traversent la frontière entre
informatique de gestion et informatique industrielle.

Cet essai conçoit et mesure un écosystème agentique ancrant un modèle de langage
sur ATT\&CK au moyen du \emph{Model Context Protocol}. Un corpus de
42 échantillons est produit en laboratoire, à l'abri de la contamination qui
fausse les bancs d'essai publics, sur deux zones --- systèmes d'information et
automates --- et deux matrices, Enterprise et ICS. Dix modèles de la génération
courante y sont évalués sans ancrage, en trois exécutions par échantillon, soit
1\,260 appels, selon une méthodologie qui sépare cinq catégories de prédiction
et deux bases de calcul. Le modèle retenu est ensuite rejoué avec le serveur, à
corpus et conditions identiques.

Sans ancrage, aucun modèle ne dépasse 0,603 de F1 au niveau technique, et la
performance chute d'un tiers sur la matrice industrielle. Les identifiants
fabriqués sont exclusivement des tactiques, et exclusivement sur des journaux
industriels : l'erreur est un défaut de repérage entre univers, non une
invention arbitraire. Avec ancrage, le F1 technique passe de 0,563 à 0,706, la
matrice industrielle double presque, et le taux d'identifiants périmés ne
diminue pas : il disparaît. La contribution ne tient pas à l'accès à une source,
mais à quatre décisions de conception --- version épinglée et jointe à chaque
réponse, résolution des révocations, granularité d'outils bornée, trace d'usage
vérifiable.

\vspace{1em}
\noindent\textbf{Mots-clés :} MITRE ATT\&CK, Model Context Protocol, grands
modèles de langage, ancrage, centre d'opérations de sécurité, technologies
opérationnelles, ICS, évaluation.

\cleardoublepage

\addcontentsline{toc}{chapter}{\contentsname}
\tableofcontents
\cleardoublepage

\addcontentsline{toc}{chapter}{\listfigurename}
\listoffigures
\cleardoublepage

\addcontentsline{toc}{chapter}{\listtablename}
\listoftables
\cleardoublepage

%----------------------------- REMERCIEMENTS ---------------------------------
\pagenumbering{arabic}
\setcounter{page}{1}

\chapter*{Remerciements}
\addcontentsline{toc}{chapter}{Remerciements}

Je tiens à remercier en premier lieu mon directeur d'essai, Kevin Bouchard,
ainsi que Fehmi Jaafar, pour leur encadrement, leur disponibilité et la
rigueur de leurs conseils tout au long de ce travail. Leurs relectures
successives et leurs questions ont orienté ce projet autant qu'elles l'ont
corrigé.

Mes remerciements vont également à Franck Jeannot, mon superviseur, dont
l'accompagnement en contexte professionnel a permis d'ancrer cette recherche
dans des préoccupations opérationnelles réelles.

Je remercie le laboratoire Cynergia et l'ensemble de son équipe pour l'accès aux
ressources et aux environnements de test, sans lesquels la campagne
expérimentale présentée dans cet essai n'aurait pas été possible, ainsi que pour
les échanges qui l'ont nourrie.

Enfin, je remercie ma famille et mes proches pour leur soutien constant et pour
leurs nombreuses relectures.

%=========================== INTRODUCTION GÉNÉRALE ===========================

\chapterwithtoc{Introduction générale}
\label{chap:introduction}

Pendant longtemps, défendre un système d'information a consisté à en garder
la porte. On admet aujourd'hui que l'intrusion finira par se produire, et que la valeur d'une défense se mesure à sa capacité de la détecter, de la comprendre et de la contenir pendant qu'elle a lieu. Au sein d'un centre d'opérations de sécurité, deux équipes sur trois déclarent ne plus suivre le rythme, et les analystes juniors se renouvellent plus vite qu'ils ne se forment~\cite{sans2024}. La recherche a donné un nom à ce phénomène, la \emph{fatigue d'alerte}~\cite{tariq}, et un diagnostic : le
goulot d'étranglement n'est pas la détection, mais la qualification de ce
qui est détecté.

Le référentiel MITRE ATT\&CK~\cite{attack} est devenu la langue commune de cet
exercice ; y cartographier une activité observée reste pourtant un travail
manuel, lent, et tributaire d'une expertise rare. Pour un exploitant
d'infrastructure énergétique, la difficulté redouble : l'incident type
traverse la frontière entre l'informatique de gestion et l'informatique
industrielle (automates, capteurs, protocoles de terrain) deux mondes
qui ne parlent ni les mêmes journaux ni la même matrice ATT\&CK, et dont la
convergence est précisément ce que visent les campagnes les plus avancées.

Les grands modèles de langage (LLM) semblent taillés pour cette tâche. Mais on ne
peut pas encore leur faire confiance : ils \emph{hallucinent}~\cite{ji}, et leur savoir est
en outre figé à la date de leur entraînement, quand le référentiel, lui,
évolue à chaque version. La solution ce n'est pas un modèle plus grand mais un
modèle \emph{ancré}, rattaché au moment de répondre à une source vérifiable et
à jour plutôt qu'à sa seule mémoire. Le \emph{Model Context Protocol} (MCP),
ouvert fin 2024 et confié depuis à une gouvernance
neutre~\cite{mcpintro,aaif}, standardise cet ancrage en donnant au modèle un
accès structuré, déterministe et journalisable à des outils et à des données.

C'est dans ce contexte que s'inscrit le présent essai, qui conçoit un écosystème agentique intégrant le MCP et MITRE ATT\&CK pour
détecter et contextualiser les attaques. Quatre questions le guident :
\begin{enumerate}\setlength{\itemsep}{0pt}\setlength{\parsep}{0pt}\setlength{\topsep}{2pt}
  \item \textbf{Q1 --- Où en sont les modèles sans ancrage ?} Dans quelle mesure
        les modèles de langage actuels savent-ils rattacher au référentiel une
        activité malveillante lue dans des journaux bruts, aux trois niveaux de
        la hiérarchie (tactique, technique, sous-technique) ? Cette capacité se
        maintient-elle lorsque l'on passe de la matrice Enterprise à la matrice
        ICS, c'est-à-dire du journal de système d'information au journal
        d'automate ? La réponse suppose un corpus produit en laboratoire, à
        l'abri de la contamination qui fausse les bancs d'essai publics.
  \item \textbf{Q2 --- Que change l'ancrage déterministe ?} À corpus, consigne
        et conditions d'appel identiques, que gagne un modèle qui interroge un
        référentiel avant de répondre par rapport au même modèle répondant de
        mémoire ? Le gain attendu porte sur trois défauts distincts : les
        identifiants fabriqués, les identifiants périmés, et la confusion entre
        les deux matrices. Encore faut-il les mesurer séparément.
  \item \textbf{Q3 --- Comment construire le serveur qui le permet ?} Quelles
        propriétés un serveur MCP doit-il réunir pour que cet ancrage soit
        exact, reproductible et daté : quelle granularité d'outils, quelle
        version du référentiel, quel traitement des identifiants révoqués,
        quelle trace jointe à chaque réponse ? La littérature fixe des
        conditions d'emploi, elle ne les éprouve pas sur cette tâche.
  \item \textbf{Q4 --- À quel coût en surface d'attaque ?} Ancrer un modèle
        déplace le risque du texte vers l'action : le serveur devient un canal
        d'instruction autant qu'une source de faits. Quelles garanties de
        conception --- lecture seule, moindre privilège, journalisation,
        validation des entrées --- rendent ce déplacement acceptable pour un
        exploitant d'infrastructure essentielle soumis à la LPCE ?
\end{enumerate}

%================================ CHAPITRES ==================================
%% ===================================================================
%% CHAPITRE 1 : Préparation
%% Revue de littérature et méthodologie d'évaluation
%% ===================================================================
\chapter{Préparation}\label{chap:preparation}
\label{chap:preparation}

Ce premier chapitre établit le cadre théorique et opérationnel de notre étude. Nous proposons d'abord une revue de littérature qui examine la capacité des approches actuelles à ancrer le raisonnement des grands modèles de langage en cyberdéfense, puis nous présentons la méthodologie d'évaluation retenue pour mesurer l'apport de notre approche.

\section{Revue de littérature}\label{sec:revue}
Cette revue se lit comme un entonnoir : on étudie les grands modèles de langage (LLM) en cybersécurité, puis l'ancrage par la récupération, l'ancrage par les outils jusqu'au MCP, puis le revers : la sécurité des systèmes agentiques.

\FloatBarrier
\subsection{Axe 1 --- Les LLM en cyberdéfense : valeur, mesure, limites}\label{sec:axe1}

Les synthèses récentes attestent d'usages allant de la détection d'intrusion à l'assistance à l'analyste ; la vaste revue de Kheddar dans \emph{Information Fusion} en fournit la cartographie pour les systèmes de détection d'intrusion \cite{kheddar}. L'axe examine successivement la tâche pivot de l'essai (le \emph{mapping} ATT\&CK), la façon de mesurer sans se tromper, la fiabilité, puis l'angle mort industriel.

\soustitre{Cartographier vers ATT\&CK : deux générations et un plafond}

\textbf{MITRE ATT\&CK} est une base de connaissances ouverte qui catalogue les comportements adverses en \emph{techniques} (p. ex. « vol d'identifiants en mémoire »), regroupées par \emph{tactiques} (l'objectif poursuivi) ; l'édition Enterprise compte environ deux cents techniques et, depuis la scission de \emph{Defense Evasion} en version 19, quinze tactiques \cite{attack}. \emph{Mapper} une menace consiste à lire une description d'activité --- rapport d'incident, journal, règle de détection --- et à retrouver la ou les techniques qu'elle illustre : c'est la brique de base de la détection contextualisée.

\textbf{Génération 1 : la classification supervisée.} L'outil emblématique est TRAM (\emph{Threat Report ATT\&CK Mapper}) du Center for Threat-Informed Defense \cite{tram} : un modèle SciBERT affiné sur des phrases étiquetées à la main ne couvre qu'une cinquantaine des techniques les plus fréquentes. Le plafond n'est pas paramétrique mais structurel : l'annotation experte est un goulot d'étranglement et la couverture reste partielle, comme le documente la synthèse de Rahman \emph{et al.} sur l'extraction automatique de renseignement \cite{rahman}.

\begin{figure}[htbp]
\centering
\figuretikz{fig-1-1-arc-ancrage}
\caption{L'arc de l'ancrage, fil conducteur des axes 1 à 3 : chaque étape corrige la précédente et en déplace la limite. L'essai porte sur la troisième étape et sur sa comparaison empirique avec la deuxième.}\label{fig:arc-ancrage}
\end{figure}

\textbf{Génération 1 bis : les graphes de connaissances.} AttacKG (ESORICS 2022) reconstruit, à partir d'un rapport, le graphe des actions de l'attaquant, puis agrège les graphes de plusieurs rapports en graphes de techniques \cite{attackg}. Sur un large corpus réel il identifie 28 262 techniques et 8393 indicateurs de compromission ; sur seize rapports annotés, il atteint des F1 de 0,887 (entités), 0,896 (dépendances) et 0,789 (techniques), au-dessus des méthodes antérieures.

\textbf{La mesure : CTIBench.} Ce banc d'essai (NeurIPS 2024, piste \emph{Datasets and Benchmarks}) évalue les LLM sur des tâches de renseignement, dont l'extraction de techniques (CTI-ATE) \cite{ctibench}. Les modèles ouverts y sont faibles (Llama-3-8B $\approx$ 0,16 de F1) ; les modèles propriétaires ne sont pas mûrs pour un usage non supervisé (ChatGPT-3.5 $\approx$ 0,31, ChatGPT-4 0,64, meilleur du panel)\footnote{Le papier est ambigu sur la métrique : le texte parle de micro-F1, le tableau de macro-F1. Les valeurs sont celles du tableau publié.}. Son prolongement AthenaBench montre que la domination des modèles propriétaires ne s'étend pas au raisonnement --- attribution d'un attaquant, choix d'une mitigation --- où tous restent faibles \cite{athenabench}.

\begin{encadre}
\textbf{Conclusion :} Classification ou graphe, la première génération bute sur un plafond qui tient à la nature de la tâche : le \emph{mapping} n'est pas une reconnaissance de mots-clés mais un raisonnement contextuel, et même le meilleur LLM généraliste s'arrête à 0,64 de F1 sans ancrage (\textbf{QR1}).
\end{encadre}

\soustitre{Mesurer sans se tromper : capacité et contamination}

CyberSecEval (Meta) mesure deux propensions : produire du code vulnérable et se laisser enrôler dans une attaque, sur 8 langages, 50 faiblesses CWE et 10 catégories de techniques ATT\&CK \cite{cyberseceval}. En moyenne les modèles suggèrent du code vulnérable dans environ 30\,\% des cas, et --- motif récurrent de cette revue --- les modèles les plus compétents sont souvent les plus enclins à l'insécurité.

La contamination des jeux de test est le piège méthodologique central. Si les exemples d'un banc d'essai, et surtout leurs réponses, figuraient dans les données d'entraînement, le modèle récite au lieu de raisonner. PrimeVul (ICSE 2025) le démontre sur la détection de vulnérabilités \cite{primevul} : avec un jeu dédupliqué et découpé chronologiquement (près de 7000 fonctions vulnérables, 228 800 saines, 140 CWE), un modèle de 7 milliards de paramètres crédité de 68,3\,\% de F1 sur BigVul s'effondre à 3,1\,\%, GPT-4 se rapproche du hasard, et 18,9\,\% des échantillons de BigVul fuitaient dans le jeu de test.

\begin{encadre}
\textbf{Conclusion :} Les bancs d'essai publics surestiment massivement les performances parce que les modèles ont déjà vu les données. Ce résultat justifie directement, dans l'essai, la génération de journaux propres en laboratoire plutôt que la réutilisation d'un jeu public (\textbf{QR1}).
\end{encadre}

\soustitre{Fiabilité : l'ancrage plutôt que la taille}

La synthèse de Ji \emph{et al.} distingue l'hallucination intrinsèque (le modèle contredit la source fournie) de l'extrinsèque (il affirme l'invérifiable) \cite{ji}. En cybersécurité les conséquences sont concrètes : attribuer une technique inexistante oriente l'investigation puis la remédiation vers la mauvaise cible, au point que Mezzi, Massacci et Tuma concluent à l'\emph{unreliability} des LLM pour le renseignement sur les menaces lorsqu'ils opèrent sans garde-fou \cite{mezzi}. Le point capital, et contre-intuitif, est que la fiabilité factuelle dépend de l'ancrage et non de la taille : le programme \emph{Inverse Scaling} (TMLR 2023) a mis au jour des tâches où un modèle plus grand fait pire, notamment parce qu'il s'appuie avec excès sur ses connaissances mémorisées \cite{inversescaling}.

Sur le terrain, enfin, la contrainte première n'est pas la détection mais le volume : la synthèse de Tariq \emph{et al.} sur la fatigue d'alerte (ACM CSUR 2025) décrit un déluge quotidien d'alertes majoritairement faussement positives, dont une part importante reste sans investigation \cite{tariq}. Les LLM y ont une valeur d'augmentation --- résumer une alerte, interpréter une commande, contextualiser un incident --- avec un humain dans la boucle pour les cas critiques.

\begin{encadre}
\textbf{Conclusion :} Ancrer le raisonnement dans une base vérifiable est plus efficace qu'augmenter la puissance de calcul ; c'est l'argument directeur de l'essai (\textbf{QR1}).
\end{encadre}

\soustitre{L'angle mort : la détection en environnement industriel (OT/ICS)}

L'informatique de gestion (IT) et l'informatique industrielle (OT : automates, capteurs, protocoles de terrain comme Modbus qui pilotent des procédés physiques) sont deux mondes ; or, pour un opérateur d'infrastructure critique, l'incident type traverse la frontière IT/OT. La quasi-totalité des travaux ci-dessus porte sur l'IT. Les rares travaux industriels sont éclairants : Kong, Saber, Youssef et Kundur évaluent un LLM prêt à l'emploi comme couche de détection explicable et supervisée sur du trafic Modbus d'un réseau électrique, en insistant sur des alarmes auditables --- comprises, tracées, archivées pour la réponse à incident \cite{kong} --- un cahier des charges très proche de celui d'Hydro-Québec. La synthèse de Kheddar confirme que le modèle est surtout employé comme assistant sur des rapports en langage naturel, rarement comme analyste de jetons de protocole industriel \cite{kheddar}.

\begin{encadre}
\textbf{Conclusion :} L'évaluation des LLM sur de la télémétrie OT reste largement inexplorée : c'est le vide que le laboratoire IT/OT de l'essai --- une simulation de procédé GRFICS v3 instrumentée par une sonde Zeek, produisant du trafic Modbus conforme au protocole réel et garanti non contaminé --- vient combler ; l'évaluation sur des journaux d'exploitation OT reste une perspective (\textbf{QR1}).
\end{encadre}

\FloatBarrier
\subsection{Axe 2 --- Ancrer par la récupération : le RAG et ses limites}\label{sec:axe2}

L'axe 1 a posé le problème : les LLM aident mais hallucinent, et leur fiabilité tient à l'ancrage. La communauté a exploré deux voies (figure~\ref{fig:arc-ancrage}) : récupérer de l'information, objet de cet axe, puis appeler des outils (axe 3).

\textbf{Définition.} La génération augmentée par récupération (RAG) a été formalisée par Lewis \emph{et al.} (NeurIPS 2020) comme le couplage d'un récupérateur dense --- qui plonge requête et documents dans un même espace vectoriel et sélectionne les \emph{k} passages les plus proches --- et d'un générateur conditionné sur ces passages \cite{lewis}. Le modèle ne répond plus « de mémoire » mais sur la foi de documents fournis à la volée, ce qui réduit l'hallucination et permet de citer. La taxonomie de Gao \emph{et al.} distingue le RAG naïf, avancé (reformulation, re-classement) et modulaire \cite{gao} ; les systèmes ci-dessous relèvent de la seconde catégorie.

\textbf{Application au \emph{mapping}.} TechniqueRAG est l'état de l'art : pour un texte donné, il récupère des exemples de techniques pertinents, les fait re-classer par un LLM, puis génère la réponse avec un modèle affiné, surpassant nettement les approches sans ancrage \cite{techniquerag}. RAM (\emph{Rule-ATT\&CK Mapper}) applique la même logique aux règles de détection d'un SIEM --- le moteur central du SOC qui collecte les journaux et déclenche les alertes --- et relie chaque règle aux techniques qu'elle détecte, aux composants de données et aux mitigations associés \cite{ram}.

\textbf{Trois limites.} (a) \emph{Similarité n'est pas pertinence} : le récupérateur renvoie ce qui ressemble à la requête, et un voisin vectoriel peut décrire une technique adjacente mais distincte, précisément là où le \emph{mapping} demande une discrimination fine. (b) \emph{Non-déterminisme} : le résultat dépend du modèle de plongement, du découpage des documents et de \emph{k} ; deux exécutions ne sont pas garanties identiques, ce qui complique l'audit d'une alerte. (c) \emph{La base est une surface d'attaque} : PoisonedRAG (USENIX Security 2025) formule l'empoisonnement comme un problème d'optimisation --- le texte injecté doit être récupéré \emph{et} induire la réponse voulue --- et obtient 90\,\% de réussite, jusqu'à 97\,\% sur certains jeux, en injectant cinq textes dans une base de plus de 2,6 millions de documents ; les défenses testées (reformulation, détection par perplexité) sont insuffisantes \cite{poisonedrag}. La figure~\ref{fig:rag-outils} oppose cette chaîne probabiliste à l'accès déterministe de l'axe 3.

\begin{figure}[htbp]
\centering
\figuretikz{fig-1-2-rag-vs-outils}
\caption{Deux façons d'ancrer le même raisonnement. Le RAG récupère un voisin plausible par similarité ; l'accès outillé récupère la définition exacte par une requête structurée. La surface d'attaque ne disparaît pas : elle se déplace de la base vers la chaîne d'approvisionnement du serveur.}\label{fig:rag-outils}
\end{figure}

\begin{encadre}
\textbf{Conclusion :} La récupération par ressemblance n'est ni pleinement fiable (un voisin peut être hors sujet, et le résultat n'est pas reproductible) ni pleinement sûre (la base est empoisonnable). D'où la recherche d'un ancrage plus déterministe (\textbf{QR2}).
\end{encadre}

\FloatBarrier
\subsection{Axe 3 --- Ancrer par les outils : de l'apprentissage par outils au MCP}\label{sec:axe3}

La seconde voie ne récupère pas un document ressemblant : elle appelle un outil qui renvoie un fait exact. C'est le paradigme de l'apprentissage par outils (\emph{tool learning}), fondement théorique direct du MCP.

\soustitre{Le paradigme de l'apprentissage par outils}

Toolformer (NeurIPS 2023) est un modèle qui a appris, de façon auto-supervisée, à décider lui-même quand insérer un appel d'outil au fil de sa génération : il insère des appels candidats dans des textes, les exécute, et ne conserve que ceux dont le résultat réduit sa perplexité sur la suite \cite{toolformer}. Cinq outils sont testés (calculatrice, recherche Wikipédia, questions-réponses, traduction, calendrier) sur un modèle de 6,7 milliards de paramètres. Sur les problèmes arithmétiques, Toolformer atteint 29,4\,\% (SVAMP) et 40,4\,\% (ASDiv) sans exemple, contre 10,0\,\% et 14,0\,\% pour GPT-3 --- 175 milliards de paramètres, soit un modèle environ 25 fois plus gros ; le contre-exemple est instructif : sur les questions factuelles, le moteur de recherche employé est trop rudimentaire et GPT-3 reprend l'avantage. L'outil ne vaut que ce que vaut sa source.

ReAct (ICLR 2023) est une méthode de sollicitation qui alterne explicitement « pensées » (raisonnement en langage naturel) et « actions » (appels d'outils) \cite{react}. En ancrant chaque étape dans un résultat concret, ReAct réduit l'hallucination sur HotpotQA et, sur deux environnements interactifs (ALFWorld, WebShop), dépasse de 34 et 10 points de taux de succès des méthodes entraînées sur 10\textsuperscript{3} à 10\textsuperscript{5} exemples, avec un ou deux exemples seulement. La synthèse de Qin \emph{et al.} (ACM CSUR 2024) généralise : doter un modèle d'outils atténue à la fois l'obsolescence de ses connaissances et ses hallucinations, et déplace le problème vers la sélection et la composition des outils \cite{qin}.

\begin{encadre}
\textbf{Conclusion :} Le bon appel d'outil vaut plus que vingt-cinq fois plus de paramètres --- à condition que l'outil renvoie un fait exact (\textbf{QR3}).
\end{encadre}

\soustitre{Le MCP : standardiser l'ancrage par outils}

\textbf{Définition et architecture.} Le \emph{Model Context Protocol}, ouvert par Anthropic en novembre 2024, n'est pas un modèle mais un protocole qui standardise l'accès d'une application LLM à des outils et à des données \cite{mcpintro}. L'analogie usuelle est le port USB-C : un connecteur unique remplace une multitude de câbles propriétaires. L'architecture (figure~\ref{fig:archi-mcp}) distingue l'hôte (l'application), qui instancie un client par connexion, et les serveurs qui exposent trois primitives : des \emph{outils} (fonctions invocables), des \emph{ressources} (données en lecture) et des \emph{invites} (gabarits) ; le client expose en retour \emph{sampling}, \emph{roots} et \emph{elicitation}. Les messages sont en JSON-RPC 2.0 sur deux transports : \emph{stdio} pour un serveur local et \emph{Streamable HTTP} pour un serveur distant, ce dernier protégé par une autorisation OAuth 2.1 \cite{mcpspec2025}.

\begin{figure}[htbp]
\centering
\figuretikz{fig-1-3-architecture-mcp}
\caption{Architecture du MCP. L'hôte instancie un client par serveur ; chaque serveur expose outils, ressources et invites sur JSON-RPC 2.0. Le serveur du bas est le connecteur SIEM de l'essai. Tout ce qui est à droite de la frontière est, par défaut, non fiable : c'est le point de départ de l'axe 4.}\label{fig:archi-mcp}
\end{figure}

\textbf{Évolution et gouvernance.} Le protocole a mûri vite (tableau~\ref{tab:revisions-mcp}). La révision 2025-03-26 a ajouté l'autorisation OAuth 2.1, le transport \emph{Streamable HTTP} et des annotations décrivant le comportement d'un outil (lecture seule, destructif) ; la révision 2025-06-18 a classé les serveurs comme \emph{resource servers} OAuth, rendu obligatoires les indicateurs de ressource (RFC 8707) pour empêcher un serveur malveillant d'obtenir des jetons, et introduit la sortie structurée et l'\emph{elicitation} \cite{mcpspec2025}. Le 9 décembre 2025, Anthropic a transféré le protocole à l'Agentic AI Foundation, fonds dirigé de la Linux Foundation cofondé avec Block et OpenAI, à un moment où plus de 10 000 serveurs publics étaient actifs \cite{aaif}. Enfin la révision 2026-07-28 --- la plus importante depuis le lancement --- rend le noyau du protocole sans état, durcit l'autorisation (validation de l'émetteur, \emph{Client ID Metadata Documents}), formalise un cadre d'extensions et une politique de dépréciation \cite{mcpspec2026} : un serveur distant devient une charge HTTP ordinaire, à laquelle s'applique l'outillage de sécurité HTTP classique.




\begin{longtable}[]{@{}
  >{\raggedright\arraybackslash}p{(\columnwidth - 4\tabcolsep) * \real{0.1580}}
  >{\raggedright\arraybackslash}p{(\columnwidth - 4\tabcolsep) * \real{0.4210}}
  >{\raggedright\arraybackslash}p{(\columnwidth - 4\tabcolsep) * \real{0.4210}}@{}}
\toprule\noalign{}
\begin{minipage}[b]{\linewidth}\raggedright
\textbf{Révision}
\end{minipage} & \begin{minipage}[b]{\linewidth}\raggedright
\textbf{Apports majeurs}
\end{minipage} & \begin{minipage}[b]{\linewidth}\raggedright
\textbf{Implication pour un serveur de SOC}
\end{minipage} \\
\midrule\noalign{}
\endhead
\bottomrule\noalign{}
\endlastfoot
2024-11-05 & Primitives (outils, ressources, invites) ; stdio et HTTP+SSE & Garde-fous de consentement en \emph{should} ; sécurité laissée à l'implémenteur \\
2025-03-26 & OAuth 2.1 ; \emph{Streamable HTTP} ; annotations d'outils ; lots JSON-RPC & Les annotations sont des indications : la spécification les déclare non fiables sauf serveur de confiance \\
2025-06-18 & Serveurs = \emph{resource servers} ; RFC 8707 obligatoire ; sortie structurée ; \emph{elicitation} ; page de bonnes pratiques & Premiers \emph{must} de sécurité ; sortie typée = validation possible côté client \\
2025-11-25 & Révision intermédiaire ; consolidation par SEP et groupes de travail & Gouvernance communautaire (AAIF) \\
2026-07-28 & Noyau sans état ; autorisation durcie (émetteur, CIMD) ; cadre d'extensions ; politique de dépréciation & Serveur distant = charge HTTP standard ; extensions à auditer séparément \\
\end{longtable}
\legendetableau{Révisions de la spécification MCP et implications de sécurité \cite{mcpintro,mcpspec2025,mcpspec2026}.}{tab:revisions-mcp}

\textbf{Application au \emph{mapping} : l'accès déterministe.} Le serveur mitre-mcp (Montimage) expose l'intégralité d'ATT\&CK sous forme d'outils interrogeables, avec des recherches en temps constant grâce à des index préconstruits \cite{mitremcp}. Lorsqu'il a besoin de la description d'une technique, le modèle ne récupère pas un voisin plausible : il exécute une requête structurée et obtient la définition exacte, de façon reproductible et journalisable. La surface d'empoisonnement par similarité disparaît ; elle se déplace vers la chaîne d'approvisionnement du serveur lui-même, objet de l'axe 4.

\textbf{Une réserve sur le référentiel.} Une analyse conjointe d'Anthropic et de Verizon, qui a cartographié 832 comptes malveillants et 13 873 actions sur la version 18 d'ATT\&CK (482 sous-techniques, 14 tactiques), souligne qu'ATT\&CK ne capture pas encore l'orchestration autonome par des agents d'IA, marque des campagnes les plus avancées \cite{anthropicverizon}. Le référentiel ancré reste donc une borne inférieure de la menace.

\soustitre{Ce que la littérature empirique mesure sur le MCP}

La littérature savante sur le MCP a moins de deux ans, mais quatre travaux portent directement sur la question de l'essai : combien d'outils exposer, ce que l'ancrage fait gagner, ce qu'il tient en conditions réelles, et ce qu'il coûte en sécurité. Le tableau~\ref{tab:litt-mcp} les rassemble.

\textbf{Trop d'outils dégrade la sélection.} Gan et Sun mesurent l'effet du nombre d'outils exposés à un modèle \cite{ragmcp}. Leur essai de charge présente au modèle $N$ descriptions de serveurs --- une seule pertinente, les autres tirées d'un annuaire de plus de 4400 serveurs publics --- et fait varier $N$ de 1 à plus de 11 000 sur vingt tâches de recherche. En deçà d'une trentaine de candidats, le bon outil est choisi dans plus de 90\,\% des cas ; entre trente et soixante-dix, les échecs apparaissent par grappes, les descriptions voisines se confondant ; au-delà d'une centaine, l'échec domine. Deux causes se cumulent : l'\emph{encombrement de l'invite}, le contexte étant saturé de distracteurs, et la \emph{charge de décision}, des outils aux fonctions proches devenant indiscernables. La mesure chiffrée est nette : présenter tous les outils d'un coup donne 13,62\,\% de sélections correctes pour 2134 jetons d'invite, un pré-filtrage par mots-clés 18,20\,\% pour 1646 jetons, et une pré-sélection sémantique n'injectant que le meilleur candidat 43,13\,\% pour 1084 jetons --- soit plus du triple de justesse pour la moitié des jetons.

\textbf{L'appel à la demande restaure l'efficacité.} Fei, Zheng et Feng inversent la logique : plutôt que d'injecter tous les schémas dès le départ, le modèle énonce lui-même le besoin qu'il ne peut satisfaire, et un routage à deux étages --- d'abord le serveur, ensuite l'outil --- ne lui renvoie que ce qu'il a demandé \cite{mcpzero}. Sur un corpus de 308 serveurs et 2797 outils issus du dépôt officiel, le bon outil est retrouvé parmi près de trois mille candidats représentant 248 100 jetons, et la consommation chute de 98\,\% sur APIBank à justesse maintenue. L'enseignement pour l'essai est direct : le coût de l'ancrage est une variable de conception, non une fatalité du protocole.

\textbf{En conditions réelles, l'apport reste partiel.} MCP-Universe (Salesforce AI Research) évalue les modèles sur 231 tâches réparties en six domaines et servies par 11 serveurs MCP réels, avec des évaluateurs qui vérifient l'exécution plutôt que la forme de la réponse \cite{mcpuniverse}. Aucun modèle ne franchit la moitié des tâches : GPT-5 réussit 43,72\,\%, Grok-4 33,33\,\% et Claude-4.0-Sonnet 29,44\,\%. Deux causes d'échec reviennent, et toutes deux concernent l'essai : le contexte enfle avec le nombre d'étapes, et les outils inconnus sont mal exploités. La tâche de l'essai --- un référentiel unique, en lecture seule, sur un nombre d'étapes borné --- se situe précisément dans le cas favorable que ce banc d'essai isole par contraste.

\textbf{La sécurité du protocole reste en retard.} Yang, Wu et Chen proposent la première formalisation systématique des surfaces d'attaque du MCP : 17 types d'attaques répartis sur quatre surfaces --- interaction avec l'utilisateur, client, transport et serveur --- éprouvés sur trois hôtes du commerce \cite{mcpsecbench}. Plus de 85\,\% de ces attaques compromettent au moins une plateforme, les failles de protocole et d'implémentation atteignant 100\,\% de réussite, et les protections intégrées n'en arrêtent qu'une fraction. Le constat annonce l'axe 4 : ce qui est standardisé, c'est l'accès, pas sa sûreté.


{\footnotesize
\begin{longtable}[]{@{}
  >{\raggedright\arraybackslash}p{(\columnwidth - 4\tabcolsep) * \real{0.2300}}
  >{\raggedright\arraybackslash}p{(\columnwidth - 4\tabcolsep) * \real{0.3600}}
  >{\raggedright\arraybackslash}p{(\columnwidth - 4\tabcolsep) * \real{0.4100}}@{}}
\toprule\noalign{}
\begin{minipage}[b]{\linewidth}\raggedright
\textbf{Travail}
\end{minipage} & \begin{minipage}[b]{\linewidth}\raggedright
\textbf{Dispositif}
\end{minipage} & \begin{minipage}[b]{\linewidth}\raggedright
\textbf{Résultat clé}
\end{minipage} \\
\midrule\noalign{}
\endhead
\bottomrule\noalign{}
\endlastfoot
Gan et Sun (2025) \cite{ragmcp} & Essai de charge, $N$ de 1 à 11 000 serveurs, annuaire de 4400+ & Effondrement au-delà de 100 outils ; 13,62\,\% $\rightarrow$ 43,13\,\% avec pré-sélection, jetons divisés par deux \\
MCP-Zero (2025) \cite{mcpzero} & 308 serveurs, 2797 outils ; appel d'outil à la demande & Bon outil parmi $\approx$3000 candidats (248 100 jetons) ; $-$98\,\% de jetons sur APIBank \\
MCP-Universe (2025) \cite{mcpuniverse} & 231 tâches, 11 serveurs réels, 6 domaines, évaluation par exécution & GPT-5 43,72\,\%, Grok-4 33,33\,\%, Claude-4.0-Sonnet 29,44\,\% ; contexte long et outils inconnus \\
MCPSecBench (2025) \cite{mcpsecbench} & 17 attaques, 4 surfaces, 3 hôtes du commerce & \textgreater85\,\% des attaques réussissent sur au moins un hôte ; protocole et implémentation à 100\,\% \\
\end{longtable}
\legendetableau{Ce que la littérature empirique établit sur le MCP (axe 3).}{tab:litt-mcp}
}

\textbf{Ce que ces quatre travaux imposent à la conception du serveur.} Ils convergent vers une règle simple : un serveur MCP utile est un serveur \emph{étroit}. Exposer peu d'outils, aux périmètres nettement séparés, préserve la justesse de sélection que le premier travail voit s'effondrer ; ne servir que la réponse demandée, sans contexte superflu, préserve le budget de jetons que le deuxième optimise ; borner le nombre d'étapes évite la dérive que le troisième observe ; et n'offrir aucune écriture réduit la surface que le quatrième mesure. Ces quatre contraintes sont celles retenues pour le serveur de l'essai.

\begin{encadre}
\textbf{Conclusion :} L'accès déterministe est l'aboutissement de l'arc de l'ancrage, et le MCP en est désormais le standard sous gouvernance neutre. La littérature empirique en fixe déjà les conditions d'emploi --- peu d'outils, servis à la demande, sur un nombre d'étapes borné --- mais elle mesure l'ancrage outillé sur des tâches génériques, jamais sur le \emph{mapping} ATT\&CK en conditions de SOC : c'est l'un des buts directs de l'essai (\textbf{QR3}).
\end{encadre}

\FloatBarrier
\subsection{Axe 4 --- Le revers : sécurité des agents outillés et du MCP}\label{sec:axe4}

Les trois premiers axes construisent l'argument pour l'ancrage outillé ; le quatrième en examine le revers. La bascule est fondamentale : un LLM isolé qui se trompe produit un texte erroné ; un LLM connecté à des outils qui se trompe peut exécuter une action sur un système réel. Le risque passe du registre de l'erreur à celui du dommage, ce que l'OWASP nomme \emph{excessive agency} dans son classement 2025 des risques des applications LLM \cite{owasp}.

\soustitre{La faille fondatrice : l'injection indirecte de prompt}

On distingue l'injection directe (l'attaquant tape l'instruction) de l'injection indirecte, plus insidieuse : l'instruction est cachée dans un contenu tiers --- page web, courriel, résultat d'outil --- que le modèle lit au cours de sa tâche. Greshake \emph{et al.} (AISec, ACM CCS 2023) en ont établi le principe et l'ont démontré sur des applications réelles : parce que le modèle ne distingue pas les données qu'il traite des instructions qu'il doit suivre, un contenu piégé le détourne à distance \cite{greshake}. Deux bancs d'essai quantifient la vulnérabilité des agents outillés, cas exact du MCP. InjecAgent (ACL 2024) compose 1054 cas à partir de 17 outils légitimes et 62 outils d'attaquant : un agent GPT-4 guidé par ReAct est détourné dans 24\,\% des cas, 47\,\% avec des invites d'attaque renforcées \cite{injecagent}. AgentDojo (NeurIPS 2024) va plus loin avec un environnement dynamique de 97 tâches réalistes et 629 cas de sécurité qui mesure conjointement utilité et résistance : sous attaque, l'utilité de GPT-4o chute d'environ 69\,\% à 50\,\%, avec des taux de réussite d'attaque atteignant 53\,\% \cite{agentdojo}.

\begin{encadre}
\textbf{Conclusion :} Les agents outillés sont, en l'état, largement vulnérables, et la vulnérabilité tient à une propriété du modèle --- il suit ce qu'il lit --- que ni la taille ni l'alignement n'ont résolue (\textbf{QR4}).
\end{encadre}

\soustitre{Ce que le MCP ajoute --- mesuré sur des serveurs réels}

\textbf{Taxonomie.} L'étude de paysage de Hou \emph{et al.} fait référence : elle organise les menaces le long du cycle de vie du serveur --- création (collision de noms, installateur usurpé, code injecté), opération (conflits de noms d'outils, chevauchement de commandes, évasion du bac à sable) et mise à jour (persistance de privilèges, redéploiement de versions vulnérables, dérive de configuration) --- et identifie un mécanisme central, la propagation de confiance : dans un hôte combinant plusieurs serveurs, un seul serveur compromis peut influencer le comportement des outils légitimes \cite{hou}. La figure~\ref{fig:surface-attaque} en donne une lecture opérationnelle.

L'empoisonnement d'outil (\emph{tool poisoning}) est la menace spécifique qui dicte la conception à venir : les instructions malveillantes sont glissées non dans les données échangées mais dans la \emph{description} de l'outil, ce texte que le modèle lit pour savoir quand et comment l'employer. Révélée par Invariant Labs en avril 2025 avec une démonstration d'exfiltration de conversations \cite{invariant}, elle a été mesurée à grande échelle par MCPTox (AAAI 2026) : à partir de 45 serveurs réels et de 353 outils authentiques, plus de 1300 cas d'attaque évalués sur 20 modèles donnent un taux de réussite moyen de 36,5\,\%, culminant à 72,8\,\% (o1-mini), le taux de refus restant partout inférieur à 3\,\% \cite{mcptox}. Fait marquant, en écho aux axes 1 et 3, les modèles les plus capables sont souvent les plus vulnérables, l'attaque exploitant leur meilleure aptitude à suivre les instructions.

\textbf{L'état de l'écosystème.} La première étude empirique à grande échelle, publiée dans ACM TOSEM, a analysé 1899 serveurs MCP libres en combinant analyse statique générale et scanner spécifique \cite{hasan} : huit classes de vulnérabilités dont trois seulement recoupent le logiciel classique ; 7,2\,\% des serveurs portent une vulnérabilité générale et 5,5\,\% un empoisonnement d'outil ; 66\,\% présentent des odeurs de code et 14,4\,\% des motifs de bogues connus. Côté vulnérabilités avérées, CVE-2025-6514 (CVSS 9,6) permettait l'exécution de commandes arbitraires via l'utilitaire mcp-remote \cite{cve6514}, et CVE-2025-49596 (CVSS 9,4) exploitait l'absence d'authentification par défaut de l'inspecteur MCP \cite{cve49596} ; un audit outillé (MCPSafetyScanner) avait dès avril 2025 obtenu exécution de code et vol d'identifiants à travers des serveurs courants \cite{radosevich}.

\begin{figure}[htbp]
\centering
\figuretikz{fig-1-4-surface-attaque}
\caption{Surface d'attaque d'un serveur MCP le long de son cycle de vie (taxonomie adaptée de Hou et al. \cite{hou}, chiffres de MCPTox \cite{mcptox} et des CVE \cite{cve6514,cve49596}) et réponses de conception retenues pour le serveur de l'essai.}\label{fig:surface-attaque}
\end{figure}
\begin{encadre}
\textbf{Conclusion :} Ces failles ne sont pas exotiques : la plupart sont des manquements aux principes classiques (validation des entrées, moindre privilège, authentification). Une seule classe est nouvelle --- la description d'un outil devient un vecteur d'instruction --- et c'est elle qui impose des contrôles spécifiques au MCP (\textbf{QR4}).
\end{encadre}

\soustitre{Un modèle de sécurité en rattrapage, et les réponses émergentes}

\textbf{La spécification.} Une part des garde-fous de consentement et de bonnes pratiques y est formulée en \emph{should} (recommandé) plutôt qu'en \emph{must} (obligatoire), ce qui les rend optionnels pour l'implémenteur ; la spécification prévient d'ailleurs elle-même que les descriptions et annotations d'outils doivent être considérées comme non fiables sauf serveur de confiance \cite{mcpspec2025}. Les révisions récentes comportent de véritables \emph{must}, surtout sur l'autorisation (tableau~\ref{tab:revisions-mcp}), mais rien n'oblige un serveur à être en lecture seule, à journaliser, ni à figer ses descriptions : ces propriétés relèvent de la conception de chaque serveur.

\textbf{Les organismes.} Le centre de sécurité de l'IA de la NSA constate, dans une fiche de mai 2026, que la prolifération du MCP a devancé le développement de son modèle de sécurité\footnote{La co-signature de la CISA n'apparaît pas dans le communiqué officiel ; la référence est attribuée à la seule NSA (AISC).} \cite{nsa}. L'OWASP a publié un classement des risques des applications LLM (2025) puis des applications agentiques (2026) \cite{owasp}. Formaliser ces menaces dans un modèle auditable suppose une taxonomie dédiée à l'IA : MITRE ATLAS, calqué sur ATT\&CK mais orienté vers les attaques contre les systèmes d'IA, couvre désormais explicitement la compromission de serveurs MCP et l'injection indirecte \cite{atlas}. Au Canada, cette exigence d'auditabilité devient légale : la \emph{Loi sur la protection des cybersystèmes essentiels} (projet de loi C-8, sanction royale en juin 2026, qui reprend le projet C-26 devenu caduc) impose aux exploitants désignés d'infrastructures critiques un programme de cybersécurité documenté et la déclaration des incidents \cite{lpce}.

\textbf{La défense outillée.} Le tableau~\ref{tab:vulnerabilite} rassemble les mesures ; côté défense active, l'automatisation de l'audit progresse. PentestGPT (USENIX Security 2024) coordonne trois modules --- raisonnement, génération de commandes, analyse des sorties --- autour d'un arbre de tâches qui maintient le fil de l'investigation, et augmente le taux de complétion des épreuves de 228,6\,\% par rapport à GPT-3.5, sous supervision humaine \cite{pentestgpt}. Des scanners spécifiques (mcp-scan, MCPSafetyScanner) détectent l'empoisonnement de descriptions \cite{invariant,radosevich}, et les passerelles de médiation qui filtrent et journalisent les appels sont la piste d'architecture la plus citée \cite{hou}.


\begin{longtable}[]{@{}
  >{\raggedright\arraybackslash}p{(\columnwidth - 4\tabcolsep) * \real{0.2370}}
  >{\raggedright\arraybackslash}p{(\columnwidth - 4\tabcolsep) * \real{0.3815}}
  >{\raggedright\arraybackslash}p{(\columnwidth - 4\tabcolsep) * \real{0.3815}}@{}}
\toprule\noalign{}
\begin{minipage}[b]{\linewidth}\raggedright
\textbf{Travail (lieu)}
\end{minipage} & \begin{minipage}[b]{\linewidth}\raggedright
\textbf{Dispositif}
\end{minipage} & \begin{minipage}[b]{\linewidth}\raggedright
\textbf{Résultat clé}
\end{minipage} \\
\midrule\noalign{}
\endhead
\bottomrule\noalign{}
\endlastfoot
PoisonedRAG (USENIX Sec'25) \cite{poisonedrag} & 5 textes injectés dans 2,6 M de documents & 90\,\% de réussite, jusqu'à 97\,\% \\
InjecAgent (ACL'24) \cite{injecagent} & 1054 cas, 17 outils légitimes, 62 d'attaquant & GPT-4 + ReAct détourné à 24\,\%, 47\,\% renforcé \\
AgentDojo (NeurIPS'24) \cite{agentdojo} & 97 tâches, 629 cas de sécurité, utilité + attaque & GPT-4o : utilité 69\,\% $\rightarrow$ 50\,\% ; attaque $\leq$ 53\,\% \\
MCPTox (AAAI'26) \cite{mcptox} & 45 serveurs réels, 353 outils, \textgreater1300 cas, 20 modèles & Moy. 36,5\,\%, max 72,8\,\% ; refus \textless{} 3\,\% \\
Hasan \emph{et al.} (TOSEM'26) \cite{hasan} & 1899 serveurs libres, analyse statique + scanner & 7,2\,\% vulnérables, 5,5\,\% empoisonnés, 66\,\% d'odeurs de code \\
CVE-2025-6514 / 49596 \cite{cve6514,cve49596} & mcp-remote ; inspecteur MCP & CVSS 9,6 et 9,4 : exécution de code à distance \\
\end{longtable}
\legendetableau{Mesures de la vulnérabilité des agents outillés et du MCP (axe 4).}{tab:vulnerabilite}

\begin{encadre}
\textbf{Le \emph{rug pull} : pourquoi la confiance ne doit jamais être acquise}

Un serveur se comporte correctement au moment de la validation initiale, puis modifie silencieusement ses outils une fois la confiance accordée \cite{hou}. Appliqué à notre cas, un outil de \emph{mapping} pourrait retirer discrètement certaines techniques de ses résultats, rendant des pans entiers d'attaques invisibles sans que l'analyste ne perçoive la dérive. C'est ce scénario, plus que l'exécution de code, qui dicte les choix du serveur développé dans l'essai : aucun outil d'écriture, référentiel en lecture seule vérifié par empreinte, descriptions d'outils gelées, bascule atomique des mises à jour et journalisation systématique de chaque appel (figure~\ref{fig:surface-attaque}).
\end{encadre}

\begin{encadre}
\textbf{Conclusion :} La voie d'un audit de sécurité outillé et ancré est ouverte : c'est ce que le chapitre suivant prolonge avec un flux multi-MCP fondé sur ATLAS, sous contrainte réglementaire canadienne (\textbf{QR4}).
\end{encadre}

\subsection{Synthèse : acquis et lacunes}\label{sec:synthese-revue}

Le tableau~\ref{tab:synthese-revue} met en regard, pour chaque question de revue, l'acquis chiffré de la littérature et la lacune qu'elle laisse ouverte.


\begin{longtable}[]{@{}
  >{\raggedright\arraybackslash}p{(\columnwidth - 4\tabcolsep) * \real{0.0858}}
  >{\raggedright\arraybackslash}p{(\columnwidth - 4\tabcolsep) * \real{0.4571}}
  >{\raggedright\arraybackslash}p{(\columnwidth - 4\tabcolsep) * \real{0.4571}}@{}}
\toprule\noalign{}
\begin{minipage}[b]{\linewidth}\raggedright
\end{minipage} & \begin{minipage}[b]{\linewidth}\raggedright
\textbf{Acquis (chiffré)}
\end{minipage} & \begin{minipage}[b]{\linewidth}\raggedright
\textbf{Lacune actuelle}
\end{minipage} \\
\midrule\noalign{}
\endhead
\bottomrule\noalign{}
\endlastfoot
\textbf{QR1} & Les LLM aident mais hallucinent ; leur mesure est piégée par la contamination (F1 68\,\% $\rightarrow$ 3\,\% \cite{primevul}) ; le meilleur modèle généraliste plafonne à 0,64 F1 en \emph{mapping} \cite{ctibench} ; l'OT est un angle mort \cite{kong,kheddar} & Aucune évaluation sur de la télémétrie OT ; les bancs d'essai portent sur du texte de renseignement en anglais, et la contamination interdit de les réutiliser tels quels \\
\textbf{QR2} & Le RAG ancre et cite, mais un voisin n'est pas une définition, le résultat n'est pas reproductible, et cinq textes suffisent à 90\,\% de réussite d'attaque \cite{poisonedrag} & Pas de comparaison contrôlée RAG vs accès déterministe sur le \emph{mapping} ATT\&CK, à modèles et journaux identiques \\
\textbf{QR3} & Un modèle outillé de 6,7 G bat GPT-3 175 G \cite{toolformer} ; le MCP standardise l'accès déterministe sous gouvernance neutre \cite{mcpintro,aaif} ; au-delà d'une centaine d'outils exposés la sélection s'effondre (13,6\,\% contre 43,1\,\% avec pré-sélection \cite{ragmcp}), et l'appel à la demande économise 98\,\% des jetons \cite{mcpzero} & Les bancs d'essai MCP portent sur des tâches génériques \cite{mcpuniverse} ; la supériorité théorique de l'accès déterministe (justesse, reproductibilité, latence) n'est pas mesurée sur le \emph{mapping} ATT\&CK en conditions de SOC \\
\textbf{QR4} & Le MCP déplace le risque du texte à l'action ; les serveurs réels sont vulnérables (jusqu'à 72,8\,\% de réussite \cite{mcptox}, 5,5\,\% empoisonnés \cite{hasan}) ; la spécification laisse la lecture seule et la journalisation à l'implémenteur & Le « serveur MCP sûr par conception pour un SOC régulé » et son audit automatisé ancré dans ATLAS restent peu traités, a fortiori sous le cadre de la LPCE \\
\end{longtable}
\legendetableau{Ce que la revue établit et ce qu'elle laisse ouvert, par question de revue.}{tab:synthese-revue}

Trois fils se dégagent. \textbf{D'abord}, un motif transversal : à chaque axe, les modèles les plus capables sont les plus exposés --- au code non sûr (CyberSecEval), à l'excès de confiance mémorisée (\emph{inverse scaling}), à l'empoisonnement d'outil (MCPTox) --- parce que l'attaque exploite l'aptitude même qui fait leur valeur, suivre finement une instruction. La réponse n'est donc pas un meilleur modèle mais un meilleur \emph{cadre} : ancrer, borner, journaliser. \textbf{Ensuite}, l'arc de l'ancrage ne supprime pas la surface d'attaque, il la déplace --- du modèle vers la base (RAG), puis de la base vers le serveur (MCP) --- et chaque déplacement rend la surface plus petite et plus auditable, à condition que le serveur soit conçu pour l'être. \textbf{Enfin}, la littérature savante sur le MCP a moins de deux ans : elle mesure abondamment sa vulnérabilité mais très peu son apport en conditions opérationnelles. C'est dans cet écart que se situe l'essai ; la manière dont il répond à chacune de ces lacunes est exposée dans sa conclusion, à la lumière des travaux réalisés.

\section{Méthodologie d'évaluation}\label{sec:methodo}
Le projet intègre le référentiel MITRE ATT\&CK dans un écosystème MCP. La mission confiée au modèle est de produire des moyens de mitigation, et la qualité d'une mitigation dépend directement de la justesse du mapping qui la précède. On évalue donc les modèles sur cette tâche, aux trois niveaux du référentiel, dans deux environnements : les technologies de l'information (TI) et les technologies opérationnelles (TO).

Les travaux les plus proches --- RAM, CTIBench --- ont été publiés entre 2024 et le début de 2026 et portent sur des modèles aujourd'hui dépassés : GPT-4-Turbo, Llama 3, Mistral 7B, IBM Granite. Ils traitent par ailleurs de tâches voisines mais distinctes, comme le mapping de vulnérabilités CVE, la génération de règles SIEM ou la conversion de rapports de renseignement. Aucun n'évalue la matrice ATT\&CK for ICS, ni la version v19.2 du référentiel. C'est cette absence d'évaluation actualisée et alignée sur notre cas d'usage qui motive la campagne.

\FloatBarrier
\subsection{Construction du corpus}\label{sec:meth-corpus}

\subsubsection{Pourquoi générer les journaux}\label{sec:1-1-2}
Les journaux sont produits en laboratoire plutôt que repris d'un jeu public. La raison est la contamination : si un jeu de données est public, le modèle a pu voir pendant son entraînement les journaux \emph{et} leurs identifiants MITRE. Il répondrait alors de mémoire et non par raisonnement, et la campagne mesurerait sa mémorisation. En générant nous-mêmes les journaux, on garantit que leur forme exacte n'a jamais été indexée.

L'évaluation porte donc sur la capacité du modèle à mapper correctement des activités malveillantes \textbf{connues}, et non sur sa capacité à distinguer le légitime du malveillant.

\subsubsection{Génération}\label{sec:1-2-2}
Le laboratoire virtuel couvre deux zones. La zone TI comprend un contrôleur de domaine Windows Server DC01 instrumenté avec Sysmon, et un serveur Linux web01 instrumenté avec auditd. La zone TO est un réseau ICS composé d'une simulation de procédé GRFICS v3 et d'une sonde réseau Zeek. La zone TO est donc simulée, mais le trafic Modbus produit est conforme au protocole réel : c'est la représentativité que la tâche exige, sans les contraintes d'accès d'un réseau de production.

Des attaques unitaires ciblent chaque zone via un harnais qui produit les traces. Les résultats sont capturés dans les journaux Sysmon et Security pour DC01, dans audit.log pour web01, dans les journaux Zeek (Modbus) et le journal d'audit de l'hôte pour la zone TO. Un script découpe ensuite ces journaux en échantillons et construit la vérité de terrain avec des identifiants validés contre la version courante du référentiel.

\begin{figure}[htbp]
\centering
\figuretikz{fig-1-5-telemetrie}
\caption{Génération de la télémétrie TO et TI.}\label{fig:telemetrie}
\end{figure}

\subsubsection{Un triplet par échantillon}\label{sec:1-3}
Chaque échantillon correspond à une seule action offensive et porte exactement un triplet attendu : une tactique, une technique, et une sous-technique lorsque la technique en possède une, la mention « aucune » sinon.

Une technique peut être rattachée à plusieurs tactiques dans le référentiel --- T1053 relève à la fois d'Execution, de Persistence et de Privilege Escalation. Dans ce cas, on retient la seule tactique correspondant à l'intention réelle du scénario exécuté. L'attaque ayant été lancée par nous, cette intention est connue sans ambiguïté.

\subsubsection{Composition du corpus}\label{sec:1-4}
\begin{longtable}[]{@{}
  >{\raggedright\arraybackslash}p{(\columnwidth - 10\tabcolsep) * \real{0.0962}}
  >{\raggedright\arraybackslash}p{(\columnwidth - 10\tabcolsep) * \real{0.2084}}
  >{\raggedright\arraybackslash}p{(\columnwidth - 10\tabcolsep) * \real{0.1389}}
  >{\raggedright\arraybackslash}p{(\columnwidth - 10\tabcolsep) * \real{0.1229}}
  >{\raggedright\arraybackslash}p{(\columnwidth - 10\tabcolsep) * \real{0.1816}}
  >{\raggedright\arraybackslash}p{(\columnwidth - 10\tabcolsep) * \real{0.2520}}@{}}
\toprule\noalign{}
\begin{minipage}[b]{\linewidth}\raggedright
\textbf{Zone}
\end{minipage} & \begin{minipage}[b]{\linewidth}\raggedright
\textbf{Source}
\end{minipage} & \begin{minipage}[b]{\linewidth}\raggedright
\textbf{Matrice}
\end{minipage} & \begin{minipage}[b]{\linewidth}\raggedright
\textbf{Nombre}
\end{minipage} & \begin{minipage}[b]{\linewidth}\raggedright
\textbf{Avec sous-tech.}
\end{minipage} & \begin{minipage}[b]{\linewidth}\raggedright
\textbf{Techniques couvertes}
\end{minipage} \\
\midrule\noalign{}
\endhead
\bottomrule\noalign{}
\endlastfoot
TI & DC01 --- Sysmon + Security & Enterprise & 12 & 7 & 11 \\
TI & web01 --- audit.log & Enterprise & 16 & 14 & 14 \\
TO & Sonde Zeek --- Modbus & ICS & 11 & 1 & 11 \\
TO & Hôte PLC --- audit & ICS & 3 & 0 & 3 \\
\textbf{Total} & & & \textbf{42} & \textbf{22} & \textbf{36} \\
\end{longtable}
\legendetableau{Composition du corpus de la vérité de terrain.}{tab:corpus}


La vérité de terrain complète est reproduite ci-dessous, échantillon par échantillon. Elle fixe le triplet attendu à chaque niveau et sert de référence à tous les calculs qui suivent.

\begin{longtable}[]{@{}
  >{\raggedright\arraybackslash}p{(\columnwidth - 8\tabcolsep) * \real{0.0833}}
  >{\raggedright\arraybackslash}p{(\columnwidth - 8\tabcolsep) * \real{0.1111}}
  >{\raggedright\arraybackslash}p{(\columnwidth - 8\tabcolsep) * \real{0.2415}}
  >{\raggedright\arraybackslash}p{(\columnwidth - 8\tabcolsep) * \real{0.2778}}
  >{\raggedright\arraybackslash}p{(\columnwidth - 8\tabcolsep) * \real{0.2863}}@{}}
\toprule\noalign{}
\begin{minipage}[b]{\linewidth}\raggedright
\textbf{Éch.}
\end{minipage} & \begin{minipage}[b]{\linewidth}\raggedright
\textbf{Mat.}
\end{minipage} & \begin{minipage}[b]{\linewidth}\raggedright
\textbf{Tactique}
\end{minipage} & \begin{minipage}[b]{\linewidth}\raggedright
\textbf{Technique}
\end{minipage} & \begin{minipage}[b]{\linewidth}\raggedright
\textbf{Sous-technique}
\end{minipage} \\
\midrule\noalign{}
\endhead
\bottomrule\noalign{}
\endlastfoot
S000 & ICS & TA0102 --- Discovery & T0846 --- Remote System Discovery & aucune \\
S001 & ICS & TA0102 --- Discovery & T0888 --- Remote System Information Discovery & aucune \\
S002 & ICS & TA0100 --- Collection & T0861 --- Point \& Tag Identification & aucune \\
S003 & ICS & TA0100 --- Collection & T0801 --- Monitor Process State & aucune \\
S004 & ICS & TA0100 --- Collection & T0802 --- Automated Collection & aucune \\
S005 & ICS & TA0100 --- Collection & T0868 --- Detect Operating Mode & aucune \\
S006 & ICS & TA0106 --- Impair Process Control & T0836 --- Modify Parameter & aucune \\
S007 & ICS & TA0106 --- Impair Process Control & T1692 --- Unauthorized Message & T1692.001 --- Command Message \\
S008 & ICS & TA0106 --- Impair Process Control & T0806 --- Brute Force I/O & aucune \\
S009 & ICS & TA0107 --- Inhibit Response Function & T0835 --- Manipulate I/O Image & aucune \\
S010 & ICS & TA0107 --- Inhibit Response Function & T0814 --- Denial of Service & aucune \\
S011 & ICS & TA0104 --- Execution & T0807 --- Command-Line Interface & aucune \\
S012 & ICS & TA0104 --- Execution & T0853 --- Scripting & aucune \\
S013 & ICS & TA0110 --- Persistence & T0889 --- Modify Program & aucune \\
S014 & Ent. & TA0002 --- Execution & T1059 --- Command and Scripting Interpreter & T1059.004 --- Unix Shell \\
S015 & Ent. & TA0002 --- Execution & T1059 --- Command and Scripting Interpreter & T1059.006 --- Python \\
S016 & Ent. & TA0004 --- Privilege Escalation & T1548 --- Abuse Elevation Control Mechanism & T1548.003 --- Sudo and Sudo Caching \\
S017 & Ent. & TA0004 --- Privilege Escalation & T1548 --- Abuse Elevation Control Mechanism & T1548.001 --- Setuid and Setgid \\
S018 & Ent. & TA0003 --- Persistence & T1053 --- Scheduled Task/Job & T1053.003 --- Cron \\
S019 & Ent. & TA0003 --- Persistence & T1543 --- Create or Modify System Process & T1543.002 --- Systemd Service \\
S020 & Ent. & TA0003 --- Persistence & T1546 --- Event Triggered Execution & T1546.004 --- Unix Shell Configuration Modification \\
S021 & Ent. & TA0003 --- Persistence & T1098 --- Account Manipulation & T1098.004 --- SSH Authorized Keys \\
S022 & Ent. & TA0003 --- Persistence & T1037 --- Boot or Logon Initialization Scripts & T1037.004 --- RC Scripts \\
S023 & Ent. & TA0006 --- Credential Access & T1003 --- OS Credential Dumping & T1003.008 --- /etc/passwd and /etc/shadow \\
S024 & Ent. & TA0006 --- Credential Access & T1552 --- Unsecured Credentials & T1552.003 --- Shell History \\
S025 & Ent. & TA0005 --- Stealth & T1070 --- Indicator Removal & T1070.003 --- Clear Command History \\
S026 & Ent. & TA0005 --- Stealth & T1140 --- Deobfuscate/Decode Files or Information & aucune \\
S027 & Ent. & TA0112 --- Defense Impairment & T1685 --- Disable or Modify Tools & T1685.006 --- Clear Linux or Mac System Logs \\
S028 & Ent. & TA0112 --- Defense Impairment & T1222 --- File and Directory Permissions Modification & T1222.002 --- Linux and Mac Permissions \\
S029 & Ent. & TA0112 --- Defense Impairment & T1686 --- Disable or Modify System Firewall & aucune \\
S030 & Ent. & TA0007 --- Discovery & T1033 --- System Owner/User Discovery & aucune \\
S031 & Ent. & TA0007 --- Discovery & T1082 --- System Information Discovery & aucune \\
S032 & Ent. & TA0007 --- Discovery & T1057 --- Process Discovery & aucune \\
S033 & Ent. & TA0007 --- Discovery & T1016 --- System Network Configuration Discovery & aucune \\
S034 & Ent. & TA0007 --- Discovery & T1087 --- Account Discovery & T1087.001 --- Local Account \\
S035 & Ent. & TA0002 --- Execution & T1059 --- Command and Scripting Interpreter & T1059.001 --- PowerShell \\
S036 & Ent. & TA0002 --- Execution & T1059 --- Command and Scripting Interpreter & T1059.003 --- Windows Command Shell \\
S037 & Ent. & TA0005 --- Stealth & T1027 --- Obfuscated Files or Information & T1027.010 --- Command Obfuscation \\
S038 & Ent. & TA0003 --- Persistence & T1547 --- Boot or Logon Autostart Execution & T1547.001 --- Registry Run Keys / Startup Folder \\
S039 & Ent. & TA0003 --- Persistence & T1053 --- Scheduled Task/Job & T1053.005 --- Scheduled Task \\
S040 & Ent. & TA0003 --- Persistence & T1543 --- Create or Modify System Process & T1543.003 --- Windows Service \\
S041 & Ent. & TA0112 --- Defense Impairment & T1112 --- Modify Registry & aucune \\
\end{longtable}
\legendetableau{Vérité de terrain par échantillon : tactique, technique et sous-technique attendues.}{tab:verite-terrain}


\textbf{Contamination d'entraînement.} Les journaux sont rédigés pour la campagne et ne sont publiés nulle part : aucun modèle n'a pu les rencontrer à l'entraînement. Les scores mesurent donc l'analyse des événements, non la mémorisation du corpus.

\subsubsection{Échantillons, exécutions et workflow d'évaluation}\label{sec:1-5}
Chaque échantillon est soumis trois fois au même modèle. Un modèle de langue ne rend pas nécessairement la même réponse deux fois de suite : répéter l'appel permet de mesurer cette reproductibilité au lieu de la supposer. Un appel au modèle --- un échantillon envoyé, une réponse reçue --- est appelé une \textbf{exécution}. La campagne compte donc \textbf{42 échantillons et 126 exécutions} par modèle.

Ces exécutions sont produites par un workflow n8n, présenté à la figure~\ref{fig:workflow-eval}. Il se lit de gauche à droite, en cinq modules.

\begin{figure}[htbp]
\centering
\figuretikz{fig-1-2-workflow-eval}
\caption{Le workflow d'évaluation n8n de la phase 1 : cinq modules, des 42 échantillons aux fichiers de résultats JSONL. Fonctionnement détaillé en annexe~B.}\label{fig:workflow-eval}
\end{figure}

\textbf{Préparation.} Le workflow lit les 42 échantillons et dresse la liste des appels à faire : chaque échantillon, trois fois, soit 126 appels pour le modèle évalué.

\textbf{Boucle.} Les appels sont traités un par un, jamais en parallèle. Une courte pause sépare deux appels, pour rester sous les limites de débit imposées par les fournisseurs.

\textbf{Appel du LLM.} L'échantillon est envoyé au modèle avec la même consigne et le même format de réponse imposé pour les dix modèles. Seule la forme technique de la requête change d'un fournisseur à l'autre, car chaque API a ses propres règles.

\textbf{Jugement.} Le workflow extrait de la réponse le triplet tactique, technique et sous-technique, puis lui attribue un statut. Il sépare deux origines d'échec. Si la faute vient du modèle --- réponse hors format, JSON invalide ---, le statut est enregistré tel quel, et l'exécution comptera comme une abstention. Si la faute vient de l'infrastructure --- quota dépassé, délai dépassé, coupure réseau ---, le modèle n'a rien à se reprocher : le workflow attend, puis relance l'appel, jusqu'à cinq tentatives. Une exécution qui échoue encore après ces cinq tentatives est écartée du calcul, et rapportée à part.

\textbf{Écriture.} Chaque exécution devient une ligne du fichier de résultats du modèle, au format JSON : le triplet proposé, le statut, la latence et les jetons consommés. La ligne est ajoutée dès que l'appel est terminé ; une panne en cours de campagne ne fait donc perdre que l'appel en cours. Le fichier contient ainsi une ligne par exécution, et non par échantillon.

Les deux unités ne servent pas aux mêmes mesures. Les taux de diagnostic se comptent en \textbf{exécutions} : un comportement qui se répète sur les trois exécutions d'un même échantillon s'est produit trois fois, et doit compter trois fois. La stabilité se compte en \textbf{échantillons}, puisqu'elle compare les trois exécutions entre elles --- il n'y a rien à compter par exécution lorsque la question porte sur leur ensemble.

Un exemple. Si un modèle s'arrête à la tactique sur deux des trois exécutions d'un même échantillon, le compteur de sous-spécification retient 2, et non 1 : le comportement s'est produit deux fois sur les 126 occasions qu'il avait de se produire. Rapporter ce 2 aux 42 échantillons donnerait un taux trois fois trop élevé.

Une fois les 126 exécutions de chaque modèle enregistrées, les fichiers de résultats sont confiés à un script, \texttt{calculer\_metriques.py}, présenté à la figure~\ref{fig:workflow-metriques}. Le script reçoit deux entrées --- les fichiers de résultats et la vérité de terrain --- et rend deux sorties : un rapport texte, qui se lit directement, et un fichier CSV qui rassemble les tableaux de résultats repris dans ce chapitre. Pour plus d'information, voir l'annexe~B.
% Lien cliquable à activer une fois l'annexe B rédigée (y placer \label{annexe:B}) :
% Pour plus d'information, voir l'\hyperref[annexe:B]{annexe~B}.

\begin{figure}[htbp]
\centering
\figuretikz{fig-1-2b-workflow-metriques}
\caption{Du workflow d'évaluation au calcul des métriques : les fichiers de résultats JSONL et la vérité de terrain sont traités par le script \texttt{calculer\_metriques.py}, qui produit un rapport texte et un fichier CSV.}\label{fig:workflow-metriques}
\end{figure}

L'ensemble des étapes d'étiquetage et de calcul des métriques est réalisé par ce script, suivant la méthode exposée ci-dessous :

\FloatBarrier
\subsection{Étiquetage des prédictions}\label{sec:meth-etiquetage}

\subsubsection{Les trois niveaux du référentiel}\label{sec:2-1}
Le cadre MITRE ATT\&CK est hiérarchique : chaque technique appartient à une tactique, et certaines techniques se déclinent en sous-techniques. La figure~\ref{fig:hierarchie} l'illustre avec T1059 et ses variantes.

\begin{figure}[htbp]
\centering
\figuretikz{fig-1-6-hierarchie}
\caption{Hiérarchie à trois niveaux du référentiel, illustrée par la technique T1059.}\label{fig:hierarchie}
\end{figure}

Une erreur n'a pas la même gravité selon le niveau où elle se produit : confondre deux sous-techniques d'une même technique est une erreur fine, se tromper de tactique est une erreur grossière. La vérité de terrain est donc préparée aux trois niveaux, et chaque niveau est évalué séparément.

Le corpus compte 42 échantillons. Chacun porte un identifiant attendu par niveau. Le modèle rend un identifiant par niveau. On veut savoir à quel point il réussit.

Il n'y a pas une seule façon d'échouer. Le modèle peut donner un mauvais identifiant qui existe. Il peut en inventer un qui n'existe pas. Il peut donner un identifiant révoqué dont le successeur est exactement la bonne réponse. Il peut ne rien donner. Ces quatre échecs ne disent pas la même chose sur le modèle, et un score unique les écraserait.

D'où le principe : \textbf{on classe chaque prédiction dans une catégorie, avant tout calcul}, et les trois niveaux sont étiquetés et mesurés séparément.

\subsubsection{Les cinq catégories}\label{sec:2-2}
\begin{itemize}
\item
  \textbf{VP} --- l'identifiant attendu est prédit.
\item
  \textbf{FP} --- un identifiant qui existe, mais pas le bon. Attendu T1059.001, prédit T1547.001. Les FP ne se ressemblent pas tous --- mauvaise matrice, identifiant périmé, scission v19 --- et des compteurs dédiés, décrits plus bas, distinguent ces familles sans changer la case.
\item
  \textbf{H} (hallucination) --- un identifiant qui n'existe nulle part. T1059.999. Un identifiant \emph{déprécié} --- retiré du référentiel sans remplaçant officiel --- n'en est pas une : il existe. Il est compté FP et suivi par un compteur dédié ; faute de remplaçant officiel à vérifier, il ne peut pas être crédité O. Ce compteur couvre aussi les identifiants d'anciennes versions ou de matrices absorbées --- PRE-ATT\&CK, ATT\&CK for ICS v8 ---, absents du paquet figé mais attestés par les paquets historiques : eux non plus ne sont pas des inventions.
\item
  \textbf{O} (obsolète) --- un identifiant révoqué dont le remplaçant officiel est celui attendu. Prédit T1086, attendu T1059.001. Le raisonnement est juste, le référentiel du modèle date. C'est la seule catégorie dont le traitement change d'une base de calcul à l'autre. Elle n'est détectable qu'aux niveaux technique et sous-technique ; au niveau tactique, le cas est indiscernable et reste compté FP.
\item
  \textbf{A} (abstention) --- rien n'est prédit, le champ est vide. Une réponse inexploitable imputable au modèle --- format non respecté, JSON invalide, réponse vide --- devient une abstention aux trois niveaux, avec un compteur dédié ; les F1 sont aussi recalculés hors rejets. Le statut exact du rejet (\emph{format non respecté}, \emph{JSON invalide}) et l'exécution concernée sont conservés dans le fichier de résultats et rapportés nominativement dans l'évaluation.
\end{itemize}

Deux propriétés comptent. \textbf{Disjointes} : une prédiction reçoit une catégorie, une seule. \textbf{Exhaustives} : toute prédiction en reçoit une. Il en découle, sur les \emph{N} échantillons évalués à un niveau, l'égalité \emph{N = VP + FP + H + O + A}. C'est l'outil de contrôle : si elle ne tombe pas juste, il y a une erreur de comptage.

\textbf{Ce qu'on perdrait sans cette séparation.} En comptant seulement des bonnes et des mauvaises réponses, trois informations disparaîtraient. Le taux d'hallucination, H étant noyé dans les mauvaises --- or un modèle qui invente des identifiants est disqualifié pour un usage opérationnel, quel que soit son score. La base tolérante, O étant noyé de la même façon --- or la résolution des identifiants révoqués est précisément ce que la campagne veut mesurer. Et la distinction entre A et l'erreur : l'abstention renvoie le cas à l'analyste, l'erreur l'oriente vers une cible fausse.

\textbf{Pourquoi FN ne figure pas dans la liste.} Les cinq catégories décrivent ce que le modèle a écrit, et la liste est complète. Le faux négatif désigne autre chose : la réponse attendue n'a pas été retrouvée. Cela se produit dans quatre des cinq catégories --- FP, H, A, et O en base stricte. Ce n'est donc pas une catégorie de plus mais un regroupement qui traverse les autres ; l'ajouter romprait l'exclusivité mutuelle et l'égalité ci-dessus. Les catégories disent ce que le modèle a écrit, les regroupements du tableau~\ref{tab:bases} disent ce que ça vaut.

\subsubsection{Pourquoi le vrai négatif n'est pas comptabilisé}\label{sec:2-3}
La version figée du référentiel (v19.2) compte 794 techniques. Pour chaque échantillon, le modèle en écarte implicitement 793. En les comptant, l'exactitude atteindrait 99,9\,\% pour tous les modèles, y compris pour un modèle qui se tromperait systématiquement : l'indicateur cesserait de discriminer.

C'est pourquoi la précision et le rappel sont les métriques adaptées ici : \textbf{ni l'une ni l'autre n'utilise le vrai négatif}. Elles sont construites pour les cas où les négatifs écrasent numériquement les positifs.

\subsubsection{Le niveau sous-technique}\label{sec:2-4}
Les niveaux tactique et technique sont évalués sur les 42 échantillons. Le niveau sous-technique est le seul où le corpus se dédouble : \textbf{22 échantillons} possèdent une sous-technique attendue, \textbf{20 n'en possèdent aucune} --- leur technique est terminale. Le F1 sous-technique est calculé sur les 22, puisque la question posée est : lorsqu'une sous-technique existe, le modèle la retrouve-t-il ? Les 20 autres ne peuvent pas y contribuer, il n'y a rien à retrouver.

Sur les 42, un modèle répondant « aucune » partout sans analyse tomberait juste sur les 20 terminaux et obtiendrait un F1 de 0,476 --- la moitié du score maximal, sans aucune capacité. Ces cas sont triviaux, connus d'avance dès que la technique l'est ; les inclure comprimerait l'échelle vers le haut. Sur 22, ce modèle obtient zéro. La capacité à reconnaître une technique terminale reste mesurée, par le \textbf{taux de sur-spécification}, rapporté sur les 20.

\textbf{« aucune » n'est pas le vide.} Le critère est la forme de la réponse, pas la vérité attendue. « aucune » est une réponse : le modèle affirme que la technique n'a pas de sous-technique. C'est vérifiable, donc faux lorsqu'une sous-technique existait, et compté FP. Le champ vide est une absence de réponse : rien n'est affirmé, c'est une abstention. Confondre les deux fausse deux indicateurs différents --- traiter « aucune » comme une abstention mélange le silence et une prise de position ; traiter le vide comme « aucune » crédite le modèle d'une réponse qu'il n'a pas donnée. La consigne impose donc un champ sous\_technique toujours présent : un identifiant, la chaîne « aucune », ou une chaîne vide.

\begin{longtable}[]{@{}
  >{\raggedright\arraybackslash}p{(\columnwidth - 6\tabcolsep) * \real{0.2350}}
  >{\raggedright\arraybackslash}p{(\columnwidth - 6\tabcolsep) * \real{0.2350}}
  >{\raggedright\arraybackslash}p{(\columnwidth - 6\tabcolsep) * \real{0.1389}}
  >{\raggedright\arraybackslash}p{(\columnwidth - 6\tabcolsep) * \real{0.3911}}@{}}
\toprule\noalign{}
\begin{minipage}[b]{\linewidth}\raggedright
\textbf{Attendu}
\end{minipage} & \begin{minipage}[b]{\linewidth}\raggedright
\textbf{Prédit}
\end{minipage} & \begin{minipage}[b]{\linewidth}\raggedright
\textbf{Catégorie}
\end{minipage} & \begin{minipage}[b]{\linewidth}\raggedright
\textbf{Lecture}
\end{minipage} \\
\midrule\noalign{}
\endhead
\bottomrule\noalign{}
\endlastfoot
\multicolumn{4}{@{}l@{}}{%
\textbf{Les 22 échantillons évalués}} \\
T1059.001 & T1059.001 & \textbf{VP} & Réponse exacte. \\
T1059.001 & T1059.003 & \textbf{FP} & Confusion entre variantes. \\
T1059.001 & « aucune » & \textbf{FP} & Absence affirmée à tort ; compteur dédié. \\
T1059.001 & champ vide & \textbf{A} & Abstention ; alimente la sous-spécification N2. \\
T1059.001 & T1086 (révoqué) & \textbf{O} & Successeur officiel = réponse attendue. \\
T1059.001 & T1059.999 & \textbf{H} & Identifiant inexistant. \\
\multicolumn{4}{@{}l@{}}{%
\textbf{Les 20 échantillons terminaux --- hors du F1}} \\
aucune & « aucune » & \textbf{---} & Comportement attendu. \\
aucune & champ vide & \textbf{---} & Non compté. \\
aucune & une sous-technique & \textbf{---} & Sur-spécification ; compteur dédié, sur 20. \\
aucune & une sous-t. révoquée & \textbf{---} & Sur-spécification ; l'identifiant a existé, ce n'est pas une invention. \\
\end{longtable}
\legendetableau{Étiquetage du niveau sous-technique.}{tab:sous-technique}


\subsubsection{Attributs de diagnostic}\label{sec:2-5}
Ces indicateurs ne sont pas des catégories : ils se superposent à la classification et expliquent \emph{comment} le modèle échoue, là où le F1 mesure \emph{combien} il réussit. Ils se comptent en exécutions --- chaque échantillon étant soumis trois fois ---, sauf la stabilité, qui se compte en échantillons.

\begin{itemize}
\item
  \textbf{Sous-spécification N1} = exécutions où la tactique est fournie sans technique / exécutions évaluées. Le modèle a compris l'intention, pas le mécanisme.
\item
  \textbf{Sous-spécification N2} = exécutions où la technique est correcte et la sous-technique omise / exécutions portant sur un échantillon à sous-technique. Bon mécanisme, granularité insuffisante pour des mitigations précises.
\item
  \textbf{Sur-spécification} = exécutions affirmant une sous-technique sur une technique terminale / exécutions portant sur un échantillon terminal. L'inverse du précédent : inventer de la précision.
\item
  \textbf{Sous-technique niée à tort} = nombre d'exécutions répondant « aucune » alors qu'une sous-technique existait. Effectif brut, en complément du F1.
\item
  \textbf{Mauvaise matrice} = nombre d'identifiants prédits hors de la matrice de l'échantillon, les trois niveaux confondus. Effectif brut.
\item
  \textbf{Scission v19} = nombre d'exécutions proposant TA0005 là où TA0112 est attendu. La v19 a scindé la tactique Defense Evasion ; ce compteur mesure l'intégration d'un changement de nomenclature daté et unique, que ni la base tolérante ni la table des révocations ne rattrapent --- TA0005 reste vivant ailleurs, aucun lien officiel ne le relie à TA0112.
\item
  \textbf{Stabilité} = échantillons dont toutes les exécutions rendent un triplet identique / échantillons évalués au moins deux fois. Mesure la reproductibilité, indépendamment de la justesse.
\end{itemize}

\textbf{Cas exclu du calcul.} Un échec imputable à l'infrastructure d'appel --- quota épuisé, requête annulée, délai dépassé, coupure réseau --- ne mesure pas le modèle. Le harnais rejoue l'appel avec attente exponentielle ; toute itération abandonnée malgré les tentatives est marquée (erreur\_infrastructure) puis écartée, son effectif étant rapporté par modèle. Ces itérations ne deviennent jamais des abstentions : l'abstention suppose que le modèle a reçu la question et n'a rien rendu d'exploitable, alors qu'ici il ne l'a jamais reçue. Les deux familles ne se croisent pas --- l'abstention pèse dans FN, l'itération écartée sort de N.

\FloatBarrier
\subsection{Calcul des métriques}\label{sec:meth-metriques}

\subsubsection{Deux questions, deux populations}\label{sec:3-1}
La précision et le rappel ne sont pas deux versions du même chiffre. Ils portent sur deux populations.

\begin{itemize}
\item
  \textbf{Précision} --- population : les prédictions émises. Question : combien étaient justes ? Dénominateur \emph{VP + FP + H + O}. Les abstentions n'y sont pas, puisque rien n'a été émis.
\item
  \textbf{Rappel} --- population : les échantillons évalués. Question : sur combien le modèle a-t-il produit la bonne réponse ? Dénominateur \emph{N}. Les abstentions y sont.
\end{itemize}

Deux dénominateurs différents, parce que deux questions différentes. Il en découle la \textbf{règle d'appariement} : une prédiction fausse compte comme échec au rappel, au même titre qu'une absence de prédiction. Dans les deux cas, l'échantillon n'a pas été qualifié.

\textbf{Pourquoi cette règle.} Un modèle qui répond aux 42 échantillons et se trompe 36 fois a 6 bonnes réponses : son rappel vaut 6/42, soit 14\,\%. Si seules les abstentions comptaient comme échec, ce modèle n'en aurait aucune et son rappel vaudrait 6/6, soit 100\,\% --- un modèle qui échoue 86\,\% du temps serait crédité d'un rappel parfait. C'est la seule règle qui empêche d'améliorer son rappel en devinant. Une prédiction fausse pénalise donc la précision, parce que le modèle a affirmé du faux, et le rappel, parce que l'échantillon reste non qualifié : ce ne sont pas deux sanctions pour une même faute, mais deux questions auxquelles la réponse est négative.

\subsubsection{Les deux bases}\label{sec:3-2}
Un mapping juste mais daté ne doit être confondu ni avec une réussite pleine ni avec une erreur ordinaire. La \textbf{base stricte} prend le référentiel courant pour juge : l'obsolète est une erreur, puisqu'un analyste qui le recevrait chercherait une entrée qui n'existe plus. La \textbf{base tolérante} prend le raisonnement pour juge : l'obsolète est crédité, puisque la table de correspondance MITRE atteste que la réponse désignait la bonne réalité. La base stricte fournit les métriques principales ; l'écart entre les deux isole la part des erreurs imputable à la fraîcheur des connaissances.

\begin{longtable}[]{@{}
  >{\raggedright\arraybackslash}p{(\columnwidth - 6\tabcolsep) * \real{0.2179}}
  >{\raggedright\arraybackslash}p{(\columnwidth - 6\tabcolsep) * \real{0.2607}}
  >{\raggedright\arraybackslash}p{(\columnwidth - 6\tabcolsep) * \real{0.2607}}
  >{\raggedright\arraybackslash}p{(\columnwidth - 6\tabcolsep) * \real{0.2606}}@{}}
\toprule\noalign{}
\begin{minipage}[b]{\linewidth}\raggedright
\end{minipage} & \begin{minipage}[b]{\linewidth}\raggedright
\textbf{Numérateur}
\end{minipage} & \begin{minipage}[b]{\linewidth}\raggedright
\textbf{Dénom. précision}
\end{minipage} & \begin{minipage}[b]{\linewidth}\raggedright
\textbf{Dénom. rappel}
\end{minipage} \\
\midrule\noalign{}
\endhead
\bottomrule\noalign{}
\endlastfoot
\textbf{Base stricte} & \emph{VP} & \emph{VP + FP + H + O} & \emph{N} \\
\textbf{Base tolérante} & \emph{VP + O} & \emph{VP + FP + H + O} & \emph{N} \\
\end{longtable}
\legendetableau{Ce qui change d'une base à l'autre.}{tab:bases}


Une seule différence entre les deux lignes : \emph{O} passe au numérateur. Les dénominateurs sont identiques --- le modèle n'a pas émis plus ou moins de prédictions, et le corpus n'a pas changé de taille. C'est cette invariance qui rend l'écart entre les deux bases interprétable : il mesure \emph{O}, et rien d'autre.

\textbf{Pourquoi un tableau plutôt qu'une réaffectation.} On écrit souvent, pour passer d'une base à l'autre, des affectations du type « FP devient FP + O ». Le symbole y désigne deux choses dans une même ligne, et la formule qui le réutilise ensuite devient ambiguë : selon la lecture, l'obsolète est compté une fois ou deux, et le résultat change. Ici les catégories ne bougent jamais ; seule la façon de les regrouper change, et chaque symbole garde un sens unique dans tout le document.

\subsubsection{Précision, rappel et F1}\label{sec:3-3}
\textbf{Précision} stricte = \emph{VP} / (\emph{VP + FP + H + O}) \textbf{Rappel} strict = \emph{VP} / \emph{N}

\textbf{Précision} tolérante = (\emph{VP + O}) / (\emph{VP + FP + H + O}) \textbf{Rappel} tolérant = (\emph{VP + O}) / \emph{N}

\textbf{F1} = 2 $\times$ Précision $\times$ Rappel / (Précision + Rappel), dans chaque base.

\textbf{Le calcul, en français simple.} Le calcul part des cinq cases de l'étiquetage --- VP, FP, H, O, A --- et se refait à chaque niveau : tactique, technique, sous-technique. Les lettres des cases gardent leur sens : le calcul ne fait que les regrouper.

On fixe d'abord, selon la base, ce qui compte comme \emph{réponse juste} et ce qui compte comme \emph{affirmation fausse}. En base stricte, les réponses justes sont les VP, et les affirmations fausses regroupent FP, H et O. En base tolérante, l'obsolète change de camp, puisque son remplaçant officiel est la bonne réponse : les réponses justes regroupent VP et O, les affirmations fausses FP et H. Rien d'autre ne bouge d'une base à l'autre.

Vient ensuite le FN, le faux négatif : les cibles manquées. L'étiquetage ne lui a pas donné de case, pour une raison simple : la bonne réponse manque dans plusieurs cases à la fois --- quand le modèle se trompe, quand il invente, quand il se tait --- et une case FN compterait deux fois les mêmes prédictions. On le reconstruit donc ici, au moment où le calcul en a besoin : les cibles manquées sont les affirmations fausses plus les abstentions. Une prédiction fausse compte donc des deux côtés --- elle affirme une erreur, et elle laisse l'échantillon sans réponse ---, et c'est voulu : deviner ne rapporte jamais.

Les trois formules ferment le calcul, identiques dans les deux bases. La précision mesure la fiabilité des affirmations : réponses justes, divisées par toutes les affirmations émises --- justes et fausses. Le rappel mesure la couverture du corpus : réponses justes, divisées par réponses justes plus cibles manquées ; et comme cette somme vaut N, le rappel est simplement la part des N exécutions résolues. Le F1 est la moyenne harmonique des deux : il ne monte que si la précision et le rappel montent ensemble.

Trois égalités permettent de contrôler n'importe quel tableau de résultats sans refaire l'analyse : réponses justes + cibles manquées = \emph{N} ; cibles manquées = affirmations fausses + \emph{A} ; réponses justes + affirmations fausses = nombre d'affirmations émises, identique dans les deux bases.

\textbf{Une propriété à connaître.} Dans chaque base, précision et rappel coïncident dès que le modèle se prononce sur tous les échantillons : les deux rapports partagent alors le même dénominateur. Seule l'abstention les fait diverger --- elle pèse au rappel, pas à la précision. Le passage en base tolérante déplace les deux indicateurs ensemble, il ne les sépare pas. C'est une conséquence arithmétique du cadre, non un défaut.

\subsubsection{Exemple de vérification}\label{sec:3-4}
Dix échantillons à un niveau donné : \emph{VP} = 5, \emph{FP} = 2, \emph{H} = 1, \emph{O} = 1, \emph{A} = 1, soit 9 affirmations émises et \emph{N} = 10.

\begin{longtable}[]{@{}
  >{\raggedright\arraybackslash}p{(\columnwidth - 12\tabcolsep) * \real{0.1603}}
  >{\raggedright\arraybackslash}p{(\columnwidth - 12\tabcolsep) * \real{0.1282}}
  >{\raggedright\arraybackslash}p{(\columnwidth - 12\tabcolsep) * \real{0.1282}}
  >{\raggedright\arraybackslash}p{(\columnwidth - 12\tabcolsep) * \real{0.1282}}
  >{\raggedright\arraybackslash}p{(\columnwidth - 12\tabcolsep) * \real{0.1517}}
  >{\raggedright\arraybackslash}p{(\columnwidth - 12\tabcolsep) * \real{0.1517}}
  >{\raggedright\arraybackslash}p{(\columnwidth - 12\tabcolsep) * \real{0.1517}}@{}}
\toprule\noalign{}
\begin{minipage}[b]{\linewidth}\raggedright
\end{minipage} & \begin{minipage}[b]{\linewidth}\raggedright
\textbf{Justes}
\end{minipage} & \begin{minipage}[b]{\linewidth}\raggedright
\textbf{Fausses}
\end{minipage} & \begin{minipage}[b]{\linewidth}\raggedright
\textbf{Manquées}
\end{minipage} & \begin{minipage}[b]{\linewidth}\raggedright
\textbf{Précision}
\end{minipage} & \begin{minipage}[b]{\linewidth}\raggedright
\textbf{Rappel}
\end{minipage} & \begin{minipage}[b]{\linewidth}\raggedright
\textbf{F1}
\end{minipage} \\
\midrule\noalign{}
\endhead
\bottomrule\noalign{}
\endlastfoot
\textbf{Stricte} & 5 & 4 & 5 & 5/9 = 0,556 & 5/10 = 0,500 & 0,526 \\
\textbf{Tolérante} & 6 & 3 & 4 & 6/9 = 0,667 & 6/10 = 0,600 & 0,632 \\
\end{longtable}
\legendetableau{Application numérique dans les deux bases.}{tab:application}


Les trois égalités se vérifient : 5 + 5 = 10 et 6 + 4 = 10 ; 5 = 4 + 1 et 4 = 3 + 1 ; 9 affirmations sur les deux lignes. L'écart de F1, 0,632 moins 0,526, mesure ici le coût de la seule prédiction obsolète.

\subsubsection{Indicateurs dérivés}\label{sec:3-5}
\begin{itemize}
\item
  \textbf{Écart d'obsolescence} = F1 tolérant - F1 strict.
\item
  \textbf{Taux d'obsolescence} = \emph{O} / affirmations émises, les trois niveaux confondus. Le référentiel interne du modèle est-il à jour ?
\item
  \textbf{Taux d'hallucination} = \emph{H} / affirmations émises, les trois niveaux confondus. Le modèle fabrique-t-il des identifiants ?
\item
  \textbf{Taux d'abstention} = \emph{A} / \emph{N}, rapporté par niveau. Un modèle peut être précis parce qu'il ne se prononce que lorsqu'il est sûr : ce taux complète la lecture de la précision.
\item
  \textbf{Latence} = horodatage de fin - horodatage de début de l'appel réussi ; les pauses du harnais et les attentes de reprise n'en font pas partie. Elle produit une valeur par exécution, soit 126 exécutions (3 exécutions $\times$ 42 échantillons). Les latences sont triées par ordre croissant, \emph{v{[}0{]}} étant la plus rapide. On utilise la convention d'interpolation de NumPy pour calculer le p50 et le p95 : le rang cherché vaut \emph{k = (n - 1) $\times$ q} et, en notant \emph{i} la partie entière de \emph{k}, la valeur est \emph{v{[}i{]} + (v{[}i+1{]} - v{[}i{]}) $\times$ (k - i)}.
\item
  \textbf{p50 (médiane)} --- \emph{q = 0,50}, soit \emph{k = 125 $\times$ 0,50 = 62,5} : la valeur se lit à mi-chemin entre la 63e et la 64e latence. La moitié des exécutions sont plus rapides : c'est le temps de traitement courant du modèle.
\item
  \textbf{p95} --- \emph{q = 0,95}, soit \emph{k = 125 $\times$ 0,95 = 118,75} : la valeur se lit aux trois quarts entre la 119e et la 120e latence. 95\,\% des exécutions sont plus rapides : c'est le pire cas ordinaire, celui qui dimensionne un usage en exploitation et qu'une moyenne masquerait.
\item
  \textbf{Jetons} = jetons d'entrée + jetons de sortie, jetons de raisonnement compris lorsque le fournisseur les facture. Rapportés en moyenne par exécution.
\item
  \textbf{Efficacité} = 1000 $\times$ F1 technique strict / jetons moyens par exécution, soit le F1 obtenu par millier de jetons consommés. Le F1 technique est retenu car les mitigations MITRE se rattachent aux techniques.
\end{itemize}

\subsubsection{Ventilation des résultats}\label{sec:3-6}
Chaque taux est rapporté avec le dénominateur fixé lors de sa définition --- sections~\ref{sec:2-4} et~\ref{sec:2-5} --- afin qu'aucune proportion ne soit lue sans son effectif. Toutes les métriques sont produites globalement puis ventilées par matrice --- Enterprise, 28 échantillons ; ICS, 14 --- avec leurs effectifs bruts. Une différence entre deux modèles n'est retenue que si elle est nette et constante sur les trois exécutions. Répéter une campagne complète à l'identique fait par ailleurs varier les F1 d'environ $\pm$0,008 : ce grain sert de seuil de lecture, tout écart inférieur relève du bruit. En particulier, deux F1 séparés de moins de ce grain sont rapportés comme un ex æquo de fait, l'ordre des lignes d'un tableau n'étant alors qu'une convention d'affichage ; un rééchantillonnage \emph{bootstrap} des 42 échantillons pourra compléter ce seuil si un départage plus fin devient nécessaire. Enfin, le F1 sous-technique de la matrice ICS, calculé sur trois échantillons seulement, est rapporté pour mémoire sans être interprété.

Ce cadre est la micro-agrégation usuelle en classification multi-classes à étiquette unique. Il est vérifiable par les trois égalités, non ambigu puisque chaque symbole garde un sens unique, et complet puisque les quatre modes d'échec --- confusion, fabrication, obsolescence, abstention --- restent mesurables séparément.

\FloatBarrier
\subsection{Sélection des modèles}\label{sec:meth-modeles}

\subsubsection{Critères de sélection}\label{sec:4-1}
Dix modèles sont retenus, issus de six éditeurs. Le choix suit cinq critères, tous liés à ce que la campagne cherche à mesurer.

\begin{itemize}
\item
  \textbf{Récence} --- les travaux voisins évaluent des modèles dépassés. On retient la génération courante, seule pertinente pour une recommandation opérationnelle.
\item
  \textbf{Deux familles} --- les modèles propriétaires correspondent à ce qu'un exploitant utilise aujourd'hui par API ; les modèles à poids ouverts à ce qu'il pourrait héberger lui-même, question qui se pose dès que des journaux d'infrastructure critique quittent le périmètre.
\item
  \textbf{Couverture des éditeurs et des gammes} --- six éditeurs sont représentés, ce qui évite de mesurer une seule lignée d'entraînement. Pour quatre d'entre eux, deux points de gamme sont évalués : un modèle de tête et un modèle plus léger de la même lignée. Comparer deux modèles d'un même éditeur isole l'effet de la taille à entraînement constant, ce qu'une comparaison entre éditeurs ne permet pas.
\item
  \textbf{Capacité de raisonnement} --- le mapping n'est pas une reconnaissance de mots-clés mais un raisonnement contextuel. On retient donc des modèles de la classe raisonnement, ce que la revue de littérature justifie.
\item
  \textbf{Appel d'outils} --- la phase 2 rejoue la campagne avec le serveur MCP. Les dix modèles doivent donc gérer l'appel d'outils, sans quoi la comparaison avec et sans ancrage serait impossible sur le même panel.
\end{itemize}

\subsubsection{Les dix modèles retenus}\label{sec:4-2}
\begin{longtable}[]{@{}
  >{\raggedright\arraybackslash}p{(\columnwidth - 6\tabcolsep) * \real{0.2137}}
  >{\raggedright\arraybackslash}p{(\columnwidth - 6\tabcolsep) * \real{0.1624}}
  >{\raggedright\arraybackslash}p{(\columnwidth - 6\tabcolsep) * \real{0.4017}}
  >{\raggedright\arraybackslash}p{(\columnwidth - 6\tabcolsep) * \real{0.2222}}@{}}
\toprule\noalign{}
\begin{minipage}[b]{\linewidth}\raggedright
\textbf{Modèle}
\end{minipage} & \begin{minipage}[b]{\linewidth}\raggedright
\textbf{Éditeur}
\end{minipage} & \begin{minipage}[b]{\linewidth}\raggedright
\textbf{Identifiant d'API}
\end{minipage} & \begin{minipage}[b]{\linewidth}\raggedright
\textbf{Accès}
\end{minipage} \\
\midrule\noalign{}
\endhead
\bottomrule\noalign{}
\endlastfoot
\multicolumn{4}{@{}l@{}}{%
\textbf{Modèles propriétaires}} \\
Claude Opus 4.8 & Anthropic & claude-opus-4-8 & API Anthropic \\
Claude Sonnet 5 & Anthropic & claude-sonnet-5 & API Anthropic \\
GPT-5.6 Sol & OpenAI & gpt-5.6-sol & API OpenAI \\
GPT-5.6 Luna & OpenAI & gpt-5.6-luna & API OpenAI \\
Gemini 3.1 Pro & Google & google/gemini-3.1-pro-preview & Vertex AI \\
Gemini 3.8 Flash & Google & google/gemini-3.8-flash & Vertex AI \\
\multicolumn{4}{@{}l@{}}{%
\textbf{Modèles à poids ouverts}} \\
Kimi K3 & Moonshot AI & kimi-k3 & API Moonshot \\
DeepSeek V4 Pro & DeepSeek & deepseek-v4-pro & API DeepSeek \\
DeepSeek V4 Flash & DeepSeek & deepseek-v4-flash & API DeepSeek \\
Qwen3.8 & Alibaba & qwen3.8-2.4t-a95b & Model Studio \\
\end{longtable}
\legendetableau{Modèles évalués, éditeurs et voies d'accès.}{tab:modeles}


Les six modèles propriétaires proviennent de trois éditeurs distincts et représentent l'état de l'art accessible par API au moment de la campagne. Les quatre modèles à poids ouverts sont interrogés par \textbf{les API de premier niveau de leurs éditeurs} --- Moonshot AI, DeepSeek, Alibaba Cloud --- et non par une passerelle d'inférence. Ce choix écarte les écarts de quantification et de réglage de service propres aux passerelles : les dix modèles sont servis dans des conditions homogènes, ce qui est la condition d'une comparaison entre familles.

Les identifiants d'API sont ceux transmis par le harnais et sont relevés dans chaque enregistrement de résultat, avec la date de la campagne. Cette trace est nécessaire : les identifiants évoluent, et \emph{gemini-3.1-pro-preview} porte la mention preview, donc peut changer ou disparaître. Le préfixe \emph{google/} des deux modèles Gemini n'appartient pas au modèle : il est imposé par l'endpoint compatible OpenAI de Vertex AI. Un accès par clé AI Studio les désignerait sans préfixe.

\subsubsection{Conditions d'appel}\label{sec:4-3}
Les conditions sont identiques pour les dix modèles là où les API le permettent : même consigne système, même format de sortie imposé, même plafond de sortie de 32 768 jetons, trois exécutions par échantillon. Le plafond est assez haut pour ne contraindre aucun modèle sur cette tâche ; le champ fin\_raison enregistré à chaque appel permet de le vérifier après coup --- aucun enregistrement ne doit porter la mention de troncature.

\textbf{La température fait exception.} Elle est fixée à 0 lorsque le paramètre est accepté, et laissée à sa valeur par défaut là où les API l'imposent ou l'ignorent : Anthropic et Moonshot n'acceptent plus de température configurable sur leurs modèles à raisonnement, et le schéma OpenAI ne la transmet pas. La valeur effectivement appliquée est enregistrée à chaque appel. La stabilité inter-exécutions se lit donc comme la reproductibilité du modèle dans ses conditions de service réelles, et non comme un test de déterminisme à température nulle.

L'effort de raisonnement est laissé aux valeurs par défaut de chaque fournisseur. Aucune échelle commune n'existe entre éditeurs : forcer un niveau « maximal » partout ne rendrait pas les modèles comparables, cela les déplacerait vers des points arbitraires différents. Les jetons de raisonnement consommés sont en revanche relevés et rapportés, ce qui rend l'écart mesurable au lieu d'être supposé.

\subsubsection{Limites assumées}\label{sec:4-4}
\begin{itemize}
\item
  La température n'est pas contrôlable sur tous les modèles : cinq des dix --- les deux Claude, les deux GPT et Kimi K3 --- sont évalués à la valeur par défaut de leur fournisseur.
\item
  L'effort de raisonnement n'est pas égalisé entre éditeurs, faute d'échelle commune.
\item
  La dégradation en contexte long n'est pas mesurée : avec 42 échantillons par modèle, les effectifs par tranche de longueur seraient trop faibles pour conclure.
\item
  Deux éditeurs sur six ne sont représentés que par un modèle --- Moonshot et Alibaba : l'effet de gamme n'y est pas mesurable.
\end{itemize}

\FloatBarrier
\subsection{Évaluation et synthèse}\label{sec:meth-evaluation}

Les tableaux qui suivent comparent les modèles sur deux axes : la qualité du mapping et l'actualité des connaissances d'abord, rapportées dans les deux bases de calcul ; la fiabilité, le temps et le coût ensuite. Aucun critère n'est éliminatoire : la mauvaise matrice et l'hallucination, les deux fautes les plus graves, sont placées en tête de colonnes comme des réserves qui se lisent avec le F1, sans retrancher aucun modèle des tableaux. Les lignes sont ordonnées par F1 technique ; deux F1 séparés de moins de $\pm$0,008 --- l'écart observé en rejouant deux fois la même campagne --- sont un ex æquo de fait ; et le même ordre de colonnes est appliqué aux deux bases, si bien que l'écart entre les deux tableaux mesure une seule chose : l'actualité du référentiel interne de chaque modèle.

\subsubsection{Étiquetage des prédictions}\label{sec:5-1}

Rappelons d'abord pourquoi cette évaluation classe chaque réponse dans une catégorie, au lieu de compter seulement les bonnes et les mauvaises réponses. Le corpus compte 42 échantillons ; chacun est soumis trois fois au même modèle, soit 126 exécutions par modèle. Si l'on se contentait de deux cases --- bonne réponse, mauvaise réponse ---, les identifiants inventés seraient noyés parmi les mauvaises réponses : le taux d'hallucination disparaîtrait du calcul, alors qu'un modèle qui fabrique des identifiants est inutilisable en exploitation, quel que soit son score. Les identifiants obsolètes seraient noyés de la même façon, et la base tolérante ne pourrait plus être calculée. C'est pour cela que chaque réponse est d'abord rangée dans une case ; le tableau~\ref{tab:etiquetage} en donne le résultat, modèle par modèle et niveau par niveau : les cinq catégories --- VP, FP, H, O, A --- et deux sous-comptes inclus dans FP : F, les identifiants d'anciennes versions, et MM, les identifiants pris hors de la matrice de l'échantillon.

\begin{longtable}[]{@{}
  >{\raggedright\arraybackslash}p{(\columnwidth - 16\tabcolsep) * \real{0.2100}}
  >{\raggedright\arraybackslash}p{(\columnwidth - 16\tabcolsep) * \real{0.2400}}
  >{\raggedright\arraybackslash}p{(\columnwidth - 16\tabcolsep) * \real{0.0786}}
  >{\raggedright\arraybackslash}p{(\columnwidth - 16\tabcolsep) * \real{0.0786}}
  >{\raggedright\arraybackslash}p{(\columnwidth - 16\tabcolsep) * \real{0.0786}}
  >{\raggedright\arraybackslash}p{(\columnwidth - 16\tabcolsep) * \real{0.0786}}
  >{\raggedright\arraybackslash}p{(\columnwidth - 16\tabcolsep) * \real{0.0786}}
  >{\raggedright\arraybackslash}p{(\columnwidth - 16\tabcolsep) * \real{0.0784}}
  >{\raggedright\arraybackslash}p{(\columnwidth - 16\tabcolsep) * \real{0.0786}}@{}}
\toprule\noalign{}
\begin{minipage}[b]{\linewidth}\raggedright
\textbf{Modèle}
\end{minipage} & \begin{minipage}[b]{\linewidth}\raggedright
\textbf{Niveau}
\end{minipage} & \begin{minipage}[b]{\linewidth}\raggedright
\textbf{VP}
\end{minipage} & \begin{minipage}[b]{\linewidth}\raggedright
\textbf{FP}
\end{minipage} & \begin{minipage}[b]{\linewidth}\raggedright
\textbf{H}
\end{minipage} & \begin{minipage}[b]{\linewidth}\raggedright
\textbf{F}
\end{minipage} & \begin{minipage}[b]{\linewidth}\raggedright
\textbf{O}
\end{minipage} & \begin{minipage}[b]{\linewidth}\raggedright
\textbf{A}
\end{minipage} & \begin{minipage}[b]{\linewidth}\raggedright
\textbf{MM}
\end{minipage} \\
\midrule\noalign{}
\endhead
\bottomrule\noalign{}
\endlastfoot
\textbf{Claude Opus 4.8} & Tactique (n = 126) & 74 & 52 & 0 & 0 & 0 & 0 & 0 \\
& Technique (n = 126) & 71 & 54 & 0 & 0 & 1 & 0 & 0 \\
& Sous-technique (n = 66) & 42 & 21 & 0 & 0 & 3 & 0 & 0 \\
\textbf{Claude Sonnet 5} & Tactique (n = 126) & 64 & 60 & 0 & 0 & 0 & 2 & 0 \\
& Technique (n = 126) & 67 & 54 & 0 & 0 & 3 & 2 & 0 \\
& Sous-technique (n = 66) & 40 & 22 & 0 & 0 & 3 & 1 & 0 \\
\textbf{GPT-5.6 Sol} & Tactique (n = 126) & 82 & 44 & 0 & 0 & 0 & 0 & 0 \\
& Technique (n = 126) & 76 & 48 & 0 & 6 & 2 & 0 & 0 \\
& Sous-technique (n = 66) & 42 & 21 & 0 & 0 & 3 & 0 & 0 \\
\textbf{GPT-5.6 Luna} & Tactique (n = 126) & 63 & 49 & 13 & 0 & 0 & 1 & 4 \\
& Technique (n = 126) & 66 & 57 & 0 & 0 & 2 & 1 & 5 \\
& Sous-technique (n = 66) & 42 & 21 & 0 & 0 & 3 & 0 & 0 \\
\textbf{Gemini 3.1 Pro} & Tactique (n = 125) & 70 & 54 & 1 & 0 & 0 & 0 & 0 \\
& Technique (n = 125) & 60 & 62 & 0 & 3 & 3 & 0 & 1 \\
& Sous-technique (n = 65) & 42 & 20 & 0 & 0 & 3 & 0 & 0 \\
\textbf{Gemini 3.8 Flash} & Tactique (n = 125) & 81 & 44 & 0 & 0 & 0 & 0 & 0 \\
& Technique (n = 125) & 75 & 47 & 0 & 1 & 3 & 0 & 0 \\
& Sous-technique (n = 66) & 44 & 19 & 0 & 0 & 3 & 0 & 0 \\
\textbf{Kimi K3} & Tactique (n = 126) & 68 & 57 & 1 & 0 & 0 & 0 & 3 \\
& Technique (n = 126) & 71 & 53 & 0 & 1 & 2 & 0 & 0 \\
& Sous-technique (n = 66) & 41 & 22 & 0 & 0 & 3 & 0 & 0 \\
\textbf{DeepSeek V4 Pro} & Tactique (n = 126) & 81 & 43 & 2 & 0 & 0 & 0 & 4 \\
& Technique (n = 126) & 72 & 52 & 0 & 1 & 2 & 0 & 1 \\
& Sous-technique (n = 66) & 44 & 19 & 0 & 0 & 3 & 0 & 0 \\
\textbf{DeepSeek V4 Flash} & Tactique (n = 126) & 75 & 48 & 3 & 0 & 0 & 0 & 3 \\
& Technique (n = 126) & 76 & 48 & 0 & 0 & 2 & 0 & 1 \\
& Sous-technique (n = 66) & 41 & 22 & 0 & 0 & 3 & 0 & 0 \\
\textbf{Qwen3.8} & Tactique (n = 126) & 62 & 62 & 2 & 1 & 0 & 0 & 6 \\
& Technique (n = 126) & 72 & 54 & 0 & 4 & 0 & 0 & 1 \\
& Sous-technique (n = 66) & 42 & 21 & 0 & 0 & 3 & 0 & 0 \\
\end{longtable}
\legendetableau{Étiquetage des prédictions des dix modèles. Entre parenthèses, le nombre d'exécutions évaluées au niveau considéré. Chaque exécution reçoit une catégorie et une seule, donc chaque ligne vérifie N = VP + FP + H + O + A. F et MM sont des sous-comptes déjà inclus dans FP, affichés pour lecture. Les effectifs de 125 des deux Gemini correspondent aux deux itérations abandonnées sur quota d'API après cinq tentatives (S015, exécution 2, pour 3.1 Pro ; S013, exécution 3, pour 3.8 Flash) : des pannes du fournisseur, écartées du calcul et jamais comptées comme des abstentions du modèle.}{tab:etiquetage}

\textbf{Hallucinations.} Six modèles ont halluciné --- DeepSeek V4 Flash (3), DeepSeek V4 Pro (2), Gemini 3.1 Pro (1), GPT-5.6 Luna (13), Kimi K3 (1) et Qwen3.8 (2) ---, uniquement au niveau tactique et uniquement sur des journaux industriels : les identifiants fabriqués (TA0113 à TA0119, plus TA0089) prolongent la numérotation des tactiques ICS au-delà de TA0111, la dernière qui existe --- le modèle reconnaît la bonne intention, mais il ne connaît pas la liste ICS et invente un numéro vraisemblable plutôt que de s'abstenir.

\textbf{Erreurs de matrice.} Six modèles ont pris au moins un identifiant dans la mauvaise matrice --- un identifiant Enterprise donné sur un journal industriel, ou l'inverse : Gemini 3.1 Pro (1), Kimi K3 (3), DeepSeek V4 Flash (4), DeepSeek V4 Pro (5), Qwen3.8 (7) et GPT-5.6 Luna (9), soit 29 identifiants au total. Les plus fréquents sont TA0007, TA0002, TA0009 et T1059 --- les équivalents Enterprise, de même nom, des tactiques ou techniques ICS attendues --- ainsi que TA0112, tactique Enterprise proposée sur des cas industriels. Les quatre autres modèles --- Claude Opus 4.8, Claude Sonnet 5, GPT-5.6 Sol et Gemini 3.8 Flash --- n'ont produit ni hallucination ni identifiant hors matrice.

\subsubsection{Qualité du mapping et actualité des connaissances}\label{sec:5-2}

Les tableaux~\ref{tab:classement-strict} et~\ref{tab:classement-tolerant} rapportent les mêmes exécutions dans les deux bases de calcul : la base stricte compte les identifiants obsolètes comme des erreurs, la base tolérante les crédite parce que leur remplaçant officiel est la bonne réponse. Comme annoncé, les lignes sont ordonnées par F1 technique, et aucun critère n'est éliminatoire : la mauvaise matrice et l'hallucination, placées en tête de colonnes, sont des réserves opérationnelles qui se lisent avec le F1.

\begin{longtable}[]{@{}
  >{\raggedright\arraybackslash}p{(\columnwidth - 14\tabcolsep) * \real{0.2500}}
  >{\raggedright\arraybackslash}p{(\columnwidth - 14\tabcolsep) * \real{0.1180}}
  >{\raggedright\arraybackslash}p{(\columnwidth - 14\tabcolsep) * \real{0.1050}}
  >{\raggedright\arraybackslash}p{(\columnwidth - 14\tabcolsep) * \real{0.1070}}
  >{\raggedright\arraybackslash}p{(\columnwidth - 14\tabcolsep) * \real{0.1070}}
  >{\raggedright\arraybackslash}p{(\columnwidth - 14\tabcolsep) * \real{0.1070}}
  >{\raggedright\arraybackslash}p{(\columnwidth - 14\tabcolsep) * \real{0.1030}}
  >{\raggedright\arraybackslash}p{(\columnwidth - 14\tabcolsep) * \real{0.1030}}@{}}
\toprule\noalign{}
\begin{minipage}[b]{\linewidth}\raggedright
\textbf{Modèle}
\end{minipage} & \begin{minipage}[b]{\linewidth}\raggedright
\textbf{Mauv. mat.}
\end{minipage} & \begin{minipage}[b]{\linewidth}\raggedright
\textbf{Halluc.}
\end{minipage} & \begin{minipage}[b]{\linewidth}\raggedright
\textbf{F1 tech.}
\end{minipage} & \begin{minipage}[b]{\linewidth}\raggedright
\textbf{F1 tact.}
\end{minipage} & \begin{minipage}[b]{\linewidth}\raggedright
\textbf{F1 s.-t.}
\end{minipage} & \begin{minipage}[b]{\linewidth}\raggedright
\textbf{Obsol.}
\end{minipage} & \begin{minipage}[b]{\linewidth}\raggedright
\textbf{Abst.}
\end{minipage} \\
\midrule\noalign{}
\endhead
\bottomrule\noalign{}
\endlastfoot
GPT-5.6 Sol & 0 & 0,0\,\% & 0,603 & 0,651 & 0,636 & 1,6\,\% & 0,0\,\% \\
DeepSeek V4 Flash & 4 & 0,9\,\% & 0,603 & 0,595 & 0,621 & 1,6\,\% & 0,0\,\% \\
Gemini 3.8 Flash & 0 & 0,0\,\% & 0,600 & 0,648 & 0,667 & 1,9\,\% & 0,0\,\% \\
DeepSeek V4 Pro & 5 & 0,6\,\% & 0,571 & 0,643 & 0,667 & 1,6\,\% & 0,0\,\% \\
Qwen3.8 & 7 & 0,6\,\% & 0,571 & 0,492 & 0,636 & 0,9\,\% & 0,0\,\% \\
Claude Opus 4.8 & 0 & 0,0\,\% & 0,563 & 0,587 & 0,636 & 1,3\,\% & 0,0\,\% \\
Kimi K3 & 3 & 0,3\,\% & 0,563 & 0,540 & 0,621 & 1,6\,\% & 0,0\,\% \\
Claude Sonnet 5 & 0 & 0,0\,\% & 0,536 & 0,512 & 0,611 & 1,9\,\% & 1,6\,\% \\
GPT-5.6 Luna & 9 & 4,1\,\% & 0,526 & 0,502 & 0,636 & 1,6\,\% & 0,6\,\% \\
Gemini 3.1 Pro & 1 & 0,3\,\% & 0,480 & 0,560 & 0,646 & 1,9\,\% & 0,0\,\% \\
\end{longtable}
\legendetableau{Qualité du mapping, base stricte. Lignes ordonnées par F1 technique ; à F1 technique égal, le F1 tactique départage l'affichage.}{tab:classement-strict}

\begin{longtable}[]{@{}
  >{\raggedright\arraybackslash}p{(\columnwidth - 14\tabcolsep) * \real{0.2500}}
  >{\raggedright\arraybackslash}p{(\columnwidth - 14\tabcolsep) * \real{0.1180}}
  >{\raggedright\arraybackslash}p{(\columnwidth - 14\tabcolsep) * \real{0.1050}}
  >{\raggedright\arraybackslash}p{(\columnwidth - 14\tabcolsep) * \real{0.1070}}
  >{\raggedright\arraybackslash}p{(\columnwidth - 14\tabcolsep) * \real{0.1070}}
  >{\raggedright\arraybackslash}p{(\columnwidth - 14\tabcolsep) * \real{0.1070}}
  >{\raggedright\arraybackslash}p{(\columnwidth - 14\tabcolsep) * \real{0.1030}}
  >{\raggedright\arraybackslash}p{(\columnwidth - 14\tabcolsep) * \real{0.1030}}@{}}
\toprule\noalign{}
\begin{minipage}[b]{\linewidth}\raggedright
\textbf{Modèle}
\end{minipage} & \begin{minipage}[b]{\linewidth}\raggedright
\textbf{Mauv. mat.}
\end{minipage} & \begin{minipage}[b]{\linewidth}\raggedright
\textbf{Halluc.}
\end{minipage} & \begin{minipage}[b]{\linewidth}\raggedright
\textbf{F1 tech.}
\end{minipage} & \begin{minipage}[b]{\linewidth}\raggedright
\textbf{F1 tact.}
\end{minipage} & \begin{minipage}[b]{\linewidth}\raggedright
\textbf{F1 s.-t.}
\end{minipage} & \begin{minipage}[b]{\linewidth}\raggedright
\textbf{Obsol.}
\end{minipage} & \begin{minipage}[b]{\linewidth}\raggedright
\textbf{Abst.}
\end{minipage} \\
\midrule\noalign{}
\endhead
\bottomrule\noalign{}
\endlastfoot
Gemini 3.8 Flash & 0 & 0,0\,\% & 0,624 & 0,648 & 0,712 & 1,9\,\% & 0,0\,\% \\
GPT-5.6 Sol & 0 & 0,0\,\% & 0,619 & 0,651 & 0,682 & 1,6\,\% & 0,0\,\% \\
DeepSeek V4 Flash & 4 & 0,9\,\% & 0,619 & 0,595 & 0,667 & 1,6\,\% & 0,0\,\% \\
DeepSeek V4 Pro & 5 & 0,6\,\% & 0,587 & 0,643 & 0,712 & 1,6\,\% & 0,0\,\% \\
Kimi K3 & 3 & 0,3\,\% & 0,579 & 0,540 & 0,667 & 1,6\,\% & 0,0\,\% \\
Claude Opus 4.8 & 0 & 0,0\,\% & 0,571 & 0,587 & 0,682 & 1,3\,\% & 0,0\,\% \\
Qwen3.8 & 7 & 0,6\,\% & 0,571 & 0,492 & 0,682 & 0,9\,\% & 0,0\,\% \\
Claude Sonnet 5 & 0 & 0,0\,\% & 0,560 & 0,512 & 0,656 & 1,9\,\% & 1,6\,\% \\
GPT-5.6 Luna & 9 & 4,1\,\% & 0,542 & 0,502 & 0,682 & 1,6\,\% & 0,6\,\% \\
Gemini 3.1 Pro & 1 & 0,3\,\% & 0,504 & 0,560 & 0,692 & 1,9\,\% & 0,0\,\% \\
\end{longtable}
\legendetableau{Qualité du mapping, base tolérante. Mêmes exécutions et même ordre de colonnes que le tableau~\ref{tab:classement-strict} ; seule différence, les identifiants obsolètes sont crédités.}{tab:classement-tolerant}

\textbf{1 --- La qualité du mapping.} Le F1 combine la précision --- la part de réponses justes parmi les réponses données --- et le rappel --- la part d'échantillons résolus. Il vaut entre 0 et 1 ; plus il est haut, mieux c'est. Au niveau technique, celui qui compte le plus puisque les mitigations MITRE s'y rattachent, trois modèles se détachent en base stricte : GPT-5.6 Sol et DeepSeek V4 Flash à 0,603, Gemini 3.8 Flash à 0,600. L'écart entre eux, 0,003, est plus petit que $\pm$0,008, l'écart observé en rejouant deux fois la même campagne : ces trois modèles sont donc ex æquo de fait. En base tolérante, Gemini 3.8 Flash passe en tête (0,624) devant Sol et DeepSeek V4 Flash (0,619) --- l'écart, 0,005, reste sous le grain : ex æquo de fait là aussi. Au niveau tactique, Sol (0,651), Gemini 3.8 Flash (0,648) et DeepSeek V4 Pro (0,643) se tiennent dans le même mouchoir. Au niveau sous-technique, les dix modèles se tiennent entre 0,611 et 0,667 en base stricte : c'est le niveau où ils se distinguent le moins ; en base tolérante, Gemini 3.8 Flash et DeepSeek V4 Pro atteignent tous deux 0,712. Enfin, deux modèles ont eu des réponses rejetées, comptées comme abstentions : deux JSON invalides pour Claude Sonnet 5 --- des guillemets non échappés dans l'extrait de journal cité ---, et un format non respecté pour GPT-5.6 Luna, qui a rendu « TA???? » à la place d'un identifiant de tactique. Recalculés hors rejets, leurs F1 techniques montent légèrement --- Sonnet 0,540 en base stricte et 0,565 en tolérante, Luna 0,528 et 0,544 --- sans changer l'ordre général.

\textbf{2 --- Le taux d'hallucination et les fausses matrices.} Quatre modèles n'inventent jamais d'identifiant : GPT-5.6 Sol, Claude Opus 4.8, Claude Sonnet 5 et Gemini 3.8 Flash, tous à 0,0\,\%. Viennent ensuite Kimi K3 et Gemini 3.1 Pro (0,3\,\%), DeepSeek V4 Pro et Qwen3.8 (0,6\,\%), DeepSeek V4 Flash (0,9\,\%), puis GPT-5.6 Luna (4,1\,\%), qui concentre à lui seul treize des vingt-deux identifiants fabriqués du panel. Côté fausses matrices, les mêmes quatre modèles sont à zéro, et les six autres vont de un identifiant (Gemini 3.1 Pro) à neuf (GPT-5.6 Luna). Que la liste des modèles propres soit identique sur les deux fautes n'est pas un hasard : les deux mesurent la même compétence --- situer l'univers, Enterprise ou ICS, avant de nommer --- l'une échouant vers un identifiant qui existe ailleurs, l'autre vers un identifiant qui n'existe nulle part.

\textbf{3 --- La stabilité des réponses.} La stabilité mesure la part des 42 échantillons pour lesquels le modèle a rendu exactement le même triplet --- tactique, technique, sous-technique --- à ses trois exécutions. Gemini 3.1 Pro est le plus constant (88,1\,\%), suivi de Gemini 3.8 Flash (81,0\,\%), Claude Opus 4.8 (76,2\,\%), Claude Sonnet 5 (66,7\,\%), GPT-5.6 Sol (64,3\,\%), Kimi K3 (59,5\,\%), Qwen3.8 (57,1\,\%), GPT-5.6 Luna (54,8\,\%), DeepSeek V4 Flash (52,4\,\%) et DeepSeek V4 Pro (42,9\,\%). La stabilité ne mesure pas la justesse : Gemini 3.1 Pro, le plus stable, a aussi le F1 technique le plus bas --- il se trompe de façon très reproductible --- et GPT-5.6 Sol, en tête de la qualité en base stricte, n'est stable qu'à 64,3\,\% : ses bons scores sont une moyenne de réponses qui se contredisent d'une exécution à l'autre, or en exploitation on ne relance pas trois fois.

\subsubsection{Temps de réponse et coût}\label{sec:5-3}

Le tableau~\ref{tab:temps-cout} rapporte, pour chaque modèle, quatre mesures simples. Le \textbf{p50} est le temps de réponse médian d'un appel réussi, en millisecondes : la moitié des appels sont plus rapides. Le \textbf{p95} est le temps sous lequel 95\,\% des appels aboutissent : c'est le pire cas ordinaire, celui qu'une moyenne cacherait. Les \textbf{jetons} sont l'unité de texte que les fournisseurs facturent ; la colonne donne la moyenne consommée par exécution, réflexion interne du modèle comprise. L'\textbf{efficacité} rapporte le F1 technique strict au coût : c'est le F1 obtenu par millier de jetons consommés. L'\textbf{argent}, enfin, est ce que la campagne complète --- 126 exécutions --- a réellement coûté, chaque composante étant facturée à son tarif du 16 septembre 2026. Les lignes sont ordonnées du modèle le plus rapide au plus lent en médiane.

\begin{longtable}[]{@{}
  >{\raggedright\arraybackslash}p{(\columnwidth - 10\tabcolsep) * \real{0.2700}}
  >{\raggedright\arraybackslash}p{(\columnwidth - 10\tabcolsep) * \real{0.1460}}
  >{\raggedright\arraybackslash}p{(\columnwidth - 10\tabcolsep) * \real{0.1460}}
  >{\raggedright\arraybackslash}p{(\columnwidth - 10\tabcolsep) * \real{0.1460}}
  >{\raggedright\arraybackslash}p{(\columnwidth - 10\tabcolsep) * \real{0.1460}}
  >{\raggedright\arraybackslash}p{(\columnwidth - 10\tabcolsep) * \real{0.1460}}@{}}
\toprule\noalign{}
\begin{minipage}[b]{\linewidth}\raggedright
\textbf{Modèle}
\end{minipage} & \begin{minipage}[b]{\linewidth}\raggedright
\textbf{p50 (ms)}
\end{minipage} & \begin{minipage}[b]{\linewidth}\raggedright
\textbf{p95 (ms)}
\end{minipage} & \begin{minipage}[b]{\linewidth}\raggedright
\textbf{Jetons}
\end{minipage} & \begin{minipage}[b]{\linewidth}\raggedright
\textbf{Effic.}
\end{minipage} & \begin{minipage}[b]{\linewidth}\raggedright
\textbf{Argent}
\end{minipage} \\
\midrule\noalign{}
\endhead
\bottomrule\noalign{}
\endlastfoot
DeepSeek V4 Flash & 3 407 & 19 812 & 5 594 & 0,108 & 0,12\,\$ \\
Claude Opus 4.8 & 4 282 & 6 050 & 6 572 & 0,086 & 4,63\,\$ \\
Claude Sonnet 5 & 4 291 & 10 994 & 6 740 & 0,080 & 2,06\,\$ \\
GPT-5.6 Luna & 4 626 & 28 758 & 5 004 & 0,105 & 0,22\,\$ \\
GPT-5.6 Sol & 6 086 & 24 145 & 4 649 & 0,130 & 3,09\,\$ \\
Gemini 3.8 Flash & 10 664 & 44 482 & 7 566 & 0,079 & 1,36\,\$ \\
Gemini 3.1 Pro & 18 532 & 103 824 & 9 023 & 0,053 & 6,23\,\$ \\
Kimi K3 & 24 034 & 202 504 & 6 027 & 0,094 & 4,76\,\$ \\
DeepSeek V4 Pro & 40 724 & 151 756 & 7 941 & 0,072 & 0,62\,\$ \\
Qwen3.8 & 82 854 & 269 617 & 9 641 & 0,059 & 4,49\,\$ \\
\end{longtable}
\legendetableau{Temps de réponse et coût, par modèle. p50 et p95 : latence médiane et 95\up{e} centile de l'appel réussi. Jetons : moyenne par exécution, raisonnement compris. Effic. : F1 technique strict par millier de jetons. Argent : coût réel des 126 exécutions aux tarifs du 16 septembre 2026 (tarifs.json) ; le tarif de GPT-5.6 Sol est promotionnel, et celui de Gemini 3.8 Flash double au 1\up{er} janvier 2027.}{tab:temps-cout}

\textbf{Le temps.} Trois modèles répondent en moins de cinq secondes en médiane : DeepSeek V4 Flash (3,4~s), Claude Opus 4.8 et Claude Sonnet 5 (4,3~s). Mais la médiane ne dit pas tout : le p95 sépare les modèles prévisibles des modèles erratiques. Claude Opus 4.8 est le seul dont le pire cas ordinaire reste proche du cas courant --- 6,1~s, soit 1,4 fois sa médiane, le plus bas des dix --- quand DeepSeek V4 Flash monte à 19,8~s, Gemini 3.8 Flash à 44,5~s et Qwen3.8 à 270~s. Dans un centre opérationnel qui traite des milliers d'alertes, ce p95 est le temps que l'analyste attend vraiment les mauvais jours.

\textbf{Le coût.} Les jetons ne sont pas la dépense, ils n'en sont que l'unité : chaque fournisseur fixe son tarif, et seul le montant facturé se paie. Le classement en argent ne suit donc ni les jetons ni le temps : GPT-5.6 Sol, le plus sobre en jetons, est le troisième plus cher des dix (3,09\,\$), et Gemini 3.8 Flash, 63\,\% plus gourmand que lui, coûte deux fois moins (1,36\,\$). Les moins chers sont DeepSeek V4 Flash (0,12\,\$), GPT-5.6 Luna (0,22\,\$) et DeepSeek V4 Pro (0,62\,\$) ; les plus chers, Gemini 3.1 Pro (6,23\,\$), Kimi K3 (4,76\,\$) et Claude Opus 4.8 (4,63\,\$). L'écart de consommation entre modèles vient surtout des jetons de raisonnement --- la réflexion interne, facturée comme de la sortie sans être restituée --- et c'est aussi lui qui creuse les latences. En efficacité, GPT-5.6 Sol domine nettement (0,130), devant DeepSeek V4 Flash (0,108) et GPT-5.6 Luna (0,105).

\textbf{Classement des modèles.} Aucun modèle ne gagne partout, alors résumons axe par axe, simplement. La meilleure qualité de mapping : GPT-5.6 Sol, DeepSeek V4 Flash et Gemini 3.8 Flash, ex æquo de fait en base stricte, Gemini 3.8 Flash en tête en base tolérante. Les seuls sans aucune hallucination ni erreur de matrice : Claude Opus 4.8, Claude Sonnet 5, GPT-5.6 Sol et Gemini 3.8 Flash. Les plus stables : Gemini 3.1 Pro, Gemini 3.8 Flash et Claude Opus 4.8. Les plus rapides et les plus prévisibles : DeepSeek V4 Flash en médiane, Claude Opus 4.8 au p95. Les moins chers : DeepSeek V4 Flash, GPT-5.6 Luna et DeepSeek V4 Pro. Le plus efficace : GPT-5.6 Sol. Deux modèles reviennent le plus souvent dans ces podiums : Gemini 3.8 Flash --- qualité, propreté, stabilité, prix contenu, mais un p95 de 44,5~s --- et Claude Opus 4.8 --- propreté, stabilité, temps prévisible, mais un F1 technique à 0,040 de la tête et le troisième prix le plus élevé. Pour la phase 2, qui rejoue la même campagne avec le serveur MCP, il faut un instrument : un modèle qui ne fabrique rien, qui rend des réponses reproductibles, dont le temps ne dérive pas quand le contexte s'alourdit d'appels d'outils, et qui laisse une marge visible à combler --- son F1 technique tombe de 0,714 sur la matrice Enterprise à 0,262 sur la matrice ICS. \textbf{Claude Opus 4.8 est retenu} ; les tableaux de cette section constituent sa ligne de base sans serveur, celle contre laquelle l'apport de l'ancrage sera mesuré.

%% TODO : conclusion du chapitre (retour sur QR1--QR4 et transition vers
%% le chapitre 2), à rédiger.

% =====================================================================
% ANCIENNE RÉDACTION (campagne recalculée le 17/09 avec le script
% corrigé ; mêmes effectifs et F1). Conservée ci-dessous, désactivée par
% iffalse/fi, comme référence pour les sections restantes — fiabilité,
% temps de traitement, coût, choix du modèle — jusqu'à leur réécriture.
% =====================================================================
\iffalse
L'évaluation part de l'étiquetage brut, avant toute agrégation. Le tableau~\ref{tab:etiquetage} rapporte, modèle par modèle et niveau par niveau, les effectifs des catégories définies en section~\ref{sec:2-2} --- VP, FP, FN dont abstentions (Ab), identifiants obsolètes (OBS) et hallucinations (H) --- auxquels s'ajoute le compteur de mauvaise matrice (MM) de la section~\ref{sec:2-5}, c'est-à-dire les identifiants pris hors de la matrice de l'échantillon.

\begin{longtable}[]{@{}
  >{\raggedright\arraybackslash}p{(\columnwidth - 16\tabcolsep) * \real{0.2032}}
  >{\raggedright\arraybackslash}p{(\columnwidth - 16\tabcolsep) * \real{0.2506}}
  >{\raggedright\arraybackslash}p{(\columnwidth - 16\tabcolsep) * \real{0.0756}}
  >{\raggedright\arraybackslash}p{(\columnwidth - 16\tabcolsep) * \real{0.0756}}
  >{\raggedright\arraybackslash}p{(\columnwidth - 16\tabcolsep) * \real{0.0756}}
  >{\raggedright\arraybackslash}p{(\columnwidth - 16\tabcolsep) * \real{0.0948}}
  >{\raggedright\arraybackslash}p{(\columnwidth - 16\tabcolsep) * \real{0.0767}}
  >{\raggedright\arraybackslash}p{(\columnwidth - 16\tabcolsep) * \real{0.0655}}
  >{\raggedright\arraybackslash}p{(\columnwidth - 16\tabcolsep) * \real{0.0824}}@{}}
\toprule\noalign{}
\begin{minipage}[b]{\linewidth}\raggedright
\textbf{Modèle}
\end{minipage} & \begin{minipage}[b]{\linewidth}\raggedright
\textbf{Niveau}
\end{minipage} & \begin{minipage}[b]{\linewidth}\raggedright
\textbf{VP}
\end{minipage} & \begin{minipage}[b]{\linewidth}\raggedright
\textbf{FP}
\end{minipage} & \begin{minipage}[b]{\linewidth}\raggedright
\textbf{FN}
\end{minipage} & \begin{minipage}[b]{\linewidth}\raggedright
\textbf{dont Ab}
\end{minipage} & \begin{minipage}[b]{\linewidth}\raggedright
\textbf{OBS}
\end{minipage} & \begin{minipage}[b]{\linewidth}\raggedright
\textbf{H}
\end{minipage} & \begin{minipage}[b]{\linewidth}\raggedright
\textbf{MM}
\end{minipage} \\
\midrule\noalign{}
\endhead
\bottomrule\noalign{}
\endlastfoot
\textbf{Claude Opus 4.8} & Tactique (n = 126) & 74 & 52 & 52 & 0 & 0 & 0 & 0 \\
& Technique (n = 126) & 71 & 55 & 55 & 0 & 1 & 0 & 0 \\
& Sous-technique (n = 66) & 42 & 24 & 24 & 0 & 3 & 0 & 0 \\
\textbf{Claude Sonnet 5} & Tactique (n = 126) & 64 & 60 & 62 & 2 & 0 & 0 & 0 \\
& Technique (n = 126) & 67 & 57 & 59 & 2 & 3 & 0 & 0 \\
& Sous-technique (n = 66) & 40 & 25 & 26 & 1 & 3 & 0 & 0 \\
\textbf{GPT-5.6 Sol} & Tactique (n = 126) & 82 & 44 & 44 & 0 & 0 & 0 & 0 \\
& Technique (n = 126) & 76 & 50 & 50 & 0 & 2 & 0 & 0 \\
& Sous-technique (n = 66) & 42 & 24 & 24 & 0 & 3 & 0 & 0 \\
\textbf{GPT-5.6 Luna} & Tactique (n = 126) & 63 & 62 & 63 & 1 & 0 & 13 & 4 \\
& Technique (n = 126) & 66 & 59 & 60 & 1 & 2 & 0 & 5 \\
& Sous-technique (n = 66) & 42 & 24 & 24 & 0 & 3 & 0 & 0 \\
\textbf{Gemini 3.1 Pro} & Tactique (n = 125) & 70 & 55 & 55 & 0 & 0 & 1 & 0 \\
& Technique (n = 125) & 60 & 65 & 65 & 0 & 3 & 0 & 1 \\
& Sous-technique (n = 65) & 42 & 23 & 23 & 0 & 3 & 0 & 0 \\
\textbf{Gemini 3.8 Flash} & Tactique (n = 125) & 81 & 44 & 44 & 0 & 0 & 0 & 0 \\
& Technique (n = 125) & 75 & 50 & 50 & 0 & 3 & 0 & 0 \\
& Sous-technique (n = 66) & 44 & 22 & 22 & 0 & 3 & 0 & 0 \\
\textbf{Kimi K3} & Tactique (n = 126) & 68 & 58 & 58 & 0 & 0 & 1 & 3 \\
& Technique (n = 126) & 71 & 55 & 55 & 0 & 2 & 0 & 0 \\
& Sous-technique (n = 66) & 41 & 25 & 25 & 0 & 3 & 0 & 0 \\
\textbf{DeepSeek V4 Pro} & Tactique (n = 126) & 81 & 45 & 45 & 0 & 0 & 2 & 4 \\
& Technique (n = 126) & 72 & 54 & 54 & 0 & 2 & 0 & 1 \\
& Sous-technique (n = 66) & 44 & 22 & 22 & 0 & 3 & 0 & 0 \\
\textbf{DeepSeek V4 Flash} & Tactique (n = 126) & 75 & 51 & 51 & 0 & 0 & 3 & 3 \\
& Technique (n = 126) & 76 & 50 & 50 & 0 & 2 & 0 & 1 \\
& Sous-technique (n = 66) & 41 & 25 & 25 & 0 & 3 & 0 & 0 \\
\textbf{Qwen3.8} & Tactique (n = 126) & 62 & 64 & 64 & 0 & 0 & 2 & 6 \\
& Technique (n = 126) & 72 & 54 & 54 & 0 & 0 & 0 & 1 \\
& Sous-technique (n = 66) & 42 & 24 & 24 & 0 & 3 & 0 & 0 \\
\end{longtable}
\legendetableau{Étiquetage des prédictions des dix modèles. Entre parenthèses, le nombre d'exécutions évaluées au niveau considéré. FP est la valeur de la base stricte : elle comprend les colonnes OBS et H, reportées à part comme sous-comptes, ainsi que les identifiants périmés --- dépréciés ou d'anciennes versions, tel le TA0024 de Qwen3.8 --- non détaillés ici ; FN = FP + dont Ab. Les effectifs de 125 des deux Gemini correspondent aux deux itérations abandonnées sur quota d'API après cinq tentatives (S015, exécution 2, pour 3.1 Pro ; S013, exécution 3, pour 3.8 Flash), écartées et rapportées conformément à la section~\ref{sec:meth-etiquetage}.}{tab:etiquetage}


Les hallucinations n'apparaissent qu'au niveau tactique, et sur six modèles ; GPT-5.6 Luna en concentre treize à lui seul. Ces six modèles sont exactement ceux qui se trompent de matrice : les deux Claude, GPT-5.6 Sol et Gemini 3.8 Flash n'ont aucun identifiant fabriqué ni hors matrice. L'abstention reste marginale, limitée à Claude Sonnet 5 et GPT-5.6 Luna. Les identifiants obsolètes, indétectables au niveau tactique, se concentrent au niveau sous-technique, où neuf modèles sur dix en produisent exactement trois. À ce même niveau, les dix modèles se tiennent entre 40 et 44 bonnes réponses sur 66 : c'est celui où ils se distinguent le moins.

\textbf{Les hallucinations sont un problème d'homonymie entre matrices.} Les vingt-deux identifiants fabriqués du panel sont tous des tactiques, et tous sur des échantillons ICS. Un vingt-troisième cas, la tactique TA0024 proposée par Qwen3.8, n'en est pas un : l'identifiant provient de PRE-ATT\&CK, matrice absorbée par Enterprise en 2020, et rejoint le compteur des identifiants périmés (section~\ref{sec:2-2}), ramenant son taux d'hallucination à 0,6\,\%. Ils ne relèvent pas d'une invention arbitraire mais d'un défaut de repérage entre les deux matrices, dont l'homonymie est la cause. Neuf tactiques portent en effet un nom identique de part et d'autre : Collection (TA0009 / TA0100), Command and Control (TA0011 / TA0101), Discovery (TA0007 / TA0102), Execution (TA0002 / TA0104), Impact (TA0040 / TA0105), Initial Access (TA0001 / TA0108), Lateral Movement (TA0008 / TA0109), Persistence (TA0003 / TA0110) et Privilege Escalation (TA0004 / TA0111). Le modèle reconnaît l'intention, mais doit encore choisir l'univers --- et c'est là qu'il échoue, de deux façons : soit il donne l'identifiant Enterprise homonyme, soit il prolonge la numérotation ICS au-delà de TA0111, dernière tactique réelle, produisant TA0113 à TA0119. Au niveau technique, la même homonymie joue --- T1046 (Network Service Discovery, Enterprise) répondu pour T0846 (Network Service Scanning, ICS), T1059 pour T0853, l'interpréteur de commandes d'Enterprise contre le Scripting industriel --- et six occurrences répondent TA0112, la tactique la plus récente d'Enterprise, sur des cas industriels : la connaissance est à jour, la discipline de matrice ne suit pas. L'hallucination, elle, n'est pas un délire mais une interpolation : la liste des tactiques ICS est courte et rare dans les corpus d'entraînement ; le modèle reconnaît le concept, sait qu'il lui faut un numéro de la gamme ICS, et en fabrique un plausible --- TA0115 sept fois chez GPT-5.6 Luna --- plutôt que de s'abstenir. Le contraste le confirme : au niveau technique, zéro invention sur tout le panel, les faux identifiants techniques étant tous des connaissances périmées, jamais des fabrications.

Un score unique masquerait les compromis : un modèle rapide peut être bavard, un modèle précis peut s'abstenir souvent, un modèle exact peut s'arrêter trop tôt dans la hiérarchie.

Les modèles LLM seront comparés entre eux sur ces différents axes :

\begin{itemize}
\item
  \textbf{Qualité du mapping et actualité des connaissances} --- rapportée en deux bases, stricte et tolérante, et ordonnée par gravité décroissante de la faute : le mauvais choix de matrice, le taux d'hallucination, le F1 par niveau --- technique, tactique, sous-technique ---, le taux d'obsolescence, puis le taux d'abstention.
\item
  \textbf{Fiabilité, temps de traitement et coût} --- ordonnée dans cet ordre : la stabilité, la latence médiane, le coût en jetons, puis l'efficacité.
\end{itemize}

\textbf{Comment le premier axe est classé.} L'ordre des critères suit la gravité de la faute, et non sa fréquence. Le mauvais choix de matrice vient en tête : un identifiant Enterprise proposé sur un journal ICS --- ou l'inverse --- ne désigne pas seulement une mauvaise technique, il désigne un mauvais univers, et la mitigation qui en découle est inapplicable au système observé. Il est compté en effectif brut, les trois niveaux confondus, parce qu'un seul de ces identifiants suffit à rendre la réponse inutilisable. Vient ensuite le taux d'hallucination : un identifiant fabriqué oriente l'analyste vers une entrée qui n'existe pas, et un modèle qui en produit est écarté d'un usage opérationnel quel que soit son F1. Le F1 ne départage donc que les modèles ayant passé ces deux filtres ; il est lu niveau par niveau, la technique d'abord --- les mitigations MITRE s'y rattachent ---, puis la tactique, puis la sous-technique. Le taux d'obsolescence et le taux d'abstention ne servent qu'à trancher les égalités restantes.

Le tri est lexicographique : chaque critère ne départage que les égalités laissées par le précédent, et aucune grandeur composite n'est fabriquée --- chaque colonne du tableau reste une métrique définie en section~\ref{sec:meth-metriques}. Le même ordre est appliqué aux deux bases, de sorte que la seule différence entre les tableaux~\ref{tab:classement-strict} et~\ref{tab:classement-tolerant} soit le traitement des identifiants obsolètes : comptés comme des erreurs en base stricte, crédités en base tolérante. L'écart entre les deux mesure l'actualité des connaissances du modèle, et rien d'autre.

\textbf{Comment le second axe est classé.} La stabilité vient en premier : un modèle qui ne rend pas deux fois la même réponse au même journal n'est pas exploitable dans un centre opérationnel, quelle que soit sa vitesse. La latence médiane (p50) suit, comme temps de traitement courant ; le 95e centile est rapporté à titre de pire cas ordinaire mais n'entre pas dans le tri, un seul appel lent ne devant pas déclasser un modèle par ailleurs régulier. Le coût vient ensuite, mesuré par les jetons moyens par exécution. L'efficacité --- F1 technique strict par millier de jetons --- est placée en dernier parce qu'elle est le seul indicateur de ce tableau à dépendre de la base de calcul : elle passe par l'étiquetage, alors que la stabilité, la latence et les jetons n'en dépendent pas. La sur-spécification ne sert que de départage ultime.

Dans les deux axes, les critères sont donc appliqués l'un après l'autre, dans l'ordre annoncé : un modèle ne peut pas compenser une faute grave par une bonne performance sur un critère plus fin.

\subsubsection{Qualité du mapping et actualité des connaissances}\label{sec:5-1-ancien}
Les tableaux~\ref{tab:classement-strict} et~\ref{tab:classement-tolerant} appliquent l'ordre de critères annoncé. Ils portent sur les mêmes exécutions et ne diffèrent que par le traitement des identifiants obsolètes ; les deux premiers critères --- mauvaise matrice et hallucination --- ne passent pas par la base de calcul et sont donc identiques dans les deux.

\begin{longtable}[]{@{}
  >{\raggedright\arraybackslash}p{(\columnwidth - 16\tabcolsep) * \real{0.0677}}
  >{\raggedright\arraybackslash}p{(\columnwidth - 16\tabcolsep) * \real{0.2257}}
  >{\raggedright\arraybackslash}p{(\columnwidth - 16\tabcolsep) * \real{0.1140}}
  >{\raggedright\arraybackslash}p{(\columnwidth - 16\tabcolsep) * \real{0.0971}}
  >{\raggedright\arraybackslash}p{(\columnwidth - 16\tabcolsep) * \real{0.1016}}
  >{\raggedright\arraybackslash}p{(\columnwidth - 16\tabcolsep) * \real{0.1016}}
  >{\raggedright\arraybackslash}p{(\columnwidth - 16\tabcolsep) * \real{0.1016}}
  >{\raggedright\arraybackslash}p{(\columnwidth - 16\tabcolsep) * \real{0.0959}}
  >{\raggedright\arraybackslash}p{(\columnwidth - 16\tabcolsep) * \real{0.0948}}@{}}
\toprule\noalign{}
\begin{minipage}[b]{\linewidth}\raggedright
\textbf{Rang}
\end{minipage} & \begin{minipage}[b]{\linewidth}\raggedright
\textbf{Modèle}
\end{minipage} & \begin{minipage}[b]{\linewidth}\raggedright
\textbf{Mauv. mat.}
\end{minipage} & \begin{minipage}[b]{\linewidth}\raggedright
\textbf{Halluc.}
\end{minipage} & \begin{minipage}[b]{\linewidth}\raggedright
\textbf{F1 tech.}
\end{minipage} & \begin{minipage}[b]{\linewidth}\raggedright
\textbf{F1 tact.}
\end{minipage} & \begin{minipage}[b]{\linewidth}\raggedright
\textbf{F1 s.-t.}
\end{minipage} & \begin{minipage}[b]{\linewidth}\raggedright
\textbf{Obsol.}
\end{minipage} & \begin{minipage}[b]{\linewidth}\raggedright
\textbf{Abst.}
\end{minipage} \\
\midrule\noalign{}
\endhead
\bottomrule\noalign{}
\endlastfoot
GPT-5.6 Sol & 0 & 0,0\,\% & 0,603 & 0,651 & 0,636 & 1,6\,\% & 0,0\,\% \\
Gemini 3.8 Flash & 0 & 0,0\,\% & 0,600 & 0,648 & 0,667 & 1,9\,\% & 0,0\,\% \\
Claude Opus 4.8 & 0 & 0,0\,\% & 0,563 & 0,587 & 0,636 & 1,3\,\% & 0,0\,\% \\
Claude Sonnet 5 & 0 & 0,0\,\% & 0,536 & 0,512 & 0,611 & 1,9\,\% & 1,6\,\% \\
Gemini 3.1 Pro & 1 & 0,3\,\% & 0,480 & 0,560 & 0,646 & 1,9\,\% & 0,0\,\% \\
Kimi K3 & 3 & 0,3\,\% & 0,563 & 0,540 & 0,621 & 1,6\,\% & 0,0\,\% \\
DeepSeek V4 Flash & 4 & 0,9\,\% & 0,603 & 0,595 & 0,621 & 1,6\,\% & 0,0\,\% \\
DeepSeek V4 Pro & 5 & 0,6\,\% & 0,571 & 0,643 & 0,667 & 1,6\,\% & 0,0\,\% \\
Qwen3.8 & 7 & 0,6\,\% & 0,571 & 0,492 & 0,636 & 0,9\,\% & 0,0\,\% \\
GPT-5.6 Luna & 9 & 4,1\,\% & 0,526 & 0,502 & 0,636 & 1,6\,\% & 0,6\,\% \\
\end{longtable}
\legendetableau{Qualité du mapping, base stricte.}{tab:classement-strict}




\textbf{Un exemple où l'erreur est entièrement imputable à l'homonymie.} L'échantillon S012 --- exécution d'un script sur l'hôte PLC, tactique attendue TA0104 Execution, matrice ICS --- l'illustre sans ambiguïté. Sept modèles répondent TA0104 aux trois exécutions. GPT-5.6 Luna répond TA0002, l'Execution d'Enterprise, aux trois exécutions ; Qwen3.8 fait de même deux fois sur trois. Aucun des deux ne s'est trompé d'intention : ils ont correctement identifié une exécution de code, puis désigné la mauvaise matrice. Trois autres échantillons présentent le même profil --- S000 et S001 pour Discovery, S003 pour Collection --- et les trois modèles concernés sont exactement ceux que le premier critère élimine.

\textbf{Conséquence opérationnelle.} La mitigation attachée à TA0002 vise des systèmes d'information ; celle de TA0104 vise des automates. Un identifiant juste de nom mais faux d'univers oriente l'analyste vers un catalogue de contre-mesures inapplicable au système observé. C'est pourquoi le mauvais choix de matrice précède l'hallucination dans l'ordre des critères, et l'hallucination le F1.

\textbf{Première étape : l'élimination par la matrice.} Six modèles sur dix produisent au moins un identifiant pris hors de la matrice de l'échantillon et sont écartés avant même que leur F1 soit lu : Gemini 3.1 Pro (1), Kimi K3 (3), DeepSeek V4 Flash (4), DeepSeek V4 Pro (5), Qwen3.8 (7) et GPT-5.6 Luna (9), soit vingt-neuf identifiants au total. Le cas de DeepSeek V4 Flash est le plus démonstratif : son F1 technique, 0,603, égale celui du premier du tableau, et il termine pourtant septième. C'est exactement l'effet recherché par l'ordre des critères --- une réponse prise dans le mauvais univers ne se rattrape pas par un bon score ailleurs, puisque la mitigation qu'elle entraîne ne s'applique pas au système observé.

\textbf{Deuxième étape : l'hallucination.} Elle ne départage personne de plus ici, les six modèles concernés ayant déjà été écartés par la matrice, mais elle confirme la dernière place de GPT-5.6 Luna : ses 4,1\,\% --- treize identifiants fabriqués sur les vingt-deux du panel --- représentent plus de quatre fois le pire des neuf autres modèles, tous sous 1\,\%.

\textbf{Troisième étape : le F1.} Il n'arbitre donc qu'entre les quatre modèles qui franchissent les deux premiers filtres sans rien concéder : GPT-5.6 Sol, Gemini 3.8 Flash, Claude Opus 4.8 et Claude Sonnet 5. En base stricte, GPT-5.6 Sol prend la tête d'un souffle sur le F1 technique (0,603 contre 0,600), les deux Claude suivant à 0,040 et 0,067.

\begin{longtable}[]{@{}
  >{\raggedright\arraybackslash}p{(\columnwidth - 16\tabcolsep) * \real{0.0677}}
  >{\raggedright\arraybackslash}p{(\columnwidth - 16\tabcolsep) * \real{0.2257}}
  >{\raggedright\arraybackslash}p{(\columnwidth - 16\tabcolsep) * \real{0.1140}}
  >{\raggedright\arraybackslash}p{(\columnwidth - 16\tabcolsep) * \real{0.0971}}
  >{\raggedright\arraybackslash}p{(\columnwidth - 16\tabcolsep) * \real{0.1016}}
  >{\raggedright\arraybackslash}p{(\columnwidth - 16\tabcolsep) * \real{0.1016}}
  >{\raggedright\arraybackslash}p{(\columnwidth - 16\tabcolsep) * \real{0.1016}}
  >{\raggedright\arraybackslash}p{(\columnwidth - 16\tabcolsep) * \real{0.0959}}
  >{\raggedright\arraybackslash}p{(\columnwidth - 16\tabcolsep) * \real{0.0948}}@{}}
\toprule\noalign{}
\begin{minipage}[b]{\linewidth}\raggedright
\textbf{Rang}
\end{minipage} & \begin{minipage}[b]{\linewidth}\raggedright
\textbf{Modèle}
\end{minipage} & \begin{minipage}[b]{\linewidth}\raggedright
\textbf{Mauv. mat.}
\end{minipage} & \begin{minipage}[b]{\linewidth}\raggedright
\textbf{Halluc.}
\end{minipage} & \begin{minipage}[b]{\linewidth}\raggedright
\textbf{F1 tech.}
\end{minipage} & \begin{minipage}[b]{\linewidth}\raggedright
\textbf{F1 tact.}
\end{minipage} & \begin{minipage}[b]{\linewidth}\raggedright
\textbf{F1 s.-t.}
\end{minipage} & \begin{minipage}[b]{\linewidth}\raggedright
\textbf{Obsol.}
\end{minipage} & \begin{minipage}[b]{\linewidth}\raggedright
\textbf{Abst.}
\end{minipage} \\
\midrule\noalign{}
\endhead
\bottomrule\noalign{}
\endlastfoot
Gemini 3.8 Flash & 0 & 0,0\,\% & 0,624 & 0,648 & 0,712 & 1,9\,\% & 0,0\,\% \\
GPT-5.6 Sol & 0 & 0,0\,\% & 0,619 & 0,651 & 0,682 & 1,6\,\% & 0,0\,\% \\
Claude Opus 4.8 & 0 & 0,0\,\% & 0,571 & 0,587 & 0,682 & 1,3\,\% & 0,0\,\% \\
Claude Sonnet 5 & 0 & 0,0\,\% & 0,560 & 0,512 & 0,656 & 1,9\,\% & 1,6\,\% \\
Gemini 3.1 Pro & 1 & 0,3\,\% & 0,504 & 0,560 & 0,692 & 1,9\,\% & 0,0\,\% \\
Kimi K3 & 3 & 0,3\,\% & 0,579 & 0,540 & 0,667 & 1,6\,\% & 0,0\,\% \\
DeepSeek V4 Flash & 4 & 0,9\,\% & 0,619 & 0,595 & 0,667 & 1,6\,\% & 0,0\,\% \\
DeepSeek V4 Pro & 5 & 0,6\,\% & 0,587 & 0,643 & 0,712 & 1,6\,\% & 0,0\,\% \\
Qwen3.8 & 7 & 0,6\,\% & 0,571 & 0,492 & 0,682 & 0,9\,\% & 0,0\,\% \\
GPT-5.6 Luna & 9 & 4,1\,\% & 0,542 & 0,502 & 0,682 & 1,6\,\% & 0,6\,\% \\
\end{longtable}
\legendetableau{Qualité du mapping, base tolérante.}{tab:classement-tolerant}

Les deux tables se lisent axe par axe, sans critère éliminatoire : les matrices se ressemblant jusqu'à l'homonymie, réussir ou rater l'univers tient parfois à très peu, et une élimination donnerait à ce très-peu un poids disproportionné.

\textbf{La qualité du mapping.} Trois modèles sortent du lot, dans un mouchoir inférieur au grain de reproductibilité ($\pm$0,008) : au niveau technique, GPT-5.6 Sol et DeepSeek V4 Flash (0,603) devant Gemini 3.8 Flash (0,600) en base stricte ; en base tolérante, Gemini 3.8 Flash passe premier devant Sol. Au niveau tactique, Sol mène (0,651) devant Gemini 3.8 Flash (0,648) ; en sous-technique, Gemini 3.8 Flash domine seul (0,712 en tolérante). Podium de l'axe : Sol, Gemini 3.8 Flash, DeepSeek V4 Flash.

\textbf{L'hallucination, en taux.} Quatre modèles n'inventent jamais --- GPT-5.6 Sol, les deux Claude et Gemini 3.8 Flash, 0,0\,\% --- puis Kimi K3 et Gemini 3.1 Pro (0,3\,\%), DeepSeek V4 Pro et Qwen3.8 (0,6\,\%), DeepSeek V4 Flash (0,9\,\%) et GPT-5.6 Luna, 4,1\,\%, qui concentre treize des vingt-deux occurrences. L'observation qui s'en dégage : le taux ne suit ni la taille ni le rang de qualité --- GPT-5.6 Luna est économe et fabricateur, Gemini 3.1 Pro lourd et presque propre. L'hypothèse la plus plausible est l'exposition du corpus d'entraînement à la nomenclature ICS, seule zone où les fabrications surviennent.

\textbf{La mauvaise matrice.} Zéro chez les quatre mêmes modèles ; de un à neuf identifiants chez les six autres (Gemini 3.1 Pro 1, Kimi K3 3, DeepSeek V4 Flash 4, DeepSeek V4 Pro 5, Qwen3.8 7, GPT-5.6 Luna 9). Que la liste des modèles propres soit identique sur les deux fautes n'est pas un hasard : elles mesurent la même compétence --- situer l'univers avant de nommer --- l'une échouant vers un identifiant existant, l'autre vers un identifiant inventé.

\textbf{La stabilité.} Gemini 3.1 Pro est le plus constant (88,1\,\%), suivi de Gemini 3.8 Flash (81,0\,\%) et de Claude Opus 4.8 (76,2\,\%). Le contre-exemple éclaire l'axe : GPT-5.6 Sol, premier en qualité stricte, n'est que 64,3\,\% stable --- ses bons scores sont une moyenne de copies qui se contredisent, et en exploitation on ne relance pas trois fois.


\textbf{Les identifiants obsolètes se concentrent sur trois attaques.} Trois échantillons concentrent l'intégralité des identifiants révoqués. S027 --- effacement des journaux système Linux --- est mappé T1070.002 par les dix modèles, aux trois exécutions : cet identifiant a été révoqué au profit de T1685.006, qui porte exactement le même nom. S007 --- message Modbus non autorisé --- est mappé T0855 par huit modèles sur dix, révoqué au profit de T1692.001. S038 --- clés de registre d'exécution automatique --- n'est touché que par GPT-5.6 Sol, qui propose T1060, révoqué au profit de T1547.001. Dans les trois cas le raisonnement est juste et seul le numéro a changé : c'est précisément ce que la base tolérante crédite.

\textbf{Comparaison sur les trois modèles retenus.} Sur S027, les trois se comportent à l'identique : T1070.002 aux trois exécutions, soit un faux positif au niveau technique et un obsolète au niveau sous-technique. Sur S007, Gemini 3.8 Flash propose T0855 aux trois exécutions --- donc crédité en base tolérante --- tandis que GPT-5.6 Sol propose T0836, une technique vivante mais fausse, jamais rattrapable ; Claude Opus 4.8 est entre les deux, T0836 deux fois et T0855 une fois. Sur S038, Claude Opus 4.8 et Gemini 3.8 Flash répondent juste aux trois exécutions, GPT-5.6 Sol deux fois T1060 puis une fois la bonne réponse. L'écart entre base stricte et base tolérante ne mesure donc pas une différence de raisonnement, mais la fraîcheur du référentiel interne.

\textbf{Ce que l'ancrage par MCP change, sur Claude Opus 4.8.} Rejouée avec le serveur MCP, la même campagne corrige les deux échantillons où le modèle était en défaut. Sur S027, les trois exécutions passent de T1070/T1070.002 à T1685/T1685.006 ; sur S007, de T0836 ou T0855 à T1692/T1692.001. S038, déjà juste, le reste. Le taux d'obsolescence tombe de 1,3\,\% à 0\,\% et le F1 progresse à tous les niveaux : 0,587 à 0,802 en tactique, 0,563 à 0,706 en technique, 0,636 à 0,758 en sous-technique. L'ancrage ne rend pas le modèle plus perspicace, il lui donne le référentiel courant.

\textbf{Les hallucinations sont un phénomène tactique et exclusivement ICS.} Les 23 hallucinations du panel sont toutes au niveau tactique, et toutes sur des échantillons ICS. Les identifiants fabriqués --- TA0113 à TA0119, plus TA0089 et TA0024 --- se situent juste au-delà de la plage réelle des tactiques ICS, qui s'arrête à TA0111. Les modèles extrapolent la numérotation au lieu de la connaître : ils savent que les tactiques ICS commencent à TA0100 et poursuivent la série. Les tactiques concernées sont Impair Process Control (S006, S007, S008), Collection (S002, S004, S005), Discovery (S001), Inhibit Response Function (S009, S010) et Persistence (S013) --- soit les tactiques propres à l'univers industriel, celles qui n'ont pas d'équivalent Enterprise dont le modèle pourrait se souvenir.

\textbf{L'homonymie entre matrices élimine des modèles par ailleurs corrects.} Neuf tactiques portent le même nom dans les deux matrices : Collection, Command and Control, Discovery, Execution, Impact, Initial Access, Lateral Movement, Persistence et Privilege Escalation. Un modèle qui reconnaît correctement l'intention mais ignore qu'il lit un journal industriel produit l'identifiant Enterprise homonyme --- TA0007 au lieu de TA0102 pour Discovery, TA0009 au lieu de TA0100 pour Collection. Quatre échantillons sont touchés : S000 et S001 (Discovery), S003 (Collection) et S012 (Execution). Trois modèles commettent cette confusion : Qwen3.8 quatre fois, DeepSeek V4 Pro et GPT-5.6 Luna trois fois chacun. Le raisonnement est bon, l'univers est faux, et la mitigation qui en découlerait serait inapplicable au système observé --- d'où son rang de premier critère d'élimination.

\textbf{Le cas de Claude Opus 4.8.} Claude Opus 4.8 ne produit aucun identifiant hors matrice ni aucune hallucination, avec ou sans serveur. Il ne doit donc pas son maintien dans le panel à l'ancrage, mais à une connaissance correcte de la frontière entre les deux univers. L'apport du MCP porte ailleurs, sur l'actualité du référentiel.

\textbf{Les tactiques concernées par l'obsolescence chez Claude Opus 4.8.} Deux tactiques du corpus concentrent les identifiants obsolètes de Claude Opus 4.8, et serviront d'exemples au chapitre 2. TA0112 Defense Impairment, matrice Enterprise, par l'échantillon S027 --- effacement des journaux système Linux --- où T1070.002 est proposé aux trois exécutions alors que l'identifiant courant est T1685.006. Et TA0106 Impair Process Control, matrice ICS, par l'échantillon S007 --- message Modbus non autorisé --- où T0855 est proposé une fois sur trois alors que l'identifiant courant est T1692.001. Les deux cas relèvent d'une révocation avec successeur officiel, donc résolubles par simple consultation du référentiel.

\begin{longtable}[]{@{}
  >{\raggedright\arraybackslash}p{(\columnwidth - 10\tabcolsep) * \real{0.2778}}
  >{\raggedright\arraybackslash}p{(\columnwidth - 10\tabcolsep) * \real{0.1603}}
  >{\raggedright\arraybackslash}p{(\columnwidth - 10\tabcolsep) * \real{0.1389}}
  >{\raggedright\arraybackslash}p{(\columnwidth - 10\tabcolsep) * \real{0.1389}}
  >{\raggedright\arraybackslash}p{(\columnwidth - 10\tabcolsep) * \real{0.1421}}
  >{\raggedright\arraybackslash}p{(\columnwidth - 10\tabcolsep) * \real{0.1421}}@{}}
\toprule\noalign{}
\begin{minipage}[b]{\linewidth}\raggedright
\textbf{Modèle}
\end{minipage} & \begin{minipage}[b]{\linewidth}\raggedright
\textbf{Enterprise}
\end{minipage} & \begin{minipage}[b]{\linewidth}\raggedright
\textbf{ICS}
\end{minipage} & \begin{minipage}[b]{\linewidth}\raggedright
\textbf{Total O}
\end{minipage} & \begin{minipage}[b]{\linewidth}\raggedright
\textbf{Affirm.}
\end{minipage} & \begin{minipage}[b]{\linewidth}\raggedright
\textbf{Taux}
\end{minipage} \\
\midrule\noalign{}
\endhead
\bottomrule\noalign{}
\endlastfoot
Claude Opus 4.8 & 3 & 1 & 4 & 318 & 1,3\,\% \\
Claude Sonnet 5 & 3 & 3 & 6 & 313 & 1,9\,\% \\
GPT-5.6 Sol & 5 & 0 & 5 & 318 & 1,6\,\% \\
GPT-5.6 Luna & 3 & 2 & 5 & 316 & 1,6\,\% \\
Gemini 3.1 Pro & 3 & 3 & 6 & 315 & 1,9\,\% \\
Gemini 3.8 Flash & 3 & 3 & 6 & 316 & 1,9\,\% \\
Kimi K3 & 3 & 2 & 5 & 318 & 1,6\,\% \\
DeepSeek V4 Pro & 3 & 2 & 5 & 318 & 1,6\,\% \\
DeepSeek V4 Flash & 3 & 2 & 5 & 318 & 1,6\,\% \\
Qwen3.8 & 3 & 0 & 3 & 318 & 0,9\,\% \\
\textbf{Total} & \textbf{32} & \textbf{18} & \textbf{50} & \textbf{3 168} & \textbf{1,6\,\%} \\
\end{longtable}
\legendetableau{Répartition des identifiants obsolètes par matrice et taux d'obsolescence, par modèle.}{tab:obsoletes}


\textbf{Lecture du tableau.} Les 32 identifiants obsolètes Enterprise proviennent presque tous du même échantillon, S027, sur lequel les dix modèles se trompent aux trois exécutions ; GPT-5.6 Sol y ajoute S038. Les 18 identifiants ICS proviennent tous de S007. L'obsolescence n'est donc pas diffuse : elle se concentre sur les rares techniques renommées entre deux versions, et frappe indifféremment les dix modèles. Aucun ne se distingue par un référentiel plus à jour --- Qwen3.8 affiche le taux le plus bas, 0,9\,\%, non parce qu'il connaît T1685.006, mais parce qu'il se trompe autrement sur S007.

\textbf{La base tolérante ne déplace que la tête du classement.} Gemini 3.8 Flash passe devant GPT-5.6 Sol (0,624 contre 0,619 en technique) : une part plus grande de ses erreurs était constituée d'identifiants révoqués dont le successeur officiel était la bonne réponse. Les quatre modèles retenus restent les mêmes, et l'élimination ne bouge pas, puisque ni la matrice ni l'hallucination ne dépendent de la base. Le gain au passage en base tolérante mesure l'actualité du référentiel interne de chaque modèle : de 0,008 pour Claude Opus 4.8 à 0,024 pour Gemini 3.8 Flash et Claude Sonnet 5. Qwen3.8 est le seul dont le F1 technique ne bouge pas d'un millième --- son référentiel est le plus à jour du panel ---, mais ses sept identifiants hors matrice l'ont déjà écarté.

\textbf{Ce que ces deux tableaux servent à la comparaison avec le MCP.} Ils constituent la ligne de base de la phase 2 : mêmes échantillons, même protocole, mêmes modèles, mais sans serveur. Trois lectures en découlent. D'abord, c'est la base tolérante qui sera la référence à battre : la résolution des identifiants révoqués est précisément la fonction que le serveur mécanise, et comparer les résultats avec MCP à la seule base stricte confondrait deux effets --- le rattrapage de version, déjà mesuré par l'écart entre les tableaux~\ref{tab:classement-strict} et~\ref{tab:classement-tolerant}, et l'apport propre de l'ancrage, c'est-à-dire consulter le référentiel avant de répondre. Tout dépassement du tableau~\ref{tab:classement-tolerant} mesurera l'ancrage, et non le simple rafraîchissement des connaissances. Ensuite, les deux colonnes de tête deviennent les indicateurs directs de l'apport du serveur : celui-ci ne sert que des identifiants vivants, pris dans la matrice interrogée, de sorte que les vingt-neuf identifiants hors matrice et les vingt-deux identifiants fabriqués relevés ici doivent disparaître avec lui ; s'ils subsistent, c'est que l'ancrage n'a pas été utilisé. Enfin, l'écart entre les deux zones désigne le terrain où le gain est attendu : au niveau technique, les dix modèles se tiennent entre 0,639 et 0,750 sur la matrice Enterprise, mais entre 0,167 et 0,452 sur la matrice ICS. La connaissance du domaine industriel est la plus datée et la plus lacunaire ; c'est là que le référentiel outillé a le plus à apporter.

\subsubsection{Fiabilité, temps de traitement et coût}\label{sec:5-2-ancien}
Aucun indicateur du tableau~\ref{tab:fiabilite} ne dépend de la base de calcul, à l'exception de la dernière colonne : l'efficacité rapporte le F1 technique strict au coût en jetons.

\begin{longtable}[]{@{}
  >{\raggedright\arraybackslash}p{(\columnwidth - 18\tabcolsep) * \real{0.0700}}
  >{\raggedright\arraybackslash}p{(\columnwidth - 18\tabcolsep) * \real{0.2144}}
  >{\raggedright\arraybackslash}p{(\columnwidth - 18\tabcolsep) * \real{0.0880}}
  >{\raggedright\arraybackslash}p{(\columnwidth - 18\tabcolsep) * \real{0.1016}}
  >{\raggedright\arraybackslash}p{(\columnwidth - 18\tabcolsep) * \real{0.1106}}
  >{\raggedright\arraybackslash}p{(\columnwidth - 18\tabcolsep) * \real{0.0926}}
  >{\raggedright\arraybackslash}p{(\columnwidth - 18\tabcolsep) * \real{0.0880}}
  >{\raggedright\arraybackslash}p{(\columnwidth - 18\tabcolsep) * \real{0.0700}}
  >{\raggedright\arraybackslash}p{(\columnwidth - 18\tabcolsep) * \real{0.0700}}
  >{\raggedright\arraybackslash}p{(\columnwidth - 18\tabcolsep) * \real{0.0948}}@{}}
\toprule\noalign{}
\begin{minipage}[b]{\linewidth}\raggedright
\textbf{Rang}
\end{minipage} & \begin{minipage}[b]{\linewidth}\raggedright
\textbf{Modèle}
\end{minipage} & \begin{minipage}[b]{\linewidth}\raggedright
\textbf{Stab.}
\end{minipage} & \begin{minipage}[b]{\linewidth}\raggedright
\textbf{p50 (ms)}
\end{minipage} & \begin{minipage}[b]{\linewidth}\raggedright
\textbf{p95 (ms)}
\end{minipage} & \begin{minipage}[b]{\linewidth}\raggedright
\textbf{Jetons}
\end{minipage} & \begin{minipage}[b]{\linewidth}\raggedright
\textbf{Sur-sp.}
\end{minipage} & \begin{minipage}[b]{\linewidth}\raggedright
\textbf{N1}
\end{minipage} & \begin{minipage}[b]{\linewidth}\raggedright
\textbf{N2}
\end{minipage} & \begin{minipage}[b]{\linewidth}\raggedright
\textbf{Effic.}
\end{minipage} \\
\midrule\noalign{}
\endhead
\bottomrule\noalign{}
\endlastfoot
1 & Gemini 3.1 Pro & 88,1\,\% & 18 532 & 103 824 & 9 023 & 25,0\,\% & 0\,\% & 0\,\% & 0,053 \\
2 & Gemini 3.8 Flash & 81,0\,\% & 10 664 & 44 482 & 7 566 & 10,2\,\% & 0\,\% & 0\,\% & 0,079 \\
3 & Claude Opus 4.8 & 76,2\,\% & 4 282 & 6 050 & 6 572 & 10,0\,\% & 0\,\% & 0\,\% & 0,086 \\
4 & Claude Sonnet 5 & 66,7\,\% & 4 291 & 10 994 & 6 740 & 15,0\,\% & 0\,\% & 0\,\% & 0,080 \\
5 & GPT-5.6 Sol & 64,3\,\% & 6 086 & 24 145 & 4 649 & 10,0\,\% & 0\,\% & 0\,\% & 0,130 \\
6 & Kimi K3 & 59,5\,\% & 24 034 & 202 504 & 6 027 & 10,0\,\% & 0\,\% & 0\,\% & 0,094 \\
7 & Qwen3.8 & 57,1\,\% & 82 854 & 269 617 & 9 641 & 11,7\,\% & 0\,\% & 0\,\% & 0,059 \\
8 & GPT-5.6 Luna & 54,8\,\% & 4 626 & 28 758 & 5 004 & 26,7\,\% & 0\,\% & 0\,\% & 0,105 \\
9 & DeepSeek V4 Flash & 52,4\,\% & 3 407 & 19 812 & 5 594 & 11,7\,\% & 0\,\% & 0\,\% & 0,108 \\
10 & DeepSeek V4 Pro & 42,9\,\% & 40 724 & 151 756 & 7 941 & 18,3\,\% & 0\,\% & 0\,\% & 0,072 \\
\end{longtable}
\legendetableau{Classement de la fiabilité, du temps de traitement et du coût.}{tab:fiabilite}

Les jetons du tableau sont une unité physique, pas une facture : chaque fournisseur fixe son tarif, et ces tarifs changent en quelques semaines. Aux tarifs de septembre 2026 (entrée/sortie en \$ par million de jetons : GPT-5.6 Sol 4/20, Gemini 3.8 Flash 0,75/3,75, Claude Opus 4.8 5/25, Claude Sonnet 5 2/10), les 126 exécutions coûtent exactement 3,09, 1,36, 4,63 et 2,06\,\$ respectivement (répartitions entrée/sortie réelles de chaque exécution, calcul du script) : l'ordre budgétaire ne suit ni les jetons ni le temps, d'où le choix d'unités physiques dans le tableau, complétées par ce rappel daté.


\textbf{Deux modèles écartés avant la campagne.} Deux modèles ont été écartés après essais. kimi-k2.7-code affichait 158 s de latence médiane, soit plus de cinq heures de campagne contre 50 minutes pour kimi-k3 de la même lignée. qwen3.5-35b-a3b consommait 47 626 jetons par appel, dont 94\,\% en raisonnement, avec deux sorties non conformes sur huit essais. Dans les deux cas, la dépense n'était pas justifiée par la qualité. Ce n'est pas un effet de taille --- DeepSeek V4 Flash, modèle léger, consomme moins que DeepSeek V4 Pro tout en obtenant de meilleurs résultats --- mais de longueur de raisonnement.

\textbf{Pourquoi aucun niveau de raisonnement n'est imposé.} L'effort de raisonnement est laissé aux valeurs par défaut et n'est pas traité comme facteur : aucune échelle commune n'existe entre éditeurs, et forcer un niveau « maximal » partout déplacerait les modèles vers des points arbitraires différents plutôt que comparables. Même logique que pour la température. Le faire varier changerait la question posée et porterait la campagne à 3 780 appels. L'effet reste mesuré : les jetons de raisonnement sont relevés et rapportés --- c'est ce qui a permis d'écarter les deux modèles ci-dessus.

\textbf{La stabilité ne mesure pas la justesse.} Gemini 3.1 Pro affiche la meilleure du panel, 88,1\,\%, tout en étant dernier en F1 technique en base stricte : il se trompe de façon remarquablement reproductible. Sa sur-spécification de 25,0\,\% --- une sous-technique inventée sur un quart des techniques terminales --- et sa médiane de 18,5 s achèvent de le disqualifier, et il avait de toute façon été écarté par la matrice.

\textbf{Parmi les quatre modèles retenus, la hiérarchie s'inverse d'un critère à l'autre.} Gemini 3.8 Flash est le plus stable (81,0\,\%), devant Claude Opus 4.8 (76,2\,\%), Claude Sonnet 5 (66,7\,\%) et GPT-5.6 Sol (64,3\,\%). Sur la vitesse, l'ordre se renverse presque : les deux Claude traitent un échantillon en 4,3 s de médiane, GPT-5.6 Sol en 6,1 s, Gemini 3.8 Flash en 10,7 s.

\textbf{Le 95e centile sépare mieux que la médiane.} Claude Opus 4.8 est à 6,0 s, le plus bas des dix modèles, soit 1,4 fois sa propre médiane ; Claude Sonnet 5 est à 11,0 s (2,6 fois), GPT-5.6 Sol à 24,1 s (4,0 fois) et Gemini 3.8 Flash à 44,5 s (4,2 fois). C'est ce rapport, plus que la médiane seule, qui dit si un centre opérationnel peut planifier sa charge : il mesure le pire cas ordinaire, celui qu'une moyenne masquerait.

\textbf{Le coût sépare franchement.} GPT-5.6 Sol consomme 4 649 jetons par exécution, Claude Opus 4.8 6 572, Claude Sonnet 5 6 740 et Gemini 3.8 Flash 7 566 --- soit 63\,\% de plus que GPT-5.6 Sol. L'efficacité classe donc presque à l'inverse de la qualité brute : 0,130 pour GPT-5.6 Sol, 0,086 pour Claude Opus 4.8, 0,080 pour Claude Sonnet 5 et 0,079 pour Gemini 3.8 Flash.

\textbf{Deux colonnes valent zéro pour les dix modèles.} N1 et N2 sont neutralisées par le format de sortie, qui impose les trois champs : la sous-spécification est empêchée par le protocole, non par une vertu des modèles. Les colonnes sont conservées parce que cette nullité est elle-même un résultat.

\subsubsection{Choix du modèle}\label{sec:5-3-ancien}
\textbf{Le coût en argent.} Les jetons ne sont pas la dépense, ils n'en sont que l'unité : chaque fournisseur fixe son tarif, et seul le montant facturé se paie. Aux tarifs du 16 septembre 2026, la campagne coûte 1,36\,\$ sur Gemini 3.8 Flash, 2,06\,\$ sur Claude Sonnet 5, 3,09\,\$ sur GPT-5.6 Sol et 4,63\,\$ sur Claude Opus 4.8 --- et, hors des quatre finalistes, 0,12\,\$ sur DeepSeek V4 Flash, 0,22\,\$ sur GPT-5.6 Luna, 0,62\,\$ sur DeepSeek V4 Pro. Le classement en argent inverse celui des jetons : le plus gourmand des quatre en jetons est le moins cher en dollars, son tarif compensant six fois sa verbosité. Deux échéances datent ces chiffres : le tarif de Sol est promotionnel, et celui de Gemini 3.8 Flash double au 1\up{er} janvier 2027.

\textbf{Le temps de traitement.} En médiane : DeepSeek V4 Flash (3,4 s), Claude Opus 4.8 et Claude Sonnet 5 (4,3 s). Le 95\up{e} centile sépare les prévisibles des erratiques : Claude Opus 4.8 y est seul de son espèce --- 6,1 s, soit 1,4 fois sa médiane, le plus bas des dix --- quand Gemini 3.8 Flash monte à 44,5 s et Qwen3.8 à 270 s. Dans un centre qui traite des milliers d'alertes, le p95 est le temps que l'analyste attend vraiment les mauvais jours.

\textbf{Pourquoi des consommations si différentes à entrées comparables.} Tous les modèles reçoivent le même journal et rendent le même triplet JSON : l'entrée et la sortie utile sont presque identiques d'un modèle à l'autre. L'écart --- de 4\,649 jetons par exécution chez GPT-5.6 Sol à 9\,641 chez Qwen3.8 --- vient pour l'essentiel des jetons de raisonnement : la réflexion interne que la plupart des fournisseurs facturent comme de la sortie sans la restituer intégralement. Kimi K3 raisonne systématiquement, sans réglage pour l'éteindre ; Qwen3.8 et Gemini 3.8 Flash raisonnent par défaut. À tarif de sortie élevé, ce raisonnement invisible devient la part dominante de la facture --- c'est lui, et non la réponse rendue, qui sépare les sobres des gourmands.

\textbf{Le modèle retenu.} Sur six axes --- qualité, matrice, hallucination, stabilité, argent, temps --- les présences aux podiums donnent Gemini 3.8 Flash et Claude Opus 4.8 à quatre, GPT-5.6 Sol, Claude Sonnet 5 et DeepSeek V4 Flash à trois. L'égalité de tête se départage par la nature des présences : celles de Claude Opus 4.8 couvrent toute la fiabilité --- matrice, hallucination, stabilité --- plus le temps ; celles de Gemini 3.8 Flash s'arrêtent avant le temps, l'axe qui se multiplie par le volume, et que la phase 2 alourdira encore en ajoutant au contexte les appels d'outils et leurs résultats. Ses absences à lui --- le podium du mapping, manqué de 0,040, et l'argent --- sont les moins graves pour une campagne d'évaluation. La phase 2 cherche un instrument : peu dispersé pour laisser l'effet de l'ancrage visible, au temps prévisible, et porteur d'une marge à combler --- ce F1 technique qui tombe de 0,714 sur Enterprise à 0,262 sur ICS. Claude Opus 4.8 est retenu ; les tableaux ci-dessus constituent sa ligne de base sans serveur, celle contre laquelle l'apport de l'ancrage sera mesuré.


\fi
\FloatBarrier

%% ===================================================================
%% CHAPITRE 2 --- L'approche MCP
%% ===================================================================
\chapter{L'approche MCP}\label{chap:approche-mcp}

\section{Introduction}\label{sec:ch2-1}

\emph{[Section à rédiger.]}

\section{Le serveur MCP}\label{sec:ch2-2}

MITRE ATT\&CK est un catalogue public de comportements d'attaquants informatiques, publié sous la forme d'un fichier de données que chacun peut télécharger. Notre serveur MCP est le connecteur qui expose les matrices Enterprise et ICS de ce référentiel aux modèles d'intelligence artificielle. Voici comment le tout fonctionne.

\subsection{Architecture : hôte, client et serveur}\label{sec:ch2-2-1}

Le MCP relie trois rôles : l'hôte porte le modèle, le client y maintient une connexion un-à-un avec chaque serveur, et le serveur expose les outils que le modèle invoque --- lequel ne s'adresse jamais directement au serveur (figure~\ref{fig:ch2-archi}).

\begin{figure}[htbp]
\centering
\figuretikz{fig-2-1-architecture}
\caption{Architecture MCP. Le modèle raisonne dans l'hôte ; le client porte la connexion ; le serveur répond à partir de ses deux sources, en lecture seule.}\label{fig:ch2-archi}
\end{figure}

\subsection{Les deux sources du serveur}\label{sec:ch2-2-2}

Le serveur sert deux sources, l'une publique, l'autre propre à l'organisation.

\textbf{Le référentiel MITRE.} La base est chargée depuis les paquets STIX publiés par MITRE, où chaque élément est un objet --- tactique, technique, sous-technique, mitigation, groupe, logiciel, campagne --- et où chaque lien est lui-même un objet (figure~\ref{fig:ch2-referentiel}) : la technique appartient à une tactique de la \emph{kill chain}, \texttt{subtechnique-of} rattache \emph{LSASS Memory} à sa parente, \texttt{mitigates} relie une mesure de défense à la technique, \texttt{uses} relie un groupe à la sous-technique qu'il emploie, et \texttt{revoked-by} conduit d'un identifiant retiré à son remplaçant vivant. La fiche renvoyée par \texttt{get\_technique} est assemblée en suivant ces liens : tactique, mitigations propres et héritées de la parente, groupes, sources de détection.

\begin{figure}[htbp]
\centering
\figuretikz{fig-2-2-referentiel}
\caption{Le référentiel MITRE tel que le serveur le lit : des objets reliés par des relations explicites.}\label{fig:ch2-referentiel}
\end{figure}

\textbf{Le référentiel interne.} Un fichier JSON de fiches aux identifiants propres à l'organisation, chacune rattachée à une technique MITRE dont le statut est signalé. Voici une fiche réelle du référentiel :

\begin{verbatim}
{"id": "THQ0001",
 "name": "Phishing imitant le portail RH interne",
 "description": "Leurre imitant le portail RH pour voler des identifiants.",
 "technique_mitre_liee": "T1566.002",
 "tactique": "initial-access",
 "source_cti": "incident-2026-03",
 "statut": "valide"}
\end{verbatim}

\subsection{Le chemin d'une requête}\label{sec:ch2-2-3}

Un analyste soumet un signal au modèle : « Comment interprètes-tu ce signal avec le référentiel MITRE ATT\&CK ? », accompagné d'une trace Zeek où une même adresse ouvre, en quelques secondes, des connexions très brèves vers le port 502 --- celui de Modbus --- de sept adresses voisines. Le modèle décide d'interroger le serveur ; l'appel passe par le client MCP, qui porte la connexion. La figure~\ref{fig:ch2-chemin} montre le chemin parcouru dans le serveur, et chaque étape est développée ci-dessous en suivant cette même requête, de son émission à sa trace d'audit. La requête est réelle : c'est la première recherche de l'exécution 1 de l'échantillon S000, celle qui ouvre notre campagne d'évaluation, et toutes les valeurs qui suivent sont tirées telles quelles du journal d'audit du serveur, les champs longs étant abrégés par \texttt{.....} La figure~\ref{fig:ch2-echanges}, plus loin, resituera cet appel parmi les quatorze échanges de cette même exécution.

\begin{figure}[htbp]
\centering
\figuretikz{fig-2-3-chemin}
\caption{Chemin d'une requête dans le serveur MCP.}\label{fig:ch2-chemin}
\end{figure}

\begin{itemize}
\tightlist
\item
  \textbf{Format des messages.} Les échanges suivent JSON-RPC 2.0, le format retenu par la spécification MCP : un objet JSON portant un identifiant, la méthode demandée et ses paramètres. Le serveur porte ce dialogue sur l'entrée et la sortie standard du processus, et sur un service HTTP local pour les clients qui ne savent pas parler ainsi ; les deux voies traversent le même code et ne changent rien à ce qui suit. Pour notre signal, le modèle formule une requête de recherche dans le vocabulaire ATT\&CK et le client émet :
\end{itemize}

\begin{verbatim}
{"jsonrpc": "2.0", "id": "toolu_01Qmsp.....",
 "method": "tools/call",
 "params": {"name": "search_techniques",
            "arguments": {
              "query": "network scan modbus port 502 enumerate devices",
              "matrix": "ICS"}}}
\end{verbatim}

\begin{itemize}
\tightlist
\item
  \textbf{Décodage.} Le champ \texttt{method} détermine le traitement. Les noms de méthodes sont imposés par la spécification MCP et le dialogue typique en enchaîne trois : \texttt{initialize} (« bonjour, qui es-tu ? »), \texttt{tools/list} (« que sais-tu faire ? »), puis \texttt{tools/call} (« exécute \texttt{search\_techniques} avec ce mot-clé »). Une quatrième, \texttt{ping}, peut être émise par l'un ou l'autre côté à tout moment : le destinataire doit répondre aussitôt par un résultat vide, ce qui atteste que la connexion est ouverte et le pair réactif ; sans réponse, l'émetteur la considère morte et la ferme. C'est par \texttt{tools/list}, en début de session, que le modèle apprend quels outils existent et quand les employer (figure~\ref{fig:ch2-outils}) ; le tableau~\ref{tab:ch2-outils} en donne la liste. Une ligne illisible est journalisée puis écartée, sans bloquer les requêtes suivantes.
\end{itemize}

\begin{figure}[htbp]
\centering
\figuretikz{fig-2-4-enumeration}
\caption{L'énumération des outils : le modèle lit là les dix-huit descriptions, figées dans le code, qui guident ensuite son choix d'outil.}\label{fig:ch2-outils}
\end{figure}


\begin{longtable}[]{@{}
  >{\raggedright\arraybackslash}p{(\columnwidth - 2\tabcolsep) * \real{0.3369}}
  >{\raggedright\arraybackslash}p{(\columnwidth - 2\tabcolsep) * \real{0.6631}}@{}}
\toprule\noalign{}
\begin{minipage}[b]{\linewidth}\raggedright
\textbf{Outil}
\end{minipage} & \begin{minipage}[b]{\linewidth}\raggedright
\textbf{Description exposée au modèle}
\end{minipage} \\
\midrule\noalign{}
\endhead
\bottomrule\noalign{}
\endlastfoot
\texttt{search\_techniques} & Recherche lexicale de techniques candidates ; une requête pour un comportement, résultats scorés avec leurs preuves \\
\texttt{get\_technique} & Détail complet d'une technique pour confirmer un candidat : tactiques, plateformes, détections, sous-techniques, groupes, logiciels, campagnes, mitigations propres et de la parente, équipements ICS ; redirige un identifiant révoqué, signale un déprécié \\
\texttt{get\_tactics} & Les tactiques Enterprise et ICS dans l'ordre de la \emph{kill chain} \\
\texttt{get\_tactic} & Une tactique et ses techniques, parentes et sous-techniques \\
\texttt{get\_mitigations} & Les contre-mesures associées à une technique, pour produire une recommandation \\
\texttt{get\_mitigation} & Une mesure de mitigation et toutes les techniques qu'elle couvre \\
\texttt{get\_groups} & Groupes APT : les plus actifs, ou ceux qui emploient une technique donnée, directement ou en campagne \\
\texttt{get\_group} & Détail d'un groupe : alias, techniques directes, techniques observées en campagne, logiciels, campagnes attribuées \\
\texttt{get\_software} & Logiciels et maliciels : les plus employés, ou ceux associés à une technique \\
\texttt{get\_campaign} & Détail d'une campagne : période, groupes attribués, techniques et logiciels observés \\
\texttt{get\_asset} & Équipement ICS --- automate, interface, unité distante, historien --- et les techniques qui le visent \\
\texttt{get\_datasources} & Techniques détectables depuis une source de journaux donnée \\
\texttt{list\_internal\_techniques} & Référentiel interne : liste complète et prochain identifiant libre, à appeler avant toute proposition de fiche \\
\texttt{get\_internal\_technique} & Détail d'une fiche interne, avec la technique MITRE liée et son statut \\
\texttt{search\_internal} & Recherche lexicale dans le référentiel interne ; obligatoire avant de proposer une fiche, pour écarter un doublon \\
\texttt{internal\_reload} & Relecture du référentiel interne sans redémarrer le serveur ; le fichier n'est jamais écrit \\
\texttt{mitre\_stats} & Version du serveur, versions ATT\&CK chargées et statistiques de la base \\
\texttt{mitre\_update} & Retéléchargement des deux matrices, avec bascule seulement si les deux réussissent \\
\end{longtable}
\legendetableau{Les dix-huit outils exposés et leur description, telle que le modèle la lit.}{tab:ch2-outils}

\begin{itemize}
\tightlist
\item
  \textbf{Un décodage qui ne se tait jamais.} Un client mal réglé ne doit pas produire un silence, car le silence laisse le modèle attendre une réponse qui ne viendra pas. Trois défauts courants sont donc traités explicitement. Un message étalé sur plusieurs lignes est assemblé dans un tampon borné à un mégaoctet, les accolades étant comptées hors des chaînes de caractères : un \texttt{\}} présent dans un extrait de journal ne fait plus croire que le message est terminé. Ce qui reste illisible reçoit l'erreur \texttt{-32700}. Une enveloppe mal formée --- message qui n'est pas un objet, \texttt{params} qui n'en est pas un, nom d'outil qui n'est pas du texte, \texttt{arguments} qui n'est pas un objet --- est refusée avec le code correspondant, \texttt{-32600} ou \texttt{-32602}. Enfin la version de protocole est négociée au lieu d'être répétée : le serveur ne rend la version demandée que s'il la connaît, et répond sinon la plus ancienne qu'il prend en charge, \texttt{2024-11-05}, plutôt que d'en accepter une qu'il ne saurait pas parler. La règle est la même partout : toute requête portant un identifiant reçoit une réponse, y compris quand elle est fautive.
\item
  \textbf{Aucun état conservé entre deux appels.} Le serveur ne retient ni session, ni historique, ni contexte de conversation, et n'exige aucun identifiant de session. Chaque message porte tout ce qu'il faut pour être traité, et plusieurs messages peuvent d'ailleurs voyager ensemble : un lot JSON-RPC reçoit un tableau de réponses, de sorte que trois outils demandés dans le même tour partent en un seul envoi. Deux conséquences : la mémoire de l'échange est portée par le client, et la même requête rejouée sur la même version de base rend exactement la même réponse, puisque le résultat ne dépend pas de ce qui a été demandé avant.
\item
  \textbf{Validation.} Avant toute exécution, l'appel est contrôlé contre le schéma que l'outil déclare : l'outil existe, les arguments requis sont présents, aucun argument étranger n'est passé, les valeurs sont du bon type et d'une longueur bornée. Un filtre mal orthographié est refusé avec la liste des valeurs admises, et non traité comme une recherche sans résultat : une liste vide signifierait « aucune correspondance » et fausserait la conclusion.
\end{itemize}

\begin{verbatim}
get_technique(id="T1003.001", matrix="ICS")
-> argument inconnu : matrix (attendus : id)

search_techniques("phishing", matrix="entreprise")
-> matrice 'entreprise' inconnue - valeurs acceptees :
  Enterprise, ICS (ou IT, OT, SCADA)

search_techniques("phishing", platform="windowz")
-> platform inconnue 'windowz' : attendu une valeur parmi Containers,
  ESXi, IaaS, Identity Provider, Linux, Network Devices, None,
  Office Suite, PRE, SaaS, Windows, macOS
\end{verbatim}

\begin{itemize}
\tightlist
\item
  Quand la correction est évidente, le serveur corrige et le dit plutôt que d'échouer : un \texttt{"limit":\ "abc"} retombe sur la valeur par défaut, avec la note \texttt{limit\ invalide\ :\ valeur\ par\ defaut\ appliquee}.
\item
  \textbf{Moteur : proposer, puis vérifier.} Le moteur ne conclut pas, il propose. Il rend des techniques candidates, chacune avec un score et avec les mots qui l'ont fait remonter ; c'est ensuite au modèle de confirmer le candidat retenu par \texttt{get\_technique}, puis de le confronter au référentiel interne. La section~\ref{sec:ch2-2-4} détaille son fonctionnement.
\item
  \textbf{Référentiel MITRE.} Le moteur y cherche les techniques candidates, puis y assemble la fiche complète de celle qui est retenue.
\item
  \textbf{Référentiel interne.} Il y vérifie ensuite si l'organisation décrit déjà ce comportement sous son propre identifiant.
\item
  \textbf{Sérialisation.} Le résultat est mis en forme en JSON compact et la version de la base servie y est jointe : la réponse est datée, donc rejouable. Voici le contenu renvoyé au client pour notre requête, avec les champs réels du serveur (enveloppe JSON-RPC omise, premier candidat seul) :
\end{itemize}

\begin{verbatim}
{"db_version": {"Enterprise": "19.2", "ICS": "19.2"},
 "query": "network scan modbus port 502 enumerate devices",
 "query_tokens": ["502", "devic", "enumer", "modbus", .....],
 "filtres": {"matrix": "ICS", "platform": "aucun",
             "tactic": "aucun"},
 "seuils": {"fort": 0.6, "faible": 0.3},
 "result_count": 10,
 "total_matches": 72,
 "results": [
   {"id": "T0846.001", "name": "Port Scan",
    "full_name": "Remote System Discovery: Port Scan",
    "matrix": "ICS", "score": 0.364,
    "is_subtechnique": true, "parent_id": "T0846",
    "tactics": ["TA0102"],
    "matched": {.....},
    "snippet": "....."}, .....],
 "note": "le candidat de tete est une SOUS-technique ; sa
          parente T0846 est aussi candidate - retenir T0846
          si l'evidence ne distingue pas la sous-technique"}
\end{verbatim}

\begin{itemize}
\tightlist
\item
  \textbf{Journal d'audit.} Vingt-quatre types d'événements y sont tracés, et pas seulement l'appel d'outil. D'abord la base elle-même : le téléchargement de chaque matrice avec sa durée et le numéro de version obtenu, puis son chargement, dont le champ \texttt{source} dit si elle vient du réseau ou du cache, puis le démarrage, qui fixe la version servie pendant toute la session. Le chargement du référentiel interne laisse une ligne de même nature, avec le nombre de fiches retenues, le nombre de fiches écartées et le prochain identifiant libre --- et son absence même est journalisée plutôt que tue, comme lors de nos campagnes, où le fichier interne était volontairement retiré. Les quatre lignes ci-dessous sont tirées telles quelles du journal : les trois premières au démarrage de la campagne du 11 septembre, la quatrième au démarrage suivant, où la même base revient du cache sous un autre \texttt{run\_id}.
\end{itemize}

\begin{verbatim}
{"ts": "2026-09-11T11:34:25.....", "run_id": "2d2e924e4a44",
 "event": "download_complete", "matrix": "Enterprise",
 "duration_ms": 3745, "release": "19.2"}
{"ts": "2026-09-11T11:34:26.....", "run_id": "2d2e924e4a44",
 "event": "internal_db_missing",
 "path": ".....referentiel_interne.json"}
{"ts": "2026-09-11T11:34:26.....", "run_id": "2d2e924e4a44",
 "event": "server_start", "pid": 15444, "version": "2.2.1",
 "transport": "http", "host": "127.0.0.1", "port": 8733,
 "db_version": {"Enterprise": "19.2", "ICS": "19.2"},
 "ready": true}
{"ts": "2026-09-14T04:35:10.....", "run_id": "7f80330e404b",
 "event": "db_loaded", "matrix": "Enterprise",
 "source": "cache", "release": "19.2", "techniques": 697}
\end{verbatim}

\begin{itemize}
\tightlist
\item
  Ensuite le protocole : \texttt{initialize} consigne la version demandée par le client, celle qui a été négociée et l'identité du client ; \texttt{tools/list} est consigné aussi, parce que c'est par cet appel que le modèle lit les descriptions d'outils sur lesquelles il fonde ses choix ; et un appel adressé à un outil qui n'existe pas produit une ligne dont le statut vaut \texttt{unknown\_tool}, car c'est la signature d'un outil inventé par le modèle et il faut pouvoir la compter.
\item
  Enfin l'appel lui-même. Chaque étape franchie ou refusée écrit une ligne horodatée portant l'identifiant d'exécution de la session. La figure~\ref{fig:ch2-audit} montre la ligne produite par notre appel : l'outil invoqué, le statut, la durée, les arguments tronqués et un résumé du résultat, jamais le résultat entier --- ici deux millisecondes de calcul côté serveur, 4 587 octets rendus, dix candidats sur 72 correspondances, \texttt{T0846.001} en tête. L'écriture est confiée à un fil d'exécution séparé, de sorte qu'elle ne retarde jamais une réponse, et si la file d'attente sature, une ligne publie le nombre d'événements perdus plutôt que de les taire.
\end{itemize}

\begin{figure}[htbp]
\centering
\figuretikz{fig-2-5-ligne-audit}
\caption{Ligne d'audit écrite pour cet appel. En bleu, l'identifiant qui corrèle la session ; en vert, la durée et la version qui datent la réponse ; en ocre, les arguments tronqués ; en rouge, le résumé du résultat.}\label{fig:ch2-audit}
\end{figure}

Le même échange se lit ainsi à trois instants : la requête émise, la réponse datée, la trace d'audit --- trois états successifs du serveur pour une seule question d'analyste.

\subsection{Le moteur de mapping}\label{sec:ch2-2-4}

Le moteur est la partie du serveur qui transforme un signal en identifiants ATT\&CK. Il travaille en deux temps, visibles en figure~\ref{fig:ch2-moteur} : il propose des techniques candidates avec un score, puis il fournit de quoi les vérifier. Il ne tranche jamais à la place de l'analyste, et le rejet est une sortie prévue.

\begin{figure}[htbp]
\centering
\figuretikz{fig-2-6-moteur-mapping}
\caption{Le moteur de mapping : proposer (chaîne du haut), puis vérifier et confronter (chaîne du bas), avec le rejet comme sortie légitime.}\label{fig:ch2-moteur}
\end{figure}

\textbf{Ce que le moteur lit dans le référentiel.} MITRE a changé sa façon de décrire la détection entre ses versions 17.1 et 18. Dans l'ancien modèle, chaque technique portait deux champs de texte : une liste de sources de données et un paragraphe de détection. Dans le nouveau, l'information est publiée sous forme d'objets reliés, sur trois niveaux : une stratégie de détection dit quoi surveiller, un ou plusieurs \emph{analytics} la mettent en œuvre, et chaque \emph{analytic} nomme les journaux à consulter. Sur \texttt{T1003.001}, cela donne :

\begin{verbatim}
DET0363  Detection of Credential Dumping from LSASS Memory
         via Access and Dump Sequence
  Analytic 1030  (Windows)
    journaux    WinEventLog:Sysmon:EventCode=10,
                WinEventLog:Security:EventCode=4673, ...
    composants  Process Access, Process Creation,
                File Creation, ...
\end{verbatim}

Le serveur lit les deux modèles et annonce lequel il a trouvé, dans le champ \texttt{detection\_model}. Sur la version 19.2, ce champ vaut \texttt{strategies} : les 697 techniques Enterprise et les 97 techniques ICS ont toutes au moins une stratégie de détection, et 652 d'un côté, 85 de l'autre portent des journaux exploitables, ce qui alimente un index de 692 libellés de sources. Si l'on place dans le cache le paquet d'une version 17.1, le champ vaut \texttt{hérité} : aucune stratégie n'y est publiée, mais 639 techniques Enterprise et 71 techniques ICS portent quand même leurs sources de données, et \texttt{T1003.001} restitue les 1 119 caractères de son ancien texte de détection. Une étude menée sur une version ancienne reste donc exploitable, et le serveur ne rend jamais une liste de détections vide en laissant croire que MITRE n'a rien publié.

\textbf{Un identifiant écrit de plusieurs façons.} MITRE ne connaît qu'une seule forme pour une sous-technique, à trois décimales. Les rapports, les tickets et les humains en écrivent d'autres. Le serveur les ramène à la forme officielle avant toute consultation, dans tous les outils qui prennent un identifiant, et aussi pour les identifiants repérés à l'intérieur du texte d'une requête :

\begin{verbatim}
get_technique("t1003.1")    -> T1003.001   LSASS Memory
get_technique("T1003.01")   -> T1003.001   LSASS Memory
get_technique("T1003.001")  -> T1003.001   LSASS Memory
\end{verbatim}

\textbf{Les identifiants d'anciennes techniques.} Le référentiel bouge : des techniques sont retirées à chaque version. Un rapport de renseignement écrit il y a cinq ans, une règle de détection ou une fiche interne peuvent encore citer un identifiant qui n'existe plus. Deux situations se présentent, et elles appellent des gestes opposés.

La première : la technique a été remplacée. Le serveur suit la relation de remplacement publiée par MITRE et redirige vers la technique vivante. La version 19.2 compte 158 redirections de ce type, dont aucune n'aboutit à une impasse. MITRE enchaîne parfois plusieurs remplacements successifs, et le chemin complet est conservé, car un ancien document peut citer justement l'identifiant intermédiaire :

\begin{verbatim}
get_technique("T1073")
-> {"redirect": true, "requested": "T1073",
   "replaced_by": "T1574.001", "via": ["T1574.002"],
   "remplacant_vivant": true,
   "note": "'T1073' est revoquee : consulter T1574.001 via
            get_technique (chaine de revocation : T1574.002)"}
\end{verbatim}

La seconde : la technique a été retirée sans remplaçante. Elles sont 24 dans la version 19.2. Le serveur le dit explicitement, au lieu de répondre « non trouvée » comme pour une faute de frappe :

\begin{verbatim}
get_technique("T1064")
-> {"deprecated": true, "requested": "T1064", "name": "Scripting",
   "note": "'T1064' (Scripting) est DEPRECIEE par MITRE sans
            remplacant : ne plus l'utiliser dans un mapping ;
            chercher le comportement avec search_techniques"}
\end{verbatim}

La distinction compte, parce que les deux cas ne se corrigent pas de la même façon : une faute de frappe se réécrit, un identifiant déprécié oblige à remapper le comportement. La recherche signale l'un et l'autre lorsque l'identifiant apparaît dans la requête.

\textbf{La base est lue, jamais écrite.} Aucun des dix-huit outils ne modifie le référentiel MITRE ni le référentiel interne. Le serveur n'écrit que deux choses sur le disque, et aucune n'est un contenu d'analyse : une copie des paquets publiés par MITRE, qui sert de cache, et le journal d'audit. Le cache est écrit par remplacement atomique, si bien qu'une interruption ne laisse jamais un fichier à moitié écrit. Le retéléchargement suit la même règle : les deux matrices sont récupérées et analysées d'abord, et la base en mémoire n'est remplacée qu'en cas de succès complet ; un échec laisse la base précédente en place et le signale dans la réponse. La bascule se fait sous verrou, de sorte qu'aucune requête ne peut voir une base à moitié échangée. Le référentiel interne obéit au même principe : le serveur le relit, mais n'y écrit pas, et la création comme la validation d'une fiche restent des actes humains, faits dans le fichier JSON versionné par l'équipe. Pour le mapping, la conséquence est directe : tout identifiant renvoyé vient de la base chargée, et le modèle n'a pas la possibilité d'en composer un qui n'existe pas.

\textbf{Le score.} La recherche est lexicale : elle compare des mots, pas des sens. Le calcul se fait en quatre étapes. D'abord la requête est découpée en mots et chaque mot est ramené à sa racine : « lsass credential dumping » devient \texttt{lsass}, \texttt{credenti}, \texttt{dump}. La même opération est appliquée au référentiel au chargement, sans quoi \texttt{dumping} dans la requête et \texttt{dump} dans le catalogue ne se rencontreraient jamais. Les mots de moins de trois lettres sont écartés, sauf s'ils contiennent un chiffre : \texttt{c2}, \texttt{2fa} et \texttt{t1003} sont retenus, \texttt{to}, \texttt{of} et \texttt{or} ne le sont pas. La raison est mesurable : ces trois mots figurent respectivement dans 790, 704 et 662 des 794 techniques. Les garder reviendrait à faire peser dans le calcul des mots qui ne distinguent rien, au point qu'une requête de pure grammaire comme « on or of » produirait un candidat. Elle ne produit maintenant aucun résultat, avec la note \texttt{aucun\ terme\ significatif\ dans\ la\ requete}.

Ensuite chaque mot reçoit un poids selon sa rareté dans le référentiel. Mesuré sur la version 19.2 : \texttt{lsass} pèse 5,89 et \texttt{dump} 4,12, tandis que \texttt{credenti} ne pèse que 2,55, \texttt{window} 2,35 et \texttt{file} 2,10. Un mot présent dans des centaines de techniques ne distingue rien.

Chaque technique est alors comparée à la requête sur six emplacements, du plus discriminant au moins discriminant : son nom et les acronymes définis dans sa description comptent pour 1,0 ; le nom de sa technique parente et les entités qui l'emploient --- groupes, logiciels, campagnes --- pour 0,8 ; les noms de ses tactiques pour 0,6 ; sa description pour 0,35.

Enfin chaque mot de la requête ne compte qu'une fois, au meilleur emplacement où il a été trouvé, et le total est rapporté au poids de la requête entière. Le score mesure donc la part de la requête retrouvée, entre 0 et 1. Sur notre exemple, le calcul se pose en entier :

\begin{verbatim}
poids requete :  lsass 5,89 + credenti 2,55 + dump 4,12  =  12,55

T1003.001  LSASS Memory (parente : OS Credential Dumping)
    nom       lsass                    1,00 x 5,89  =  5,89
    parente   credenti, dump           0,80 x 6,67  =  5,34
                                       ---------------------
                                  11,22 / 12,55  =  0,894

T1003      OS Credential Dumping
    nom       credenti, dump           1,00 x 6,67  =  6,67
                                   6,67 / 12,55  =  0,531
\end{verbatim}

Le mot \texttt{lsass} est absent du nom \emph{OS Credential Dumping} : c'est exactement ce qui fait passer la sous-technique devant sa parente ici, et le modèle peut le lire dans la réponse plutôt que de le supposer.

\textbf{Deux plafonds corrigent les cas où cette couverture atteindrait 1 sans rien prouver.} Le premier vise la requête d'un seul mot : « windows » se retrouve dans le nom de dizaines de techniques, et le mot unique de la requête y serait couvert en entier. Le score est donc plafonné à 0,59, juste sous le seuil fort, sauf si la requête est un identifiant vérifiable comme \texttt{T1003.001} ou un nom strictement égal à celui d'une technique, comme « masquerading ». Le second vise la requête réduite au nom d'une entité : « mimikatz » désigne un logiciel, pas un comportement, et les techniques qui remontent sont seulement celles que ce logiciel emploie. Elles sont plafonnées à 0,45 et marquées \texttt{matched.pivot}. Dès qu'une preuve lexicale accompagne l'entité, le plafond tombe : « apt29 credential dumping » rend bien ses candidats à 0,9.

\textbf{Un nom trop court pour être un mot.} Écarter les mots de moins de trois lettres a un effet de bord : la technique \texttt{T1053.002} s'appelle « At », deux caractères, et elle devenait introuvable par son propre nom. Plutôt que d'assouplir la règle et de réintroduire les mots vides, un index des noms exacts est construit au chargement, indépendant du découpage en mots. Il couvre le nom court et le nom complet \texttt{Parente:\ Sous-technique}. Résultat mesuré sur les 794 techniques interrogées par leur propre nom : aucune n'est introuvable, et 96,6 \% sortent en première position.

Deux seuils sont rappelés dans chaque réponse : 0,6 pour un candidat fort, 0,3 en dessous duquel le serveur écrit lui-même que l'absence de mapping est une conclusion valide. À score égal, le départage est fixe et déterministe : nom exactement égal à la requête, puis nom le mieux couvert par la requête, puis technique parente avant sous-technique, puis matrice --- Enterprise d'abord ---, puis nom le plus court, puis identifiant. Deux exécutions de la même requête rendent donc la même liste dans le même ordre, condition pour qu'une analyse soit reproductible.

\begin{verbatim}
search_techniques("valid accounts")
-> T1078      Valid Accounts   1,0  Enterprise   nom égal, matrice IT
  T0859      Valid Accounts   1,0  ICS          nom égal, matrice OT
  T1078.003  Local Accounts   1,0  Enterprise   sous-technique, après
  T1078.004  Cloud Accounts   1,0  Enterprise   sous-technique, après

  "homonymes": [{"name": "Valid Accounts",
      "techniques": [{"id": "T1078", "matrix": "Enterprise"},
                     {"id": "T0859", "matrix": "ICS"}]}]
\end{verbatim}

Le critère de matrice mérite qu'on s'y arrête, parce qu'il porte un risque propre à ce serveur. Vingt-six noms de techniques existent à l'identique dans Enterprise et dans ICS, et « Valid Accounts » en fait partie. À nom identique, le score l'est aussi, et il faut bien trancher d'une façon ou d'une autre : le classement place donc l'informatique de gestion en premier, faute de savoir dans quel monde travaille l'analyste. Le champ \texttt{homonymes} nomme l'ambiguïté quand elle se produit, et une note demande de préciser la matrice. Cela n'en fait qu'un amortisseur : le serveur ne connaît que le paramètre \texttt{matrix} qu'on lui transmet. Sur douze comportements typés, trois des premiers résultats tombaient dans la mauvaise matrice sans ce paramètre, et aucun avec. Le passer systématiquement est la règle.

\textbf{Les preuves.} Chaque candidat est accompagné des mots qui l'ont fait remonter et de l'emplacement où ils ont été trouvés, dans le champ \texttt{matched}. Le modèle peut ainsi écarter le bruit avant même de consulter la fiche complète. Deux exemples mesurés montrent ce que cela apporte :

\begin{verbatim}
search_techniques("lateral movement rdp")
-> T1021.001  Remote Desktop Protocol  0,758
    matched: {"alias": ["rdp"], "tactic": ["later", "movement"]}

search_techniques("mimikatz")
-> T1098  Account Manipulation  0,45
    matched: {"software": ["S0002 Mimikatz"],
              "pivot": "entite seule : piste, pas un mapping"}
\end{verbatim}

Dans le premier cas, \texttt{rdp} n'apparaît pas dans le nom de la technique : il vient de l'acronyme défini dans sa description, « Remote Desktop Protocol (RDP) », et les mots \texttt{lateral\ movement} ont été reconnus comme un nom de tactique. Le second cas montre l'inverse, une preuve qui n'en est pas une : aucun mot de la requête n'est dans le nom de la technique, c'est le logiciel \emph{Mimikatz} qui a été reconnu, et les dix-huit techniques qui remontent sont simplement celles que ce logiciel emploie. Le champ \texttt{pivot} le dit en toutes lettres, le plafond de 0,45 empêche de les déclarer fortes, et une note renvoie à \texttt{get\_software}, qui fait foi pour cette liste. Le champ \texttt{matched} dit donc dans les deux cas d'où vient le score, ce qui permet de juger sa solidité.

\textbf{Ce que le moteur ne fait pas.} Il ne comprend pas le sens des mots. Une requête en français ne trouve rien d'utile, et la réponse le dit :

\begin{verbatim}
search_techniques("imprimante thermique defectueuse")
-> 0 candidat
  note: aucun candidat : l'absence de mapping est une conclusion
        valide

search_techniques("vol d'identifiants dans la memoire")
-> T1591.004  Identify Roles  0,333
  note: aucun candidat MITRE fort : confronter le comportement au
        referentiel interne avec search_internal
\end{verbatim}

Le second cas est le plus instructif : la requête décrit exactement le comportement recherché, et pourtant le meilleur candidat plafonne à 0,333 --- et ce n'est même pas le bon, puisque c'est le mot \texttt{identifiants} qui a été rapproché de \texttt{Identify}. Aucun des mots français employés n'existe dans le catalogue. Les requêtes doivent donc être formulées dans le vocabulaire anglais d'ATT\&CK, un comportement à la fois. C'est la règle que le modèle lit dans la description de l'outil avant d'appeler, et c'est aussi pourquoi la vérification par \texttt{get\_technique} reste obligatoire : un score élevé prouve que les mots correspondent, pas que le comportement est le bon.

\section{Évaluation de l'approche}\label{sec:ch2-3}

La méthodologie d'évaluation est celle du chapitre I, sans modification : même corpus, même étiquetage des prédictions, mêmes deux bases de calcul, mêmes formules de précision, de rappel et de F1, mêmes attributs de diagnostic et mêmes dénominateurs. Elle n'est pas reprise ici ; seuls les changements le sont. Un seul modèle est évalué, Claude Opus 4.8, retenu au terme de la sélection : la question n'est plus lequel choisir, mais ce que gagne celui qui a été choisi.

Une seconde question s'ajoute, qui n'a de sens qu'avec un serveur. Les métriques du chapitre I demandent \emph{le mapping est-il juste ?} ; les indicateurs ajoutés ici demandent \emph{la réponse vient-elle du référentiel ou de la mémoire du modèle ?} Un F1 qui progresse sans que la seconde soit posée ne prouverait rien : le modèle pourrait avoir mieux répondu sans jamais consulter la base.

\subsection{La ligne à battre}\label{sec:ch2-3-1}

Le tableau~\ref{tab:ch2-ligne-base} rassemble les résultats de Claude Opus 4.8 sans serveur, sur les 42 échantillons et les 126 exécutions de la campagne du chapitre I. C'est le point de départ, et la seule référence contre laquelle la phase 2 se mesure.

\begin{longtable}[]{@{}
  >{\raggedright\arraybackslash}p{(\columnwidth - 4\tabcolsep) * \real{0.4018}}
  >{\raggedright\arraybackslash}p{(\columnwidth - 4\tabcolsep) * \real{0.2991}}
  >{\raggedright\arraybackslash}p{(\columnwidth - 4\tabcolsep) * \real{0.2991}}@{}}
\toprule\noalign{}
\begin{minipage}[b]{\linewidth}\raggedright
\end{minipage} & \begin{minipage}[b]{\linewidth}\raggedright
\textbf{Base stricte}
\end{minipage} & \begin{minipage}[b]{\linewidth}\raggedright
\textbf{Base tolérante}
\end{minipage} \\
\midrule\noalign{}
\endhead
\bottomrule\noalign{}
\endlastfoot
F1 technique & 0,563 & 0,571 \\
F1 tactique & 0,587 & 0,587 \\
F1 sous-technique & 0,636 & 0,682 \\
F1 technique --- Enterprise / ICS & 0,714 / 0,262 & --- \\
Hors matrice · hallucination & 0 · 0,0 \% & 0 · 0,0 \% \\
Obsolescence · abstention & 1,3 \% · 0,0 \% & 1,3 \% · 0,0 \% \\
\end{longtable}
\legendetableau{Claude Opus 4.8 sans serveur : la ligne à battre.}{tab:ch2-ligne-base}


Hors qualité du mapping, la même campagne relève une stabilité de 76,2 \%, une latence médiane de 4 282 ms pour un 95e centile de 6 050 ms, 6 572 jetons par exécution et une efficacité de 0,086 --- le F1 technique obtenu par millier de jetons.

\textbf{C'est la colonne tolérante qui fait référence.} Les deux bases portent sur les mêmes réponses et ne diffèrent que sur un point : la base stricte compte comme une erreur un identifiant que MITRE a retiré depuis, la base tolérante le crédite, puisque la table de correspondance publiée par MITRE atteste qu'il désignait bien le comportement attendu. Or résoudre ces identifiants retirés est exactement ce que le serveur mécanise. Comparer la phase 2 à la seule base stricte mélangerait donc deux effets : le rattrapage de version, que le serveur fait à coup sûr, et l'apport propre de la consultation. Le seuil à dépasser est donc \textbf{0,571 en technique}, \textbf{0,587 en tactique} et \textbf{0,682 en sous-technique}.

\textbf{Deux chiffres ne peuvent pas s'améliorer.} Ce modèle ne produit, sans serveur, aucun identifiant inventé ni aucun identifiant pris dans l'autre matrice : ces deux compteurs sont à zéro et ne peuvent donc que se dégrader. Leur dégradation serait un signal fort, car le serveur ne rend que des identifiants lus dans la base chargée --- un identifiant inexistant apparaissant en phase 2 ne pourrait venir que de la mémoire du modèle, contre l'avis de l'outil qu'il vient d'interroger.

\textbf{La marge se lit sur la matrice industrielle.} Ventilé par matrice, le F1 technique passe de \textbf{0,714 sur Enterprise} à \textbf{0,262 sur ICS}. L'écart n'est pas un accident de mesure : les comportements industriels sont moins documentés, moins présents dans les données d'entraînement, et le référentiel ICS a évolué plus récemment. C'est donc là qu'un accès à un référentiel à jour a le plus à apporter, et c'est sur cette ventilation que le gain se lira en premier.

\subsection{Ce que la phase 2 ajoute au protocole}\label{sec:ch2-3-2}

\textbf{Le workflow : un seul changement.} Le harnais est celui du chapitre 1, à l'identique --- mêmes 42 échantillons, trois exécutions, format, statuts, relances et exclusions. Une seule chose change : entre le modèle et sa réponse s'intercale un client MCP connecté au serveur de la section~\ref{sec:ch2-2}. Au démarrage, le workflow sonde le serveur, découvre ses outils et les traduit au format de chaque fournisseur ; la version du référentiel chargée est inscrite dans chaque enregistrement. L'appel devient une boucle, plafonnée à huit tours : à chaque tour, le modèle rend son triplet ou demande des outils, dont les résultats sont réinjectés dans le contexte ; épuiser les tours sans conclure vaut une abstention. Le harnais enregistre en plus les tours, les appels par outil et la version de base ; la latence couvre la boucle entière, appels d'outils compris --- consulter le serveur est un coût de l'approche. La figure~\ref{fig:workflow-mcp} situe cet ajout.

\begin{figure}[htbp]
\centering
\figuretikz{fig-2-3-workflow-mcp}
\caption{Le workflow de la phase 2 : le harnais du chapitre 1, augmenté du client MCP et de sa boucle d'outils (huit tours au plus).}\label{fig:workflow-mcp}
\end{figure}

Une exécution n'est plus un appel et une réponse, mais une conversation : le modèle reçoit les journaux, demande un ou plusieurs outils, reçoit leurs résultats, et recommence jusqu'à pouvoir conclure. Quatre décisions en découlent.

\begin{itemize}
\tightlist
\item
  \textbf{L'appel d'outil est imposé au premier tour}, les tours suivants rendant la main. Sans cette contrainte, le modèle pourrait répondre de mémoire et l'exécution ne mesurerait pas l'ancrage mais son absence. \textbf{Conséquence à déclarer :} le taux d'utilisation du serveur vaut 100 \% par construction et n'apprend donc rien. Ce sont le taux d'appels pertinents et le taux de réponses ancrées qui portent l'information.
\item
  \textbf{Le modèle choisit ses outils}. Le serveur en publie dix-huit, et le modèle les découvre lui-même par \emph{tools/list} : décider lequel appeler, et dans quel ordre, fait partie de ce que la campagne mesure. Restreindre le périmètre aux seuls outils utiles au mapping rendrait le taux d'appels pertinents égal à 100 \% par construction, comme le taux d'utilisation. Une seule exclusion, qui n'est pas un choix de périmètre mais un garde-fou : \emph{mitre\_update} remplace la base en mémoire et changerait donc la version ATT\&CK servie au milieu de la campagne. Les dix-sept autres sont en lecture seule.
\item
  \textbf{La consigne système gagne un alinéa} : consulter le référentiel plutôt que répondre de mémoire, préciser la matrice, vérifier chaque candidat retenu, ne rien changer au format de sortie. Le reste de la consigne est recopié sans modification, et les deux morceaux sont enregistrés séparément pour que la différence soit vérifiable. Coût méthodologique assumé : la comparaison mesure l'effet conjoint de l'accès au référentiel et de la consigne qui invite à s'en servir ; les séparer demanderait une troisième campagne, consigne retirée.
\item
  \textbf{Le budget de tours du harnais}, fixé à huit allers-retours, s'ajoute aux cas déjà écartés du calcul. Une exécution qui n'y converge pas est close et écartée, au même titre qu'une réponse tronquée faute de jetons : ce sont deux limites du banc d'essai, pas deux échecs du modèle. L'effectif écarté est rapporté.
\end{itemize}

Le chapitre I compte en échantillons et en exécutions. S'y ajoute \textbf{l'appel d'outil} : une exécution en produit plusieurs, et compter en exécutions un comportement qui se répète dans la même conversation le sous-estimerait. Un indicateur portant sur la façon d'appeler se compte donc en appels, un indicateur portant sur la conduite d'une exécution se compte en exécutions.

\begin{verbatim}
E  exécutions retenues      C  appels d’outils émis
T  tours de conversation    R  appels de recherche (R <= C)
\end{verbatim}

\subsection{Les indicateurs ajoutés}\label{sec:ch2-3-3}

Onze indicateurs sont calculés depuis la trace que le harnais écrit pour chaque appel : l'outil invoqué, ses arguments, son issue, et les identifiants que le serveur a rendus. Aucun ne modifie un F1 ni n'entre dans un classement --- ce sont des attributs de diagnostic, qui expliquent \emph{comment} le modèle s'est servi du serveur là où les F1 mesurent ce qu'il en a tiré.

\begin{longtable}[]{@{}
  >{\raggedright\arraybackslash}p{(\columnwidth - 4\tabcolsep) * \real{0.2889}}
  >{\raggedright\arraybackslash}p{(\columnwidth - 4\tabcolsep) * \real{0.5305}}
  >{\raggedright\arraybackslash}p{(\columnwidth - 4\tabcolsep) * \real{0.1806}}@{}}
\toprule\noalign{}
\begin{minipage}[b]{\linewidth}\raggedright
\textbf{Indicateur}
\end{minipage} & \begin{minipage}[b]{\linewidth}\raggedright
\textbf{Formule}
\end{minipage} & \begin{minipage}[b]{\linewidth}\raggedright
\textbf{Unité}
\end{minipage} \\
\midrule\noalign{}
\endhead
\bottomrule\noalign{}
\endlastfoot
Taux d'utilisation & exécutions avec $\geq$ 1 appel / E & exécution \\
Appels par exécution & C / E & appel \\
Tours par exécution & T / E & tour \\
Taux d'appels pertinents & appels vers un outil rendant des identifiants / C & appel \\
Taux d'hallucination d'outil & appels vers un outil inexistant / C & appel \\
Appels en erreur & effectif brut --- arguments refusés & appel \\
Démentis du serveur & effectif brut --- identifiant déclaré introuvable & appel \\
Conformité de matrice & recherches portant la bonne matrice / R & appel \\
Exécutions conformes & exéc. sans recherche fautive / exéc. avec $\geq$ 1 recherche & exécution \\
Averti puis faux & effectif brut --- homonymie signalée, identifiant hors matrice retenu & exécution \\
Taux de réponses ancrées & exéc. dont l'identifiant retenu a été rendu par un outil / exéc. ayant répondu & exécution \\
\end{longtable}
\legendetableau{Les indicateurs ajoutés, leur formule et leur unité de compte.}{tab:ch2-indicateurs}


\textbf{Le taux d'appels pertinents} vise les quatre outils dont la sortie porte des identifiants de tactique ou de technique. Les treize autres --- mitigations, groupes, logiciels, campagnes, sources de journaux --- ne sont pas fautifs : ils ne servent simplement pas cette tâche. L'indicateur mesure donc la pertinence de la \emph{sélection} d'outil, ce qui suppose que le modèle ait eu le choix.

\textbf{Le taux de réponses ancrées est le plus important des onze.} Une réponse est dite ancrée lorsque l'identifiant finalement retenu figure parmi ceux que les outils ont effectivement rendus au cours de l'exécution. Dans le cas contraire, le modèle l'a produit de mémoire : pas nécessairement faux, mais non vérifié --- or la vérification est précisément ce que le serveur devait apporter. Croisé avec l'étiquetage du chapitre I, cet indicateur sépare quatre situations que le seul F1 confondrait : ancrée et juste, l'ancrage a fonctionné ; ancrée et fausse, c'est le moteur de recherche qui a proposé le mauvais candidat ; non ancrée et juste, le modèle savait déjà et le serveur n'a rien apporté sur cet échantillon ; non ancrée et fausse, c'est exactement le cas que l'ancrage devait empêcher.

\textbf{La conformité de matrice tranche entre deux causes d'un même symptôme.} Le serveur ignore dans quel monde travaille l'analyste : il ne connaît que le paramètre \emph{matrix} que le modèle lui transmet, et sans ce paramètre une recherche interroge les deux matrices à la fois. Comme vingt-six noms de techniques existent à l'identique dans Enterprise et dans ICS, un identifiant industriel peut alors être rendu pour un comportement bureautique, ou l'inverse. Un identifiant hors matrice subsistant en phase 2 admet donc deux explications opposées : ou bien le modèle n'a pas consulté le serveur, ou bien il l'a consulté sans lui dire dans quelle matrice chercher. La première se corrige par la consigne, la seconde par le harnais ; les confondre conduirait à corriger la mauvaise. Cette conformité est mesurée et non forcée : le harnais pourrait injecter la matrice dans chaque appel, mais tenir compte d'une contrainte énoncée dans sa consigne fait partie de ce que la campagne mesure.

\textbf{Deux indicateurs existants changent de contenu, non de formule.} Les \textbf{jetons} portent sur la somme de tous les tours d'une exécution, chacun rechargeant la consigne, les descriptions d'outils et les résultats déjà reçus. La \textbf{latence} englobe les allers-retours ; la part imputable au serveur est isolée, et se compte en millisecondes quand l'exécution se compte en secondes --- le surcoût de l'ancrage n'est donc pas le temps de réponse du référentiel, c'est le temps que le modèle passe à relire ce qu'il en a reçu.

\subsection{Ce que la phase 2 doit montrer}\label{sec:ch2-3-4}

Les attentes sont posées avant les résultats, pour qu'elles ne puissent pas être ajustées après coup.

\begin{itemize}
\tightlist
\item
  \textbf{Le taux d'obsolescence doit s'effondrer.} ATT\&CK évolue à chaque version : des techniques sont retirées et remplacées par d'autres, et MITRE publie alors une relation de remplacement qui désigne l'identifiant vivant prenant la suite. La version 19.2 en compte 158. Un modèle dont les données d'entraînement sont antérieures à l'une de ces révisions continue de produire l'ancien identifiant : son raisonnement est juste, son numéro est périmé --- c'est ce que mesurent les 1,3 \% du tableau~\ref{tab:ch2-ligne-base}. Le serveur, lui, suit ces relations et ne rend jamais qu'un identifiant vivant. Un modèle qui consulte avant de répondre ne devrait donc plus en produire du tout. Si le taux ne baisse pas, c'est que la réponse n'est pas passée par le serveur, et le taux de réponses ancrées le confirmera sur l'exécution en cause.
\item
  \textbf{Le gain doit se voir sur ICS avant tout.} L'écart de 0,714 à 0,262 entre les deux matrices désigne le terrain où la connaissance propre du modèle est la plus lacunaire. Un gain général et uniforme serait d'ailleurs suspect : il indiquerait un effet de consigne plutôt qu'un effet d'ancrage.
\item
  \textbf{Les deux compteurs à zéro doivent le rester} --- aucun identifiant inventé, aucun identifiant hors matrice.
\item
  \textbf{Le coût doit être chiffré, pas éludé.} La conversation multiplie les jetons et allonge la latence. L'efficacité --- le F1 technique obtenu par millier de jetons, 0,086 sans serveur --- dira si ce surcoût achète de la qualité ou seulement du temps.
\end{itemize}

\textbf{Le résultat qui invaliderait la lecture} est un taux de réponses ancrées bas. Si le modèle appelle les outils parce que le harnais l'y oblige, puis répond de mémoire, les F1 mesureront sa connaissance du référentiel et non l'apport du serveur, et l'écart entre les deux phases perdrait son sens. C'est pourquoi cet indicateur est rapporté avant les F1, et non après.

\subsection{Résultats}\label{sec:ch2-3-5}

42 échantillons, 3 exécutions, 126 exécutions par condition. Aucune itération écartée : ni panne d'infrastructure, ni budget de sortie épuisé, ni boucle d'outils non convergente. Le serveur a servi la version 19.2 des deux matrices sur les 126 exécutions, sans dérive.

\begin{longtable}[]{@{}
  >{\raggedright\arraybackslash}p{(\columnwidth - 16\tabcolsep) * \real{0.2235}}
  >{\raggedright\arraybackslash}p{(\columnwidth - 16\tabcolsep) * \real{0.0880}}
  >{\raggedright\arraybackslash}p{(\columnwidth - 16\tabcolsep) * \real{0.0948}}
  >{\raggedright\arraybackslash}p{(\columnwidth - 16\tabcolsep) * \real{0.0948}}
  >{\raggedright\arraybackslash}p{(\columnwidth - 16\tabcolsep) * \real{0.0948}}
  >{\raggedright\arraybackslash}p{(\columnwidth - 16\tabcolsep) * \real{0.0880}}
  >{\raggedright\arraybackslash}p{(\columnwidth - 16\tabcolsep) * \real{0.0880}}
  >{\raggedright\arraybackslash}p{(\columnwidth - 16\tabcolsep) * \real{0.0880}}
  >{\raggedright\arraybackslash}p{(\columnwidth - 16\tabcolsep) * \real{0.1400}}@{}}
\toprule\noalign{}
\begin{minipage}[b]{\linewidth}\raggedright
\textbf{Niveau}
\end{minipage} & \begin{minipage}[b]{\linewidth}\raggedright
\textbf{MCP}
\end{minipage} & \begin{minipage}[b]{\linewidth}\raggedright
\textbf{N}
\end{minipage} & \begin{minipage}[b]{\linewidth}\raggedright
\textbf{VP}
\end{minipage} & \begin{minipage}[b]{\linewidth}\raggedright
\textbf{FP}
\end{minipage} & \begin{minipage}[b]{\linewidth}\raggedright
\textbf{H}
\end{minipage} & \begin{minipage}[b]{\linewidth}\raggedright
\textbf{O}
\end{minipage} & \begin{minipage}[b]{\linewidth}\raggedright
\textbf{A}
\end{minipage} & \begin{minipage}[b]{\linewidth}\raggedright
\textbf{MM}
\end{minipage} \\
\midrule\noalign{}
\endhead
\bottomrule\noalign{}
\endlastfoot
Tactique & non & 126 & 74 & 52 & 0 & 0 & 0 & 0 \\
Technique & non & 126 & 71 & 54 & 0 & 1 & 0 & 0 \\
Sous-technique & non & 66 & 42 & 21 & 0 & 3 & 0 & 0 \\
Tactique & oui & 126 & 101 & 25 & 0 & 0 & 0 & 0 \\
Technique & oui & 126 & 89 & 37 & 0 & 0 & 0 & 0 \\
Sous-technique & oui & 66 & 50 & 16 & 0 & 0 & 0 & 0 \\
\end{longtable}
\legendetableau{Étiquetage des prédictions, sans puis avec serveur. N = VP + FP + H + O + A, vérifié sur les six lignes.}{tab:ch2-etiquetage}


\begin{longtable}[]{@{}
  >{\raggedright\arraybackslash}p{(\columnwidth - 18\tabcolsep) * \real{0.1016}}
  >{\raggedright\arraybackslash}p{(\columnwidth - 18\tabcolsep) * \real{0.0790}}
  >{\raggedright\arraybackslash}p{(\columnwidth - 18\tabcolsep) * \real{0.1016}}
  >{\raggedright\arraybackslash}p{(\columnwidth - 18\tabcolsep) * \real{0.1072}}
  >{\raggedright\arraybackslash}p{(\columnwidth - 18\tabcolsep) * \real{0.1072}}
  >{\raggedright\arraybackslash}p{(\columnwidth - 18\tabcolsep) * \real{0.1072}}
  >{\raggedright\arraybackslash}p{(\columnwidth - 18\tabcolsep) * \real{0.0993}}
  >{\raggedright\arraybackslash}p{(\columnwidth - 18\tabcolsep) * \real{0.0959}}
  >{\raggedright\arraybackslash}p{(\columnwidth - 18\tabcolsep) * \real{0.1005}}
  >{\raggedright\arraybackslash}p{(\columnwidth - 18\tabcolsep) * \real{0.1005}}@{}}
\toprule\noalign{}
\begin{minipage}[b]{\linewidth}\raggedright
\textbf{MCP}
\end{minipage} & \begin{minipage}[b]{\linewidth}\raggedright
\textbf{MM}
\end{minipage} & \begin{minipage}[b]{\linewidth}\raggedright
\textbf{Hall.}
\end{minipage} & \begin{minipage}[b]{\linewidth}\raggedright
\textbf{F1 tech.}
\end{minipage} & \begin{minipage}[b]{\linewidth}\raggedright
\textbf{F1 tact.}
\end{minipage} & \begin{minipage}[b]{\linewidth}\raggedright
\textbf{F1 s.-t.}
\end{minipage} & \begin{minipage}[b]{\linewidth}\raggedright
\textbf{Ent.}
\end{minipage} & \begin{minipage}[b]{\linewidth}\raggedright
\textbf{ICS}
\end{minipage} & \begin{minipage}[b]{\linewidth}\raggedright
\textbf{Obsol.}
\end{minipage} & \begin{minipage}[b]{\linewidth}\raggedright
\textbf{Abst.}
\end{minipage} \\
\midrule\noalign{}
\endhead
\bottomrule\noalign{}
\endlastfoot
non & 0 & 0,0 \% & 0,563 & 0,587 & 0,636 & 0,714 & 0,262 & 1,3 \% & 0,0 \% \\
oui & 0 & 0,0 \% & 0,706 & 0,802 & 0,758 & 0,798 & 0,524 & 0,0 \% & 0,0 \% \\
\end{longtable}
\legendetableau{Qualité du mapping, base stricte. Ent. et ICS : F1 technique ventilé par matrice.}{tab:ch2-strict}


\begin{longtable}[]{@{}
  >{\raggedright\arraybackslash}p{(\columnwidth - 16\tabcolsep) * \real{0.1016}}
  >{\raggedright\arraybackslash}p{(\columnwidth - 16\tabcolsep) * \real{0.0790}}
  >{\raggedright\arraybackslash}p{(\columnwidth - 16\tabcolsep) * \real{0.1072}}
  >{\raggedright\arraybackslash}p{(\columnwidth - 16\tabcolsep) * \real{0.1185}}
  >{\raggedright\arraybackslash}p{(\columnwidth - 16\tabcolsep) * \real{0.1185}}
  >{\raggedright\arraybackslash}p{(\columnwidth - 16\tabcolsep) * \real{0.1185}}
  >{\raggedright\arraybackslash}p{(\columnwidth - 16\tabcolsep) * \real{0.1298}}
  >{\raggedright\arraybackslash}p{(\columnwidth - 16\tabcolsep) * \real{0.1129}}
  >{\raggedright\arraybackslash}p{(\columnwidth - 16\tabcolsep) * \real{0.1140}}@{}}
\toprule\noalign{}
\begin{minipage}[b]{\linewidth}\raggedright
\textbf{MCP}
\end{minipage} & \begin{minipage}[b]{\linewidth}\raggedright
\textbf{MM}
\end{minipage} & \begin{minipage}[b]{\linewidth}\raggedright
\textbf{Hall.}
\end{minipage} & \begin{minipage}[b]{\linewidth}\raggedright
\textbf{F1 tech.}
\end{minipage} & \begin{minipage}[b]{\linewidth}\raggedright
\textbf{F1 tact.}
\end{minipage} & \begin{minipage}[b]{\linewidth}\raggedright
\textbf{F1 s.-t.}
\end{minipage} & \begin{minipage}[b]{\linewidth}\raggedright
\textbf{Écart obs.}
\end{minipage} & \begin{minipage}[b]{\linewidth}\raggedright
\textbf{Obsol.}
\end{minipage} & \begin{minipage}[b]{\linewidth}\raggedright
\textbf{Abst.}
\end{minipage} \\
\midrule\noalign{}
\endhead
\bottomrule\noalign{}
\endlastfoot
non & 0 & 0,0 \% & 0,571 & 0,587 & 0,682 & +0,008 & 1,3 \% & 0,0 \% \\
oui & 0 & 0,0 \% & 0,706 & 0,802 & 0,758 & 0,000 & 0,0 \% & 0,0 \% \\
\end{longtable}
\legendetableau{Qualité du mapping, base tolérante. Écart obs. = F1 tolérant - F1 strict, au niveau technique.}{tab:ch2-tolerant}


La référence à battre était 0,571 en technique, 0,587 en tactique et 0,682 en sous-technique. Les trois sont dépassées : \textbf{+0,135 en technique}, \textbf{+0,215 en tactique}, \textbf{+0,076 en sous-technique}. Le taux d'obsolescence tombe à zéro et l'écart entre les deux bases disparaît : il n'y a plus d'identifiant retiré à créditer. Aucun identifiant fabriqué ni hors matrice n'apparaît, de part et d'autre. Le gain est porté par la matrice industrielle, qui double presque --- 0,262 à 0,524 ---, tandis qu'Enterprise progresse de 0,084 : le gain n'est ni général ni uniforme, ce qui écarte l'hypothèse d'un simple effet de consigne. Deux diagnostics achèvent la lecture : les 9 confusions de la scission v19 disparaissent complètement, et les sous-techniques niées à tort passent de 9 à 6.

\textbf{L'obsolescence ne diminue pas : elle disparaît.} Le taux d'obsolescence ne baisse pas : il disparaît. Les quatre identifiants retirés que Claude Opus 4.8 produisait sans serveur provenaient de deux échantillons seulement, et aucun des deux n'en produit plus. Sur S027 --- effacement des journaux système Linux ---, le modèle répondait T1070.002 aux trois exécutions, un identifiant révoqué au profit de T1685.006 qui porte exactement le même nom ; avec le serveur, les trois exécutions rendent T1685.006. Sur S007 --- message Modbus non autorisé ---, il répondait T0836 deux fois et T0855 une fois, ce dernier révoqué au profit de T1692.001 ; avec le serveur, les trois exécutions rendent T1692.001. Le troisième échantillon du corpus exposé à une révocation, S038, était déjà juste sans serveur et le reste.

\begin{longtable}[]{@{}
  >{\raggedright\arraybackslash}p{(\columnwidth - 6\tabcolsep) * \real{0.1261}}
  >{\raggedright\arraybackslash}p{(\columnwidth - 6\tabcolsep) * \real{0.1667}}
  >{\raggedright\arraybackslash}p{(\columnwidth - 6\tabcolsep) * \real{0.3536}}
  >{\raggedright\arraybackslash}p{(\columnwidth - 6\tabcolsep) * \real{0.3536}}@{}}
\toprule\noalign{}
\begin{minipage}[b]{\linewidth}\raggedright
\textbf{Éch.}
\end{minipage} & \begin{minipage}[b]{\linewidth}\raggedright
\textbf{Matrice}
\end{minipage} & \begin{minipage}[b]{\linewidth}\raggedright
\textbf{Sans serveur}
\end{minipage} & \begin{minipage}[b]{\linewidth}\raggedright
\textbf{Avec serveur}
\end{minipage} \\
\midrule\noalign{}
\endhead
\bottomrule\noalign{}
\endlastfoot
S027 & Enterprise & T1070 / T1070.002 (trois exécutions) & T1685 / T1685.006 (trois exécutions) \\
S007 & ICS & T0836 $\times$2, puis T0855 (une exécution) & T1692 / T1692.001 (trois exécutions) \\
S038 & Enterprise & T1547 / T1547.001 (déjà juste) & T1547 / T1547.001 (inchangé) \\
\end{longtable}
\legendetableau{Les échantillons du corpus exposés à une révocation, avant et après ancrage.}{tab:ch2-revocation}


\textbf{Ce que ces deux cas démontrent.} Ces deux corrections ne relèvent pas d'un meilleur raisonnement. Sur S027, le modèle avait déjà identifié la bonne action --- effacer des journaux --- et nommait la bonne technique sous son ancien numéro ; il lui manquait de savoir que ce numéro avait changé. Sur S007, le cas est plus net encore : deux exécutions sur trois donnaient T0836, une technique vivante mais fausse, que la base tolérante ne rattrape pas. L'ancrage corrige donc deux défauts distincts avec le même mécanisme --- un référentiel périmé et une confusion ordinaire --- parce que dans les deux cas le modèle a consulté avant de répondre au lieu de puiser dans sa mémoire.

\textbf{Portée du résultat.} C'est le résultat le plus net de la campagne, et le plus directement attribuable au serveur. Un identifiant retiré est un défaut que le modèle ne peut pas corriger seul : rien dans le journal ne lui signale que le référentiel a bougé depuis son entraînement. Le serveur, lui, suit les relations de remplacement publiées par MITRE et ne rend jamais qu'un identifiant vivant. Le passage de 1,3 \% à 0 \% n'est donc pas une amélioration graduelle mais la suppression d'une classe entière d'erreurs.

\begin{longtable}[]{@{}
  >{\raggedright\arraybackslash}p{(\columnwidth - 14\tabcolsep) * \real{0.1129}}
  >{\raggedright\arraybackslash}p{(\columnwidth - 14\tabcolsep) * \real{0.1275}}
  >{\raggedright\arraybackslash}p{(\columnwidth - 14\tabcolsep) * \real{0.1332}}
  >{\raggedright\arraybackslash}p{(\columnwidth - 14\tabcolsep) * \real{0.1163}}
  >{\raggedright\arraybackslash}p{(\columnwidth - 14\tabcolsep) * \real{0.1332}}
  >{\raggedright\arraybackslash}p{(\columnwidth - 14\tabcolsep) * \real{0.1174}}
  >{\raggedright\arraybackslash}p{(\columnwidth - 14\tabcolsep) * \real{0.1185}}
  >{\raggedright\arraybackslash}p{(\columnwidth - 14\tabcolsep) * \real{0.1411}}@{}}
\toprule\noalign{}
\begin{minipage}[b]{\linewidth}\raggedright
\textbf{Util.}
\end{minipage} & \begin{minipage}[b]{\linewidth}\raggedright
\textbf{App./ex.}
\end{minipage} & \begin{minipage}[b]{\linewidth}\raggedright
\textbf{Tours/ex.}
\end{minipage} & \begin{minipage}[b]{\linewidth}\raggedright
\textbf{Pertin.}
\end{minipage} & \begin{minipage}[b]{\linewidth}\raggedright
\textbf{Hall. outil}
\end{minipage} & \begin{minipage}[b]{\linewidth}\raggedright
\textbf{Ancrées}
\end{minipage} & \begin{minipage}[b]{\linewidth}\raggedright
\textbf{Matrice}
\end{minipage} & \begin{minipage}[b]{\linewidth}\raggedright
\textbf{Démentis}
\end{minipage} \\
\midrule\noalign{}
\endhead
\bottomrule\noalign{}
\endlastfoot
100 \% & 2,9 & 3,4 & 100 \% & 0,0 \% & 100 \% & 100 \% & 0 \\
\end{longtable}
\legendetableau{Usage du serveur, campagne avec MCP.}{tab:ch2-usage}


\textbf{Le serveur a fait sa part.} Les trois F1 montent, les identifiants périmés disparaissent, rien n'est inventé, rien n'est pris hors de la matrice, et chaque réponse s'appuie sur ce que les outils ont rendu. Reste une question : pourquoi une approche déterministe --- un référentiel figé, interrogé en direct --- ne donne-t-elle pas un score plus élevé ? Le tableau de qualité ne peut pas y répondre : il note la réponse finale, sans dire ce que le modèle avait sous les yeux au moment de la donner.

\textbf{Une analyse a donc été ajoutée au script.} Le serveur ne rend pas une réponse, il rend une liste de candidats, et c'est le modèle qui choisit dedans. Le script fait maintenant deux vérifications par exécution, et à chacun des trois niveaux. Premièrement, il rassemble tous les identifiants que les outils ont rendus pendant l'exécution et regarde si la bonne réponse s'y trouve : si oui, le serveur l'a \emph{fournie}. Deuxièmement, il regarde la réponse finale du modèle : si elle est égale à la bonne réponse, le modèle l'a \emph{retenue}. Le croisement des deux donne quatre situations : fournie et retenue ; fournie mais le modèle en retient une autre ; jamais fournie et le modèle est juste quand même ; jamais fournie et le modèle se trompe. Ces compteurs sont écrits dans le rapport, et le fichier CSV porte, pour chaque exécution et chaque niveau, une colonne qui dit si le serveur avait fourni la bonne réponse.

Le tableau~\ref{tab:ch2-serveur-modele} rapporte ces deux mesures côte à côte. Elles ne se lisent pas comme le tableau de qualité : aucune précision n'est calculée pour le serveur, puisqu'il propose plusieurs candidats et non une réponse unique. La ligne du serveur est le plafond de la condition --- le modèle ne peut pas faire mieux que ce qui lui a été fourni ; la ligne du modèle est son score réel ; l'écart entre les deux est la part des exécutions où la bonne réponse était disponible et n'a pas été retenue.

\begin{longtable}[]{@{}
  >{\raggedright\arraybackslash}p{(\columnwidth - 6\tabcolsep) * \real{0.3400}}
  >{\raggedright\arraybackslash}p{(\columnwidth - 6\tabcolsep) * \real{0.2200}}
  >{\raggedright\arraybackslash}p{(\columnwidth - 6\tabcolsep) * \real{0.2200}}
  >{\raggedright\arraybackslash}p{(\columnwidth - 6\tabcolsep) * \real{0.2200}}@{}}
\toprule\noalign{}
\begin{minipage}[b]{\linewidth}\raggedright
\textbf{}
\end{minipage} & \begin{minipage}[b]{\linewidth}\raggedright
\textbf{Tactique}
\end{minipage} & \begin{minipage}[b]{\linewidth}\raggedright
\textbf{Technique}
\end{minipage} & \begin{minipage}[b]{\linewidth}\raggedright
\textbf{Sous-technique}
\end{minipage} \\
\midrule\noalign{}
\endhead
\bottomrule\noalign{}
\endlastfoot
Serveur MCP --- réponse fournie & 96,8 \% & 89,7 \% & 90,9 \% \\
Modèle Opus --- réponse retenue & 80,2 \% & 70,6 \% & 75,8 \% \\
Perte de sélection & 16,7 pts & 19,0 pts & 15,2 pts \\
Jamais fournie, modèle faux & 3,2 \% & 10,3 \% & 9,1 \% \\
Exécutions comptées & 126 & 126 & 66 \\
\end{longtable}
\legendetableau{Ce que le serveur a fourni, ce que le modèle a retenu, aux trois niveaux, en pourcentage des exécutions. Ligne 1 : la bonne réponse figurait parmi les identifiants rendus par les outils. Ligne 2 : la réponse finale était la bonne. Ligne 3 : l'écart entre les deux, en points. Ligne 4 : le serveur n'avait pas fourni la bonne réponse et le modèle s'est trompé.}{tab:ch2-serveur-modele}

\textbf{Ce qui a réussi.} Quatre choses, toutes comptées exécution par exécution. Les identifiants périmés ont disparu, parce que le serveur ne rend que des identifiants vivants. Rien n'a été inventé, et rien n'a été pris dans la mauvaise matrice. Chaque réponse s'appuie sur un identifiant qu'un outil a rendu pendant l'exécution : aucune n'est venue de la mémoire du modèle. Et la recherche a bien fonctionné : au niveau technique, la bonne réponse figurait dans ce que le serveur a rendu pour 89,7~\% des exécutions.

\textbf{Ce qui reste en défaut.} 37 exécutions sur 126 rendent une mauvaise technique. Le tableau~\ref{tab:ch2-serveur-modele} les coupe en deux : dans 24, la bonne réponse était dans la liste et le modèle en a retenu une autre ; dans 13, aucun outil ne l'a jamais rendue. Le dépouillement des traces --- ce que le serveur a rendu, appel par appel, et ce que le modèle en a retenu --- ramène ces 37 exécutions à trois défauts. Un quatrième joue au niveau de la sous-technique. Pour chacun : un exemple, ce que le serveur a répondu, ce que le modèle a retenu, et la cause.

\emph{Défaut 1 --- le journal porte deux comportements ; le modèle en retient un, la clé attendait l'autre (13 exécutions, toutes Enterprise).} L'échantillon S016 le montre en entier. Le journal : \texttt{sudo -n cat /etc/shadow}. Cette commande fait deux choses à la fois : elle élève les privilèges (\texttt{sudo}) et elle lit le fichier des empreintes de mots de passe (\texttt{cat /etc/shadow}). Le modèle a lancé deux recherches, une par comportement. Le serveur a répondu juste aux deux : « sudo abuse privilege escalation elevate » rend \texttt{T1548.003} en tête de liste --- la réponse attendue --- et « read /etc/shadow password file dump » rend \texttt{T1003.008} en tête. Le modèle avait donc les deux bonnes fiches sous les yeux. Il a retenu l'effet, la lecture des empreintes, et sa justification l'assume : le \texttt{sudo} n'est pour lui « que le moyen » d'y parvenir. La clé attendait le moyen. À la troisième exécution, il a même ouvert les deux fiches avant de trancher, dans le même sens. Le serveur n'a donc rien manqué ; il ne peut pas non plus trancher, puisque chacune des deux réponses est exacte pour son comportement. Le choix entre les deux est une convention, et cette convention n'est écrite nulle part, ni dans la consigne ni dans la clé. La preuve que c'est bien une convention : le même défaut joue dans les deux sens à l'intérieur de la campagne. Sur S035, la clé attend l'interpréteur qui exécute la commande (\texttt{T1059}) et le modèle répond ce que la commande fait (\texttt{T1057}, la découverte de processus) ; sur S037, la clé attend cette fois le contenu obscurci que la commande transporte (\texttt{T1027}) et le modèle répond l'interpréteur (\texttt{T1059}). Les six sous-techniques niées de la campagne (S014, S035, S036) sont un dégât dérivé de ce même défaut : la technique retenue étant déjà une autre, la sous-technique attendue tombe avec elle.

\emph{Défaut 2 --- en ICS, la requête ne sépare pas des techniques voisines (11 exécutions).} Les onze autres pertes de sélection sont toutes industrielles, et sept tombent sur le même identifiant. L'échantillon S005 : le journal montre des lectures Modbus répétées (\texttt{READ\_COILS}, \texttt{READ\_DISCRETE\_INPUTS}) vers un automate ; la clé attend \texttt{T0868}, la détection du mode de marche. Ce que le serveur a rendu : \texttt{T0868} figurait dans les listes --- septième identifiant de la première recherche, troisième de la seconde --- mais \texttt{T0801}, la lecture d'état générale, sortait en tête, parce que la requête décrivait l'acte de lire (« read coils modbus process state monitor ») et non ce que cette lecture cherche à savoir. Ce que le modèle a retenu : il a ouvert les fiches de \texttt{T0801} et \texttt{T0877}, une fois celle de \texttt{T0861}, jamais celle de \texttt{T0868} ; et la fiche de \texttt{T0801} cite précisément les codes fonction \emph{Read} de Modbus, ce qui a suffi à le convaincre. La cause : ces techniques de collecte se distinguent par la finalité de la lecture, pas par le protocole ; une requête écrite dans les mots du protocole les rend indiscernables, le score met la plus générale en tête, et le modèle ne vérifie que la tête de liste. S002 (identification de points, \texttt{T0861}) et S004 (collecte automatisée, \texttt{T0802}) tombent de la même façon, toujours sur \texttt{T0801} --- sept exécutions en tout. Les quatre dernières (S006, S009) confondent, côté écriture cette fois, le message non autorisé \texttt{T1692} avec la modification qu'il produit (\texttt{T0836}, \texttt{T0835}).

\emph{Défaut 3 --- la requête omet l'artefact qui distingue, et la bonne fiche ne remonte jamais (13 exécutions).} Ici l'erreur précède le choix : aucun outil n'a rendu la bonne réponse, et aucune exécution n'a été rattrapée par la mémoire du modèle --- la ligne « jamais fournie, modèle juste » est à zéro aux trois niveaux. L'échantillon S028, à ses trois exécutions : le journal porte \texttt{sudo chown root:root /tmp/cible}, la clé attend \texttt{T1222}, la modification des permissions. Ce que le modèle a demandé : une seule recherche par exécution, toujours sur le même thème --- « sudo elevate privileges root » ---, jamais le mot \texttt{chown}, jamais les permissions. Ce que le serveur a rendu : la famille \texttt{T1548}, en toute logique, et jamais \texttt{T1222}, qui ne pouvait pas remonter puisque rien dans la requête ne le désignait. Le modèle a ouvert la fiche de \texttt{T1548.003}, l'a trouvée conforme à \texttt{sudo}, et a conclu, sans reformuler ni relancer. La cause tient en une phrase : l'artefact le plus discriminant du fragment, \texttt{chown}, n'est entré dans aucune requête. On notera que S028 porte la même paire moyen-effet que S016, avec l'attente inverse --- la clé y retient l'effet --- ce qui confirme le constat du défaut 1. La correction, elle, relève de la méthode d'interrogation, pas du moteur : inclure l'artefact discriminant et reformuler une fois avant de conclure, ce que la procédure écrite de la section~\ref{sec:ch2-3-6} imposera.

\emph{Défaut 4 --- la sous-technique en trop (11 exécutions, au niveau sous-technique).} Le modèle ajoute une sous-technique là où la clé n'en attend aucune. Trois exécutions seulement le font sur la bonne technique --- S000, aux trois exécutions --- les huit autres l'ajoutent à une technique déjà fausse, où elle ne change plus rien au compte. Le cas S000 est le plus instructif : le serveur avait rendu la parente \texttt{T0846} et la sous-technique \texttt{T0846.001}, sa note invitait à retenir la parente si l'évidence ne distingue pas la sous-technique (section~\ref{sec:ch2-2-3}), le modèle a ouvert les deux fiches, jugé que le balayage de ports distinguait bien \texttt{T0846.001}, et l'a affirmée. Tout est vérifié, rien n'est inventé ; le désaccord porte sur le niveau attendu, pas sur l'identifiant.

\textbf{Le compte, posé simplement.} Au niveau technique : 70,6~\% d'exécutions justes ; 19,0 points perdus au choix, défauts 1 et 2 ; 10,3 points perdus à la recherche, défaut 3. La part du choix est la plus grande. C'est elle qu'une méthode écrite peut viser --- ouvrir les fiches sœurs avant de trancher, reformuler quand la liste est pauvre, suivre une convention déclarée quand un fragment porte deux comportements --- et c'est l'objet de la section~\ref{sec:ch2-3-6}.

Le modèle a appelé un outil à chaque exécution, 2,9 fois en moyenne sur 3,4 tours ; le plafond de huit tours n'a jamais été approché, le maximum observé étant six. Il n'a jamais inventé d'outil ni donné d'argument refusé, et les 163 recherches portaient toutes la bonne matrice --- 121 exécutions sur 126 ont émis au moins une recherche, les cinq autres sont allées directement à la lecture d'une fiche. Le serveur n'a signalé aucune homonymie, précisément parce que la matrice était jointe à chaque recherche ; aucun identifiant soumis n'a été déclaré introuvable. Huit appels ont déclenché le garde-fou des identifiants révoqués, tous sur T0855 : la redirection vers T1692.001 corrige entièrement S007, et remplace sur S008 et S009 un identifiant mort par un identifiant vivant --- sans être pour autant le bon. Trois outils sur les dix-sept exposés ont été employés : \emph{get\_technique} 184 fois, \emph{search\_techniques} 163 fois, \emph{get\_tactic} 22 fois. Le serveur lui-même ne pèse presque rien dans le temps de réponse : environ 274 millisecondes par exécution, sur une médiane de 14,1 secondes --- le coût vient des tours de conversation, pas de l'outil.

Le taux d'ancrage de 100 \% valide la lecture des tableaux précédents : aucune réponse n'a été produite de mémoire, l'écart mesure donc bien l'apport du serveur.

\begin{longtable}[]{@{}
  >{\raggedright\arraybackslash}p{(\columnwidth - 14\tabcolsep) * \real{0.0730}}
  >{\raggedright\arraybackslash}p{(\columnwidth - 14\tabcolsep) * \real{0.1300}}
  >{\raggedright\arraybackslash}p{(\columnwidth - 14\tabcolsep) * \real{0.1300}}
  >{\raggedright\arraybackslash}p{(\columnwidth - 14\tabcolsep) * \real{0.1360}}
  >{\raggedright\arraybackslash}p{(\columnwidth - 14\tabcolsep) * \real{0.1530}}
  >{\raggedright\arraybackslash}p{(\columnwidth - 14\tabcolsep) * \real{0.1300}}
  >{\raggedright\arraybackslash}p{(\columnwidth - 14\tabcolsep) * \real{0.1240}}
  >{\raggedright\arraybackslash}p{(\columnwidth - 14\tabcolsep) * \real{0.1240}}@{}}
\toprule\noalign{}
\begin{minipage}[b]{\linewidth}\raggedright
\textbf{MCP}
\end{minipage} & \begin{minipage}[b]{\linewidth}\raggedright
\textbf{Stabilité}
\end{minipage} & \begin{minipage}[b]{\linewidth}\raggedright
\textbf{p50}
\end{minipage} & \begin{minipage}[b]{\linewidth}\raggedright
\textbf{p95}
\end{minipage} & \begin{minipage}[b]{\linewidth}\raggedright
\textbf{Jetons/exéc.}
\end{minipage} & \begin{minipage}[b]{\linewidth}\raggedright
\textbf{Sur-spéc.}
\end{minipage} & \begin{minipage}[b]{\linewidth}\raggedright
\textbf{Efficacité}
\end{minipage} & \begin{minipage}[b]{\linewidth}\raggedright
\textbf{Argent}
\end{minipage} \\
\midrule\noalign{}
\endhead
\bottomrule\noalign{}
\endlastfoot
non & 76,2 \% & 4 282 ms & 6 050 ms & 6 572 & 10,0 \% & 0,086 & 4,63 \$ \\
oui & 71,4 \% & 14 146 ms & 22 554 ms & 47 441 & 18,3 \% & 0,015 & 31,32 \$ \\
\end{longtable}
\legendetableau{Fiabilité, vitesse et coût.}{tab:ch2-fiabilite}


\textbf{Ce que disent les chiffres du coût.} Tout augmente, sauf la stabilité. Le modèle consomme \textbf{7,2 fois plus de jetons} (6 572 contre 47 441), répond en \textbf{3,3 fois plus de temps} (4,3 contre 14,1 secondes), et la campagne coûte \textbf{6,8 fois plus cher} (4,63 contre 31,32 \$). Le pire cas ordinaire suit la même pente : le 95\up{e} centile passe de 6,1 à 22,6 secondes. Rapporté au gain, l'ancrage achète 0,135 point de F1 technique au prix de 40 869 jetons par exécution : l'efficacité tombe de 0,086 à 0,015. La stabilité recule de 4,8 points, la boucle ouvrant plusieurs chemins vers une même conclusion, et la sur-spécification double presque (10,0~\% à 18,3~\%) --- c'est le défaut 4 décrit plus haut. D'où vient cette consommation, la figure~\ref{fig:ch2-echanges} le montre échange par échange.

\begin{figure}[htbp]
\centering
\figuretikz{fig-2-9-echanges-execution}
\caption{Les quatorze échanges de la première exécution de S000, numérotés dans l'ordre. Quatre appels au modèle (échanges 1, 5, 9 et 13), trois allers-retours vers le serveur. L'entrée de chaque appel au modèle contient l'intégralité des échanges précédents.}\label{fig:ch2-echanges}
\end{figure}

\textbf{D'où viennent ces jetons : une exécution décortiquée.} La figure~\ref{fig:ch2-echanges} numérote les quatorze échanges de la première exécution de S000 --- celle dont la section~\ref{sec:ch2-2-3} suivait déjà la première recherche. Une exécution avec serveur n'est pas un appel : c'est une suite d'appels au modèle, et chaque appel lui renvoie tout ce qui précède.

Échange 1 : n8n envoie au modèle la consigne, les dix-huit descriptions d'outils et le journal à analyser. Cette première entrée pèse 6 842 jetons : 2 309 pour la consigne et le journal --- c'est, à l'identique, l'appel unique de la campagne sans serveur --- et 4 533 pour le bloc d'outils, un poids fixe d'une exécution à l'autre. Échange 2 : le modèle ne répond pas encore ; il demande \texttt{search\_techniques}. Échanges 3 et 4 : n8n transmet la demande au serveur, qui répond en 2 millisecondes : 4 587 octets, dix candidats, \texttt{T0846.001} en tête. Échange 5 : n8n rappelle le modèle en lui renvoyant l'échange 1 en entier, sa propre demande, et le résultat de l'outil --- le modèle relit donc ce qu'il a déjà lu. Échanges 6 à 8 : il demande la fiche de \texttt{T0846}, que le serveur rend (2 987 octets). Échanges 9 à 12 : même chose pour \texttt{T0846.001} (3 139 octets). Échange 13 : le quatrième appel au modèle porte à lui seul 11 526 jetons d'entrée, la conversation entière. Échange 14 : le modèle rend son triplet, \texttt{TA0102} / \texttt{T0846} / \texttt{T0846.001}, en 218 jetons.

Le compte est alors simple. Les entrées des quatre appels s'additionnent : 37 035 jetons facturés en entrée, pour 540 jetons de sortie. Les 6 842 jetons du premier envoi ont été refacturés quatre fois. Le serveur, lui, a calculé 2 millisecondes en tout --- 83 millisecondes aller-retour compris --- et rendu 10 713 octets. Les 13,2 secondes de latence et l'essentiel de la facture viennent donc du modèle, qui relit une conversation de plus en plus longue, pas du serveur, qui répond. Ce mécanisme, répété sur 126 exécutions de 3,4 tours en moyenne, produit les 47 441 jetons du tableau~\ref{tab:ch2-fiabilite}. C'est aussi lui que le cache de prompt de la section~\ref{sec:ch2-3-6} viendra facturer autrement : les jetons relus ne disparaîtront pas, ils changeront de prix.

\textbf{Ce que la campagne démontre.} L'apport du serveur est réel, et il est mesuré : les trois F1 dépassent la référence, et le gain est le plus fort exactement là où le modèle était le plus faible --- la matrice industrielle passe de 0,262 à 0,524. Les identifiants périmés disparaissent entièrement, parce que le serveur ne rend que des identifiants vivants. Aucun identifiant inventé ni hors matrice n'apparaît. Et chaque réponse s'appuie sur ce que le serveur a rendu : rien n'a été répondu de mémoire. Ce ne sont pas des promesses d'architecture, ce sont des compteurs, vérifiés exécution par exécution. Reste à savoir d'où viennent les erreurs subsistantes : c'est l'objet du tableau~\ref{tab:ch2-serveur-modele}.

\subsection{Optimisation du mapping}\label{sec:ch2-3-6}

Le serveur a mis le référentiel à portée du modèle ; il ne lui a jamais dit comment s'en servir. Le tableau~\ref{tab:ch2-serveur-modele} l'établit sans ambiguïté : au niveau technique, la bonne réponse figurait dans ce que les outils avaient rendu pour 89,7 \% des exécutions, et n'a été retenue que dans 70,6 \%. Deux ajouts sont évalués ici, indépendants du serveur et indépendants l'un de l'autre : une procédure écrite, qui vise la qualité du mapping, et une facturation différente du contexte, qui vise son coût.

\textbf{Le fichier de compétences.} Un \emph{skill} est un fichier Markdown déposé à un emplacement connu du workflow --- ici \texttt{skills/mitre-attack-mapping/SKILL.md} --- dont l'en-tête déclare un nom et une description, et dont le corps décrit une marche à suivre. Il n'ajoute aucune connaissance : on n'y trouve pas une seule correspondance entre un comportement et un identifiant ATT\&CK, et rien qui remplacerait ce que le serveur rend. Il fixe une méthode, et rien d'autre. Sa règle d'ouverture énonce le principe de la phase 2 sous forme d'obligation --- la base fait foi, pas la mémoire --- et n'admet que trois issues : un mapping vérifié contre la base, un mapping multiple documenté lorsque l'observation porte réellement plusieurs techniques, ou une abstention motivée. Un identifiant qui n'a pas été vu dans la sortie d'un outil pendant la session n'est jamais une issue valide.

Neuf étapes en découlent. Le modèle commence par inventorier les outils dont il dispose, relever la version de la base et repérer les conventions que l'hôte impose ; à défaut d'outil ATT\&CK, un mode dégradé lui interdit de simuler une vérification. Il isole ensuite \emph{un} comportement atomique et en extrait le fragment exact --- ligne de commande, fonction protocolaire, chemin --- contre lequel tout le reste se jugera. Il fixe la matrice \emph{avant} de chercher, puisque la base ne connaît du contexte que ce qu'on lui passe. Il formule alors sa requête en mots-clés anglais du vocabulaire ATT\&CK, trois au minimum, en y incluant l'artefact le plus discriminant. Il consulte enfin la fiche complète des deux ou trois premiers candidats, jamais du seul premier, et c'est là que se trouve la règle la plus caractéristique du fichier : \textbf{écarter le candidat classé premier exige de nommer, dans la justification, l'élément du fragment qui l'exclut}. Sans un tel élément, le premier candidat vérifié est retenu. Suivent un ordre de départage explicite --- convention de l'hôte, spécificité de l'artefact, fonction exacte avant famille, co-occurrence déclarée, technique parente par défaut ---, six contrôles obligatoires avant de répondre (existence, matrice, tactique rattachée, sous-technique rattachée, absence de révocation, et l'identifiant vu dans une sortie d'outil), et un protocole d'abstention qui interdit de combler un vide par l'identifiant « le plus proche ».

Le fichier déclare lui-même ce qu'il ne garantit pas, et la distinction porte sur la suite de cette section : il rend la \emph{procédure} déterministe et auditable, non le \emph{résultat} infaillible. Sur une observation qui porte à la fois un moyen et son effet, la bonne réponse dépend d'une convention d'étiquetage qui appartient à l'hôte ; le skill détecte ce cas et le déclare, il ne le devine pas.

\textbf{Le cache de prompt.} Le second ajout ne touche pas au raisonnement. Une exécution de la phase 2 n'est pas un appel mais une conversation de 3,4 tours en moyenne, et chaque tour renvoie au modèle l'intégralité de ce qui précède : la consigne, les dix-sept descriptions d'outils, le journal soumis et tous les résultats déjà reçus. Mesuré sur nos campagnes, \textbf{81 \% des jetons d'entrée sont ainsi du renvoi à l'identique}. Le cache de prompt les fait reconnaître par le fournisseur comme un préfixe déjà vu et les facture 0,50 \$ le million au lieu de 5 \$, l'inscription initiale coûtant en contrepartie 1,25 fois le tarif d'entrée. Deux points de coupe sont posés : un explicite en fin de bloc système --- l'ordre de mise en cache étant outils, puis système, puis messages, il couvre aussi les descriptions d'outils --- et le cache automatique de tête pour la conversation qui grandit. Le point décisif pour la lecture des résultats est que le modèle reçoit \emph{exactement} le même contexte : le cache n'a aucun effet sur la génération, il en change seulement le prix.

\textbf{Ce que la phase 3 doit montrer.} Les deux ajouts ne se jugent donc pas au même endroit. Du skill, on attend qu'il déplace les erreurs restantes plutôt qu'il ne les masque : relever le plafond du serveur, c'est-à-dire diminuer la part des exécutions où la bonne réponse n'a jamais été rendue, sans dégrader ni l'ancrage ni la stabilité. Du cache, on attend une baisse de la facture à qualité constante --- et l'égalité des F1 est ici une condition de validité, non un résultat décevant : si les scores bougeaient nettement, il faudrait suspecter que les deux campagnes n'ont pas reçu le même contexte.

\textbf{Comment nous procédons.} Deux campagnes s'ajoutent aux deux précédentes, sur le même corpus de 42 échantillons, trois exécutions chacun, le même harnais, le même serveur --- version 2.2.1, base ATT\&CK 19.2 --- et le même modèle. La première ajoute le skill : son corps est injecté intégralement dans le prompt système, et l'alinéa ajouté en phase 2 est \emph{vidé}, de sorte que le gain soit attribuable à la procédure et non à la superposition de deux consignes. Chaque enregistrement porte le nom du fichier, son mode de chargement et son empreinte, si bien qu'une campagne se rattache à une version exacte de la méthode. La seconde campagne est la première rejouée avec un seul drapeau changé, celui du cache : c'est une réplication, ce qui en fait aussi une mesure de la variation ordinaire du modèle d'une campagne à l'autre. Les indicateurs sont ceux de la section~\ref{sec:ch2-3-3}, sans ajout, augmentés de deux compteurs : les jetons lus dans le cache et les jetons qui y sont inscrits.

\textbf{Une réserve de mesure à déclarer.} L'indicateur d'efficacité --- F1 technique par millier de jetons --- cesse d'être comparable dès qu'un cache est actif. L'interface ne compte en jetons d'entrée que ce qui n'a pas été servi par le cache : lue sur cette seule colonne, la campagne avec cache afficherait 835 jetons par exécution au lieu de 69 471, et son efficacité serait multipliée par cinquante-huit sans que rien n'ait changé. Le tableau~\ref{tab:ch2-cout-phase3} rapporte donc la consommation réelle --- entrée, sortie, cache lu et cache inscrit additionnés --- ainsi que le montant facturé, et l'efficacité en est retirée.

\begin{figure}[htbp]
\centering
\figuretikz{fig-2-8-workflow-skill-cache}
\caption{Le workflow de la phase 3 : celui de la figure~\ref{fig:workflow-mcp}, repris tel quel, avec pour seul ajout les deux points de coupe du cache. Le skill ne change pas la forme du workflow : il entre par le prompt système, au moment de préparer l'appel.}\label{fig:workflow-skill-cache}
\end{figure}

La figure~\ref{fig:workflow-skill-cache} situe le cache dans le harnais, repris tel quel de la phase 2 ; le skill, lui, n'y ajoute aucun nœud visible, puisque son corps entre par le prompt système au moment de préparer l'appel. L'annexe~\ref{ann:workflow-skill-cache} donne le détail nœud par nœud, ainsi que le texte intégral du fichier de compétences.

\textbf{La qualité du mapping.} Le tableau~\ref{tab:ch2-qualite-phase3} reprend les colonnes du tableau~\ref{tab:ch2-strict} et y ajoute les deux campagnes de la phase 3, en base stricte. Les deux premières lignes sont celles déjà commentées, rappelées pour la lecture.

\begin{longtable}[]{@{}
  >{\raggedright\arraybackslash}p{(\columnwidth - 18\tabcolsep) * \real{0.1800}}
  >{\raggedright\arraybackslash}p{(\columnwidth - 18\tabcolsep) * \real{0.0560}}
  >{\raggedright\arraybackslash}p{(\columnwidth - 18\tabcolsep) * \real{0.0790}}
  >{\raggedright\arraybackslash}p{(\columnwidth - 18\tabcolsep) * \real{0.0930}}
  >{\raggedright\arraybackslash}p{(\columnwidth - 18\tabcolsep) * \real{0.0930}}
  >{\raggedright\arraybackslash}p{(\columnwidth - 18\tabcolsep) * \real{0.0890}}
  >{\raggedright\arraybackslash}p{(\columnwidth - 18\tabcolsep) * \real{0.0930}}
  >{\raggedright\arraybackslash}p{(\columnwidth - 18\tabcolsep) * \real{0.0930}}
  >{\raggedright\arraybackslash}p{(\columnwidth - 18\tabcolsep) * \real{0.0860}}
  >{\raggedright\arraybackslash}p{(\columnwidth - 18\tabcolsep) * \real{0.0860}}@{}}
\toprule\noalign{}
\begin{minipage}[b]{\linewidth}\raggedright
\textbf{Dispositif}
\end{minipage} & \begin{minipage}[b]{\linewidth}\raggedright
\textbf{MM}
\end{minipage} & \begin{minipage}[b]{\linewidth}\raggedright
\textbf{Hall.}
\end{minipage} & \begin{minipage}[b]{\linewidth}\raggedright
\textbf{F1 tech.}
\end{minipage} & \begin{minipage}[b]{\linewidth}\raggedright
\textbf{F1 tact.}
\end{minipage} & \begin{minipage}[b]{\linewidth}\raggedright
\textbf{F1 s.-t.}
\end{minipage} & \begin{minipage}[b]{\linewidth}\raggedright
\textbf{Ent.}
\end{minipage} & \begin{minipage}[b]{\linewidth}\raggedright
\textbf{ICS}
\end{minipage} & \begin{minipage}[b]{\linewidth}\raggedright
\textbf{Obsol.}
\end{minipage} & \begin{minipage}[b]{\linewidth}\raggedright
\textbf{Abst.}
\end{minipage} \\
\midrule\noalign{}
\endhead
\bottomrule\noalign{}
\endlastfoot
Sans serveur & 0 & 0,0 \% & 0,563 & 0,587 & 0,636 & 0,714 & 0,262 & 1,3 \% & 0,0 \% \\
MCP & 0 & 0,0 \% & 0,706 & 0,802 & 0,758 & 0,798 & 0,524 & 0,0 \% & 0,0 \% \\
MCP + skill & 0 & 0,0 \% & 0,738 & 0,873 & 0,758 & 0,857 & 0,500 & 0,0 \% & 0,0 \% \\
MCP + skill + cache & 0 & 0,0 \% & 0,730 & 0,865 & 0,758 & 0,845 & 0,500 & 0,0 \% & 0,0 \% \\
\end{longtable}
\legendetableau{Qualité du mapping aux quatre dispositifs, base stricte. Ent. et ICS : F1 technique ventilé par matrice. 126 exécutions par ligne, 66 au niveau sous-technique.}{tab:ch2-qualite-phase3}


\textbf{La consommation et le temps de traitement.} Le tableau~\ref{tab:ch2-cout-phase3} rapporte les mêmes campagnes du point de vue de l'exploitation. La colonne \emph{Jetons} donne la consommation réelle par exécution, cache compris ; la colonne \emph{dont cache} isole la part servie par le préfixe déjà vu.

\begin{longtable}[]{@{}
  >{\raggedright\arraybackslash}p{(\columnwidth - 14\tabcolsep) * \real{0.1800}}
  >{\raggedright\arraybackslash}p{(\columnwidth - 14\tabcolsep) * \real{0.1050}}
  >{\raggedright\arraybackslash}p{(\columnwidth - 14\tabcolsep) * \real{0.1050}}
  >{\raggedright\arraybackslash}p{(\columnwidth - 14\tabcolsep) * \real{0.1050}}
  >{\raggedright\arraybackslash}p{(\columnwidth - 14\tabcolsep) * \real{0.1250}}
  >{\raggedright\arraybackslash}p{(\columnwidth - 14\tabcolsep) * \real{0.1150}}
  >{\raggedright\arraybackslash}p{(\columnwidth - 14\tabcolsep) * \real{0.1150}}
  >{\raggedright\arraybackslash}p{(\columnwidth - 14\tabcolsep) * \real{0.1300}}@{}}
\toprule\noalign{}
\begin{minipage}[b]{\linewidth}\raggedright
\textbf{Dispositif}
\end{minipage} & \begin{minipage}[b]{\linewidth}\raggedright
\textbf{Stab.}
\end{minipage} & \begin{minipage}[b]{\linewidth}\raggedright
\textbf{p50}
\end{minipage} & \begin{minipage}[b]{\linewidth}\raggedright
\textbf{p95}
\end{minipage} & \begin{minipage}[b]{\linewidth}\raggedright
\textbf{Jetons}
\end{minipage} & \begin{minipage}[b]{\linewidth}\raggedright
\textbf{dont cache}
\end{minipage} & \begin{minipage}[b]{\linewidth}\raggedright
\textbf{Sur-spéc.}
\end{minipage} & \begin{minipage}[b]{\linewidth}\raggedright
\textbf{Argent}
\end{minipage} \\
\midrule\noalign{}
\endhead
\bottomrule\noalign{}
\endlastfoot
Sans serveur & 76,2 \% & 4 282 ms & 6 050 ms & 6 572 & --- & 10,0 \% & 4,63 \$ \\
MCP & 71,4 \% & 14 146 ms & 22 554 ms & 47 441 & --- & 18,3 \% & 31,32 \$ \\
MCP + skill & 85,7 \% & 19 487 ms & 30 110 ms & 69 234 & --- & 16,7 \% & 45,71 \$ \\
MCP + skill + cache & 83,3 \% & 20 652 ms & 30 937 ms & 69 471 & 55 777 & 15,0 \% & 16,23 \$ \\
\end{longtable}
\legendetableau{Consommation et temps de traitement aux quatre dispositifs. Jetons : consommation réelle moyenne par exécution --- entrée, sortie, cache lu et cache inscrit additionnés. Argent : coût des 126 exécutions, chaque composante facturée à son tarif du 16 septembre 2026.}{tab:ch2-cout-phase3}


\textbf{Le constat.} Les deux tableaux se lisent ensemble, et ils racontent deux histoires différentes.

Le skill améliore la qualité, modestement au niveau technique et nettement ailleurs. Le F1 technique passe de 0,706 à 0,738, soit quatre exécutions justes de plus sur 126 ; le F1 tactique gagne au contraire 0,071, et Enterprise 0,059. La stabilité est le gain le plus net : elle remonte de 71,4 \% à 85,7 \%, c'est-à-dire qu'elle dépasse enfin celle du modèle sans serveur (76,2 \%). La procédure écrite corrige donc précisément ce que la boucle d'outils avait dégradé en phase 2, où plusieurs chemins menaient à des conclusions différentes. La sur-spécification recule elle aussi, de 18,3 \% à 16,7 \%, puis à 15,0 \%, sans revenir au niveau du modèle seul. Les quatre garanties de la phase 2 tiennent : aucune hallucination, aucun identifiant hors matrice, aucun identifiant périmé, et 100 \% de réponses ancrées dans une sortie d'outil.

Le croisement serveur contre modèle explique où le skill agit. Le plafond du serveur au niveau technique passe de 89,7 \% à 94,4 \% des exécutions, et la part des exécutions où la bonne réponse n'a jamais été rendue tombe de 10,3 \% à 5,6 \% : \textbf{le skill divise par deux les échecs de recherche}. La perte de sélection, elle, ne diminue pas --- elle progresse même de 19,0 à 21,4 points. L'explication est mécanique : mieux conduite, la recherche ramène davantage de matière, 28,7 candidats distincts par exécution contre 22,1, et une liste plus riche est une liste plus difficile à trancher. Le modèle mobilise aussi plus d'outils, six au lieu de trois --- l'inventaire de l'étape 0 l'y conduit --- pour 3,7 appels par exécution contre 2,9. La procédure a donc réglé la moitié gauche du problème, celle de la recherche, et laissé intacte la moitié droite, celle du choix final.

Une seule régression apparaît, et elle mérite d'être dite : \textbf{le F1 technique sur la matrice industrielle passe de 0,524 à 0,500}. L'écart vaut une exécution sur 42 ; il n'est pas significatif, mais il suffit à établir que le doublement du score ICS constaté en phase 2 vient de l'ancrage et non de la procédure. Le skill n'a rien apporté à la matrice industrielle.

Le cache, lui, tient exactement sa promesse, et la tient sans rien coûter à la qualité. Les jetons réellement consommés ne bougent pas --- 69 471 contre 69 234, une variation de trois pour mille : le cache ne raccourcit rien, il change le tarif de ce qui est relu. La facture, elle, passe de 45,71 \$ à 16,23 \$, soit \textbf{une baisse de 64 \%}, et de 31,32 \$ pour le serveur seul, soit une baisse de 48 \% pour un dispositif plus précis. Rapporté au résultat utile, le coût d'un mapping technique juste tombe de 0,491 \$ à 0,176 \$. Et les scores confirment la neutralité attendue : 92 techniques justes contre 93, 109 tactiques contre 110, la sous-technique inchangée à 0,758 --- une exécution d'écart aux deux premiers niveaux, c'est-à-dire la variation ordinaire du modèle d'une campagne à l'autre. Le contrôle de validité est donc satisfait.

Une nuance doit accompagner ce chiffre. Le dispositif complet reste 3,5 fois plus cher que la réponse de mémoire (16,23 \$ contre 4,63 \$) et 4,8 fois plus lent (20,7 contre 4,3 secondes en médiane). Le cache abaisse le prix de l'ancrage ; il ne le ramène pas à celui d'une réponse non vérifiée, et ce n'est pas son objet. Le poste dominant de ce qui reste n'est d'ailleurs plus la lecture de cache, qui ne pèse que 22 \% de la facture, mais son inscription, qui en représente 62 \% : la conversation grandissant à chaque tour, chaque tour inscrit un nouveau préfixe. C'est là que se trouve la marge suivante.

\textbf{Les erreurs qui persistent.} Un F1 technique de 0,730 laisse 34 exécutions fausses sur 126. Le dépouillement des traces les range dans quatre familles, dont trois se concentrent sur la sélection finale. Chacune est illustrée par un cas du corpus.

\emph{Le rabattement sur la technique voisine la plus générale.} Prolongement direct du défaut 2 de la phase 2, c'est la famille la plus nombreuse et elle est entièrement industrielle : dix exécutions concluent à \texttt{T0801 --- Monitor Process State} là où quatre échantillons différents attendaient \texttt{T0861}, \texttt{T0868}, \texttt{T0802} ou \texttt{T0814}. Sur S002, le journal Zeek montre douze requêtes Modbus \texttt{READ\_HOLDING\_REGISTERS} consécutives vers un automate, et la vérité de terrain attend \texttt{T0861 --- Point \& Tag Identification}. Le modèle a pourtant fait tout ce que la procédure demande : il a cherché avec la matrice ICS, consulté les fiches de trois candidats --- \texttt{T0801}, \texttt{T0877} et \texttt{T0861} ---, énuméré la tactique \texttt{TA0100} pour voir les techniques sœurs, et nommé l'élément qui exclut chacun des écartés, comme l'exige la règle d'écartement de l'étape 4. Sa justification écarte \texttt{T0861} au motif que la technique vise l'identification des tags et non la lecture de leurs valeurs. Le raisonnement est défendable ; il est simplement en désaccord avec la clé. L'étape 5.3 du skill vise précisément ce piège --- « fonction exacte avant famille » --- et ne suffit pas à le lever : la distinction entre lire un état et cartographier des points ne se décide pas sur le code fonction protocolaire, seul discriminant que le journal offre.

\emph{La co-occurrence d'un moyen et de son effet.} Le défaut 1 de la phase 2 est intact : trois exécutions sur S016 concluent à \texttt{T1003 --- OS Credential Dumping} là où la clé attend \texttt{T1548 --- Abuse Elevation Control Mechanism}. Le fragment est un unique \texttt{EXECVE} : \texttt{sudo -n cat /etc/shadow}. Il porte réellement les deux techniques --- l'élévation de privilège est le moyen, la lecture des empreintes est l'effet --- et le serveur a rendu les deux dès les deux premières recherches. Le modèle a retenu l'effet, en écrivant dans sa justification que le \texttt{sudo} n'était ici « que le moyen d'obtenir l'accès root pour cette lecture ». Le cas est exactement celui que l'étape 5.4 du skill décrit et que sa dernière section déclare non garanti : en l'absence de convention d'étiquetage écrite, aucune procédure ne peut deviner laquelle des deux techniques l'évaluateur a retenue. L'échantillon S028 en donne le miroir : sur \texttt{sudo chown root:root /tmp/cible}, la clé attend cette fois l'effet, \texttt{T1222 --- File and Directory Permissions Modification}, et le modèle retient le moyen, \texttt{T1548.003}. Les deux réponses suivent la même logique et sont fausses toutes les deux, parce que la convention change d'un échantillon à l'autre. Une part de l'erreur résiduelle est donc une propriété de la clé, non du dispositif --- constat qui appelle, pour une campagne ultérieure, d'écrire la règle de co-occurrence dans les conventions de l'hôte, ce que l'étape 0 du skill prévoit explicitement.

\emph{L'échec de recherche.} Le défaut 3 de la phase 2 recule de treize à sept exécutions sans s'éteindre : l'erreur est en amont du choix, la bonne réponse n'a jamais été rendue par un outil. Le cas le plus net est S028, où le fragment porte à la fois \texttt{sudo} et \texttt{chown root:root}. Les deux requêtes émises --- « sudo privilege escalation elevated execution », puis « unix shell command interpreter execution » --- décrivent le contexte d'exécution et jamais l'action sur les permissions. \texttt{chown}, l'artefact le plus discriminant du fragment, n'apparaît dans aucune requête ; \texttt{T1222} ne pouvait donc pas remonter. L'étape 3 du skill impose pourtant d'inclure l'artefact le plus discriminant : la consigne existe, elle n'a pas été appliquée. C'est une limite de l'exécution de la procédure, pas de sa rédaction.

\emph{La sur-spécification.} Neuf exécutions ajoutent une sous-technique là où la clé n'en attend aucune. Sur S000, le modèle rend la bonne technique, \texttt{T0846 --- Remote System Discovery}, puis y ajoute \texttt{T0846.001 --- Port Scan} au motif, exact, que le journal montre un balayage du port 502 sur une plage d'adresses. La sous-technique existe, elle est rattachée à la bonne parente, et le serveur l'a rendue : les six contrôles de l'étape 6 passent tous. La note du serveur --- visible dans l'exemple de la section~\ref{sec:ch2-2-3}, qui suit précisément cette requête --- l'invitait à ne retenir la sous-technique que si l'évidence la distingue ; le modèle a jugé que c'était le cas, et son argument se défend. Le désaccord porte sur le niveau attendu, non sur l'identifiant. Le phénomène s'est d'ailleurs aggravé avec l'ancrage --- 10,0 \% sans serveur, 18,3 \% avec --- pour une raison compréhensible : mieux informé des sous-techniques existantes, le modèle en affirme davantage. Le skill le contient sans le supprimer, par sa règle du niveau parent par défaut : 16,7 \%, puis 15,0 \%.

Ces quatre familles ont un point commun : aucune n'est une défaillance de l'ancrage. Dans les quatre cas, l'identifiant rendu existe, appartient à la bonne matrice, n'est pas révoqué et provient d'une sortie d'outil. Ce qui reste à traiter n'est plus la fiabilité du référentiel, c'est la décision finale --- et, pour une part qu'il faut assumer, la précision de la clé d'évaluation elle-même.

\section{Les avantages de l'approche}\label{sec:ch2-4}

Quatre bénéfices découlent de cette architecture. Les trois premiers portent sur la qualité du mapping, le quatrième sur le déploiement.

\textbf{Le versionnage du référentiel.} C'est l'avantage principal. ATT\&CK n'est pas figé : des techniques apparaissent, d'autres fusionnent ou sont retirées à chaque révision. Un modèle qui répond de mémoire restitue un état du référentiel qu'il ne sait pas dater, et l'analyste n'a aucun moyen de le vérifier. Notre serveur sert au contraire une version nommée, choisie et épinglée, et la joint à chaque réponse : une analyse est datée, et rejouable à l'identique des mois plus tard. Les identifiants retirés ne sont pas perdus pour autant, puisqu'un identifiant révoqué est redirigé vers son remplaçant vivant --- \texttt{T1086}, par exemple, conduit à \texttt{T1059.001}. Un rapport ancien reste donc exploitable sans travail de conversion. Sur notre campagne d'évaluation, le taux d'identifiants obsolètes est passé de \textbf{1,3 \%} à \textbf{0 \%}.

\textbf{La personnalisation interne.} ATT\&CK décrit des comportements observés à l'échelle mondiale ; une organisation en observe qui lui sont propres, ou que le référentiel public ne couvre qu'avec une granularité insuffisante. Le référentiel interne permet de les décrire sous des identifiants maison, alimentés par les rapports de red team et les productions de renseignement internes, chaque fiche demeurant rattachée à une technique ATT\&CK (section~\ref{sec:ch2-2-2}). Les deux niveaux cohabitent sans se contredire : la précision locale est ajoutée là où le référentiel public s'arrête, et le vocabulaire commun reste intact, de sorte qu'un interlocuteur externe comprend toujours de quoi il est question.

\textbf{La précision maîtrisée.} Le serveur ne rend pas le modèle plus savant : il rend ses réponses vérifiables. Tout identifiant renvoyé provient de la base chargée, ce qui retire au modèle la possibilité d'en composer un qui n'existe pas. Chaque candidat arrive accompagné des indices qui l'ont fait remonter, et l'absence de correspondance constitue une réponse explicite plutôt qu'un identifiant forcé. L'effet est mesurable : sur la même campagne, le score F1 du mapping est passé de \textbf{0,571} à \textbf{0,706}.

\textbf{La polyvalence.} Le serveur ignore qui l'appelle : il parle MCP, et tout client qui parle MCP peut s'y brancher. Changer d'usage revient donc à changer de client, sans rien modifier au serveur. Le premier cas d'utilisation en donne la mesure : une chaîne conversationnelle bâtie dans n8n (figure~\ref{fig:ch2-n8n}). La question de l'analyste entre par un nœud de discussion et parvient à un agent porté par Claude Sonnet, auquel sont rattachés une mémoire de session et le serveur MCP MITRE déclaré comme outil. S'y ajoute une intelligence procédurale : un fichier de compétences décrit à l'agent la marche à suivre, c'est-à-dire la méthode de mapping elle-même. Chaque appel est ensuite assemblé puis transmis par HTTP au service de journalisation, avant que la réponse ne revienne dans le fil, sourcée et datée. L'ensemble se comporte comme un assistant d'analyste : il dispose du référentiel public et des fiches internes, et répond à toute question posée dans la conversation.

\begin{figure}[htbp]
\centering
\figurefichier{figures/figure-2-7_chaine-conversationnelle-n8n}
\caption{Cas d'utilisation : chaîne conversationnelle n8n. L'agent reçoit la question, dispose du serveur MCP MITRE comme outil, et chaque appel est journalisé avant le retour de la réponse.}\label{fig:ch2-n8n}
\end{figure}

\section{Conclusion}\label{sec:ch2-5}

Ce chapitre a construit un serveur MCP pour MITRE ATT\&CK, décrit son moteur de mapping, puis mesuré ce qu'il apporte au \emph{mapping} de journaux dans des conditions proches de celles d'un centre opérationnel. Trois campagnes ont été rejouées sur le même corpus de 42 échantillons, le même harnais et le même modèle, contre une ligne de base établie au chapitre~\ref{chap:preparation} : le serveur seul, le serveur augmenté d'une procédure écrite, puis le même dispositif avec cache de prompt.

\textbf{Ce qui est établi.} L'ancrage tient ses quatre garanties, exécution par exécution et non par construction : aucun identifiant inventé, aucun identifiant pris hors de la matrice de l'échantillon, aucun identifiant périmé --- contre 1,3 \% pour le modèle seul --- et 100 \% de réponses appuyées sur un identifiant que le serveur avait effectivement rendu. Aucune réponse n'a été produite de mémoire alors que les outils étaient disponibles. La qualité suit : le F1 technique passe de 0,563 à 0,730, le F1 tactique de 0,587 à 0,865, et le gain est le plus fort exactement là où le modèle était le plus faible, la matrice industrielle passant de 0,262 à 0,500. La procédure écrite y ajoute ce que la boucle d'outils avait coûté --- la stabilité remonte à 83,3 \%, au-dessus du modèle sans serveur --- et divise par deux les échecs de recherche. Le cache, enfin, ramène la facture de 45,71 \$ à 16,23 \$ sans toucher à un seul des scores, ce qui rend le coût de l'ancrage supportable là où il était rédhibitoire.

Ces résultats comblent une partie de ce que la revue laissait ouvert. La supériorité de l'accès déterministe sur le \emph{mapping} ATT\&CK n'était pas mesurée en conditions de SOC (\textbf{QR3}) : elle l'est ici, sur des indicateurs qui séparent ce que le serveur fournit de ce que le modèle retient. L'évaluation sur de la télémétrie industrielle restait un angle mort (\textbf{QR1}) : les 42 échantillons sont générés en laboratoire, donc garantis non contaminés, et un tiers d'entre eux relève de la zone OT. Et la propriété que le RAG ne pouvait pas offrir --- un résultat daté, reproductible, rattaché à une version nommée du référentiel (\textbf{QR2}) --- est ici une conséquence directe de l'architecture, non un effet de bord.

\textbf{Ce qui reste ouvert.} Trois limites sont assumées. La première est la sélection : 21,4 points séparent encore ce que le serveur fournit de ce que le modèle retient au niveau technique, et c'est désormais le premier poste d'erreur. La deuxième est la matrice industrielle, qui plafonne à 0,500 et que la procédure écrite n'a pas fait progresser --- les techniques voisines d'une même tactique ICS se distinguent par des fonctions que le journal ne discrimine pas toujours. La troisième tient à la clé d'évaluation elle-même : sur les observations qui portent un moyen et son effet, la convention d'étiquetage n'est pas écrite et varie d'un échantillon à l'autre, de sorte qu'une part de l'erreur résiduelle n'est pas imputable au dispositif. À quoi s'ajoute la portée du banc d'essai --- un seul modèle, 42 échantillons, trois exécutions --- et un coût d'exploitation qui demeure 3,5 fois celui d'une réponse non vérifiée, pour un temps de traitement 4,8 fois plus long.

\textbf{Ce qui suit.} Le serveur a été conçu en lecture seule, journalisé et versionné, parce que la revue avait montré que le MCP déplace le risque du texte vers l'action (\textbf{QR4}). Ces propriétés ont été décrites, non éprouvées. Le chapitre suivant les met à l'épreuve : il prolonge l'approche par un flux multi-MCP d'audit de sécurité, fondé sur le référentiel ATLAS et placé sous contrainte réglementaire canadienne.


% Décommenter au fur et à mesure que les chapitres sont ajoutés :
% \input{chapitres/chapitre3.tex}

%============================ CONCLUSION GÉNÉRALE ============================

\chapterwithtoc{Conclusion générale}
\label{chap:conclusion}

Cet essai partait d'un constat : le goulot d'étranglement d'un centre
d'opérations de sécurité n'est pas la détection, mais la qualification de ce qui
est détecté. Il a conçu et mesuré un écosystème ancrant un modèle de langage sur
MITRE ATT\&CK par le \emph{Model Context Protocol}.

\textbf{Les modèles seuls ne suffisent pas, et l'écart se creuse côté
industriel.} Sur 42 échantillons produits en laboratoire, hors contamination, dix
modèles courants ont été évalués en 1\,260 appels. Aucun ne dépasse 0,603 de F1
technique en base stricte et les cinq premiers se tiennent en moins de 0,04 : il
n'y a pas de vainqueur net. La ventilation par matrice est plus instructive que
le classement --- le meilleur passe de 0,679 en Enterprise à 0,452 en ICS. Les
vingt-trois identifiants fabriqués du panel sont tous des tactiques, et tous sur
des journaux industriels : l'angle mort signalé par la revue est mesuré, et il
est d'abord un défaut de repérage entre univers.

\textbf{L'ancrage corrige ce que le modèle ne peut pas corriger seul.} À corpus
et conditions identiques, le modèle retenu progresse de 0,563 à 0,706 en
technique et de 0,587 à 0,802 en tactique. Le gain n'est pas uniforme : la
matrice industrielle double presque, de 0,262 à 0,524, ce qui écarte
l'hypothèse d'un effet de consigne. Surtout, le taux d'identifiants obsolètes ne
diminue pas, il disparaît --- une classe entière d'erreurs supprimée, qu'aucun
raisonnement ne rattrape puisque rien dans le journal ne signale que le
référentiel a bougé depuis l'entraînement.

\textbf{Le mérite revient à la conception du serveur}, non au simple accès à une
source : une version nommée et jointe à chaque réponse, la résolution des
révocations, une granularité d'outils bornée, une trace d'usage vérifiable. Le
risque, lui, est déplacé et non supprimé : le serveur devient un canal
d'instruction autant qu'une source de faits, et l'audit de cette surface reste à
mener.

Trois limites bornent ces résultats : un seul modèle évalué avec ancrage, un
corpus de 42 échantillons, et un effort de raisonnement laissé aux valeurs par
défaut faute d'échelle commune entre éditeurs. Reste l'enseignement principal :
les modèles les plus capables sont les plus exposés, parce que l'attaque exploite
l'aptitude même qui fait leur valeur. La réponse n'est pas un modèle plus grand,
mais un cadre qui ancre, borne et journalise.


%================================= ANNEXES ===================================
\appendix
%% ===================================================================
%% ANNEXE A --- Génération d'attaque : l'exemple T0846
%% Exécution, capture Zeek et construction du jeu de données.
%% ===================================================================
\chapter{Génération d'attaque : l'exemple T0846}
\label{ann:t0846}

Cette annexe montre, sur l'exemple de \textbf{T0846 — Remote System Discovery}
(MITRE ATT\&CK for ICS, tactique \emph{Discovery}), comment une attaque réelle est
\emph{exécutée}, \emph{capturée} par Zeek, puis transformée en un \emph{échantillon
étiqueté}. Le procédé est identique pour toutes les techniques.

% --------------------------- Schéma d'ensemble ---------------------------
\paragraph{Chaîne de bout en bout.} Le harnais \lstinline|atk| \textbf{réside dans
kali} : pour chaque attaque, il \textbf{enregistre} deux horodatages UTC ($t_0$
avant, $t_1$ après) et l'étiquette (ID + fenêtre) dans la \textbf{vérité-terrain}.
Pendant ce temps, Zeek journalise le trafic \emph{sans connaître} la technique. Le
script recolle le tout : il découpe la télémétrie sur $[t_0,t_1]$ en un échantillon
\lstinline|S###.jsonl| (\textbf{sans étiquette}) et place la réponse dans
\lstinline|verite_terrain.jsonl|. Cette séparation télémétrie\,/\,étiquette rend
l'évaluation honnête.

\begin{figure}[h]\centering
\begin{tikzpicture}[node distance=5mm and 11mm]
\node[atk]  (kali) {Attaquant — kali \texttt{.90.6}\\ \scriptsize contient le \textbf{harnais \lstinline|atk|}};
\node[tar,right=of kali] (cibles) {Cibles Modbus\\ \scriptsize simulation \texttt{.95.10–.15:502}};
\node[zeek,below=of cibles] (zeek) {Sonde \textbf{Zeek}\\ \scriptsize \texttt{conn.log}, \texttt{modbus.log}};
\node[gt,below=of kali] (trace) {\texttt{trace\_attaques\_zeek.jsonl}\\ \scriptsize fenêtre $t_0,t_1$ + ID};
\node[script,right=24mm of cibles] (script) {\lstinline|preparer_dataset.py|\\ \scriptsize valide l'ID MITRE,\\[-1pt] découpe par fenêtre};
\node[sortie,above=7mm of script] (samples) {\texttt{S000.jsonl}\\ \scriptsize \textbf{sans} étiquette};
\node[gt,below=7mm of script] (verite) {\texttt{verite\_terrain}\\ \scriptsize la réponse};
\draw[afleche,cAtk] (kali) -- node[aetiq,above]{balayage} (cibles);
\draw[afleche,cZeek] (cibles) -- node[aetiq,right]{capture} (zeek);
\draw[aflechep,cGT] (kali) -- node[aetiq,right]{enregistre} (trace);
\draw[afleche,cZeek] (zeek.east) to[out=0,in=180] (script.west);
\draw[afleche,cGT]  (trace.east) to[out=0,in=205] (script.west);
\draw[afleche,cOut] (script) -- (samples);
\draw[afleche,cGT]  (script) -- (verite);
\begin{scope}[on background layer]
\node[draw=cTar!35,dashed,rounded corners,inner sep=7pt,fill=cBg,
  fit=(kali)(cibles)(zeek)(trace),
  label={[font=\scriptsize\bfseries,text=cTar]above:Lab TO (génération + capture)}]{};
\node[draw=cScript!35,dashed,rounded corners,inner sep=7pt,fill=cScript!3,
  fit=(script)(samples)(verite),
  label={[font=\scriptsize\bfseries,text=cScript]above:Poste (dataset)}]{};
\end{scope}
\end{tikzpicture}
\caption{Chaîne complète. Le harnais \texttt{atk} réside dans kali et
\emph{enregistre} la fenêtre dans la trace ; Zeek capture le balayage. Télémétrie
(bleu) et étiquette (vert) ne se rejoignent que dans le script.}
\end{figure}

% --------------------------- Script d'attaque + séquence ---------------------------
\paragraph{Script de l'attaque (depuis \texttt{kali}).} On balaie
\texttt{192.168.95.0/24} et on teste, pour chaque adresse, si le port TCP
\textbf{502} (Modbus) est ouvert. Un port ouvert se termine proprement
(\texttt{conn\_state=SF}, serveur Modbus) ; un port fermé est rejeté
(\texttt{conn\_state=REJ}). Les six serveurs de la simulation
(\texttt{.10}–\texttt{.15}) répondent \texttt{SF}, \texttt{.45} répond \texttt{REJ}.

\begin{lstlisting}[style=bash]
atk T0846 - kali docker exec -i kali python3 -c '
import socket
for i in range(1,255):                       # teste .1 a .254
    s = socket.socket(); s.settimeout(0.2)
    try: s.connect(("192.168.95.%d"%i,502)); print("MODBUS ouvert: .%d"%i)
    except: pass
    finally: s.close()'
sleep 30                                      # 30 s : isole la fenetre de l'attaque
\end{lstlisting}

\begin{figure}[h]\centering
\begin{tikzpicture}[x=1cm,y=1cm]
\def\xK{0}\def\xC{5.6}\def\xZ{10.6}
\node[atk]  (K) at (\xK,0) {kali \texttt{.90.6}};
\node[tar]  (C) at (\xC,0) {cibles \texttt{.95.x:502}};
\node[zeek] (Z) at (\xZ,0) {Zeek $\rightarrow$ \texttt{conn.log}};
\foreach \x/\c in {\xK/cAtk,\xC/cTar,\xZ/cZeek}
  \draw[\c!45,thick,dashed] (\x,-0.45) -- (\x,-4.75);
\draw[gt,fill=cGT!8] (\xK-1.0,-0.72) rectangle (\xZ+1.0,-1.05);
\node[font=\scriptsize\itshape,text=cGT!60!black] at (\xK,-0.885) {$t_0$ noté par \lstinline|atk|};
\node[font=\scriptsize\bfseries,text=cGT!60!black,anchor=west] at (\xK-1.0,-1.38)
  {Cas A — \texttt{.11} (serveur présent)};
\draw[afleche,cAtk] (\xK,-1.85) -- node[aetiq,above]{SYN vers \texttt{.11:502}} (\xC,-1.85);
\draw[afleche,cTar] (\xC,-2.30) -- node[aetiq,above]{SYN/ACK, ACK} (\xK,-2.30);
\draw[afleche,cZeek,densely dashed] (\xC,-2.68) -- node[aetiq,above]{journalise} (\xZ,-2.68);
\node[gt,font=\scriptsize,anchor=west] at (\xZ+0.12,-2.68){\texttt{SF} = OUVERT};
\node[font=\scriptsize\bfseries,text=cAtk,anchor=west] at (\xK-1.0,-3.10)
  {Cas B — \texttt{.45} (aucun serveur)};
\draw[afleche,cAtk] (\xK,-3.57) -- node[aetiq,above]{SYN vers \texttt{.45:502}} (\xC,-3.57);
\draw[afleche,cTar] (\xC,-4.02) -- node[aetiq,above]{RST (rejet)} (\xK,-4.02);
\draw[afleche,cZeek,densely dashed] (\xC,-4.40) -- node[aetiq,above]{journalise} (\xZ,-4.40);
\node[atk,font=\scriptsize,anchor=west] at (\xZ+0.12,-4.40){\texttt{REJ} = FERMÉ};
\node[font=\scriptsize\itshape,text=cComment] at (\xC,-5.0) {\dots\ répété pour $i=1\dots254$, puis $t_1$ noté \dots};
\end{tikzpicture}
\caption{Exécution de T0846 (diagramme de séquence). Chaque tentative laisse une
ligne dans \texttt{conn.log} ; l'ensemble forme l'empreinte d'un balayage.}
\end{figure}

% --------------------------- Capture Zeek ---------------------------
\paragraph{Zeek : capture que l'attaque.} La sonde zeek partage le même espace de
noms réseau de la simulation EWS 192.168.95.x : elle voit donc aussi la boucle de fond permanent de l'automate PLC$\leftrightarrow$contrôleur, très bruyant. Elle est \textbf{préconfigurée pour
ne pas capter ce bruit} : comme la source d'attaque est connue d'avance
(\texttt{kali 192.168.90.6}), on limite la détection de l'interface du segment
\texttt{.95}, puis applique un filtre BPF qui ne garde que les échanges dont un
attaquant est une extrémité, et charge les analyseurs Modbus/DNP3 via le script ci-dessous :

\begin{lstlisting}[style=bash]
# entrypoint.sh de la sonde Zeek (reglage obtenu par essais successifs)
IFACE=$(ip -o -4 addr show | awk '/192\.168\.95\./ {print $2; exit}')   # interface du segment .95
FILTRE="host 192.168.90.6 or host 192.168.95.5"     # ne garder que le trafic d'un attaquant
exec zeek -C -i "$IFACE" -f "$FILTRE" icsnpp-modbus icsnpp-dnp3 LogAscii::use_json=T
\end{lstlisting}

\noindent Ce filtre est un \textbf{succès} : seuls les signaux d'attaque sont
enregistrés. La capture reste continue, mais comme seul le trafic de l'attaquant
subsiste, la télémétrie utile se réduit à la fenêtre $[t_0,t_1]$ de chaque tir, et
la pause de 30\,s rend les fenêtres disjointes. Pour T0846, l'échange se limite à
des poignées de main TCP : le balayage se lit dans \texttt{conn.log}.

% --------------------------- Harnais + trace ---------------------------
\paragraph{Harnais \texttt{atk} (dans kali) et vérité-terrain.} Pour chaque tir, le
harnais note $t_0$, exécute la commande d'attaque, note $t_1$, puis
\textbf{enregistre une ligne} dans \texttt{trace\_attaques\_zeek.jsonl} :
l'\textbf{identifiant}, la source et la fenêtre (le \emph{nom} sera résolu plus tard
par le script).

\begin{lstlisting}[style=json]
// trace_attaques_zeek.jsonl (aetiquette : ID + fenetre, jamais melangee a la telemetrie)
{ "technique":"T0846", "sub_technique":"-", "src":"192.168.90.6",
  "t0":"2026-08-31T18:32:57.7Z", "t1":"2026-08-31T18:33:27.8Z", "rc":0 }
\end{lstlisting}

% --------------------------- Construction + sorties ---------------------------
\paragraph{Construction du dataset.} Le script \lstinline|preparer_dataset.py| est
\textbf{préconfiguré pour reconnaître la zone et les journaux à analyser} : il
détecte la zone TO via la trace, puis localise les logs Zeek connus
(\texttt{modbus.log}, \texttt{conn.log}, \texttt{modbus\_read\_device\_identification.log}) ;
pour ce balayage réseau, c'est \texttt{conn.log} qu'il analyse. Il télécharge puis
met en \textbf{cache} le catalogue ATT\&CK for ICS et \textbf{vérifie chaque ID} de
la trace : s'il est révoqué, il le \textbf{met à jour} par son successeur ; sinon il
en résout le nom (\emph{Remote System Discovery}) et la tactique (\emph{Discovery},
\texttt{TA0102}). Il récupère alors \textbf{exactement} les lignes dont l'horodatage
tombe dans $[t_0,t_1]$, les renumérote (\texttt{S000-0001}\dots), et écrit
\textbf{deux fichiers distincts} : l'\textbf{échantillon} (les logs Zeek, sans
étiquette) et la \textbf{réponse} \texttt{verite\_terrain.jsonl}, utilisée à
l'évaluation.

\noindent\textbf{Contenu complet de l'échantillon} \texttt{S000.jsonl} — les 7
événements capturés (télémétrie brute, \textbf{sans étiquette}) : six connexions
acceptées vers les serveurs Modbus \texttt{.10}–\texttt{.15}
(\texttt{conn\_state=SF}) et une connexion rejetée vers \texttt{.45}
(\texttt{conn\_state=REJ}).

\begin{lstlisting}[style=jsonfull]
{"id_journal": "S000-0001", "zone": "TO", "matrice": "ics", "source": "zeek", "horodatage": "2026-08-31T18:32:57.787604Z", "contenu": {"ts": 1788201177.787604, "uid": "CU6F1C32Nw2Gfb6gPe", "id.orig_h": "192.168.90.6", "id.orig_p": 43820, "id.resp_h": "192.168.95.10", "id.resp_p": 502, "proto": "tcp", "duration": 0.00045490264892578125, "orig_bytes": 0, "resp_bytes": 0, "conn_state": "SF", "local_orig": true, "local_resp": true, "missed_bytes": 0, "history": "ShAFf", "orig_pkts": 4, "orig_ip_bytes": 216, "resp_pkts": 2, "resp_ip_bytes": 112, "ip_proto": 6}}
{"id_journal": "S000-0002", "zone": "TO", "matrice": "ics", "source": "zeek", "horodatage": "2026-08-31T18:32:57.787725Z", "contenu": {"ts": 1788201177.787725, "uid": "CqKcsvtwkRGmAP6Rj", "id.orig_h": "192.168.90.6", "id.orig_p": 59570, "id.resp_h": "192.168.95.11", "id.resp_p": 502, "proto": "tcp", "duration": 0.0014679431915283203, "orig_bytes": 0, "resp_bytes": 0, "conn_state": "SF", "local_orig": true, "local_resp": true, "missed_bytes": 0, "history": "ShAFaf", "orig_pkts": 4, "orig_ip_bytes": 216, "resp_pkts": 3, "resp_ip_bytes": 164, "ip_proto": 6}}
{"id_journal": "S000-0003", "zone": "TO", "matrice": "ics", "source": "zeek", "horodatage": "2026-08-31T18:32:57.788217Z", "contenu": {"ts": 1788201177.788217, "uid": "CHsh9K3bivHb8At539", "id.orig_h": "192.168.90.6", "id.orig_p": 58610, "id.resp_h": "192.168.95.12", "id.resp_p": 502, "proto": "tcp", "duration": 0.0012409687042236328, "orig_bytes": 0, "resp_bytes": 0, "conn_state": "SF", "local_orig": true, "local_resp": true, "missed_bytes": 0, "history": "ShAFaf", "orig_pkts": 4, "orig_ip_bytes": 216, "resp_pkts": 3, "resp_ip_bytes": 164, "ip_proto": 6}}
{"id_journal": "S000-0004", "zone": "TO", "matrice": "ics", "source": "zeek", "horodatage": "2026-08-31T18:32:57.788312Z", "contenu": {"ts": 1788201177.788312, "uid": "CJPqby2g4XNTldykp", "id.orig_h": "192.168.90.6", "id.orig_p": 43274, "id.resp_h": "192.168.95.13", "id.resp_p": 502, "proto": "tcp", "duration": 0.0016789436340332031, "orig_bytes": 0, "resp_bytes": 0, "conn_state": "SF", "local_orig": true, "local_resp": true, "missed_bytes": 0, "history": "ShAFaf", "orig_pkts": 4, "orig_ip_bytes": 216, "resp_pkts": 3, "resp_ip_bytes": 164, "ip_proto": 6}}
{"id_journal": "S000-0005", "zone": "TO", "matrice": "ics", "source": "zeek", "horodatage": "2026-08-31T18:32:57.788433Z", "contenu": {"ts": 1788201177.788433, "uid": "CoZPeJ226jvmdQY7D9", "id.orig_h": "192.168.90.6", "id.orig_p": 52516, "id.resp_h": "192.168.95.14", "id.resp_p": 502, "proto": "tcp", "duration": 0.0007989406585693359, "orig_bytes": 0, "resp_bytes": 0, "conn_state": "SF", "local_orig": true, "local_resp": true, "missed_bytes": 0, "history": "ShAFaf", "orig_pkts": 4, "orig_ip_bytes": 216, "resp_pkts": 3, "resp_ip_bytes": 164, "ip_proto": 6}}
{"id_journal": "S000-0006", "zone": "TO", "matrice": "ics", "source": "zeek", "horodatage": "2026-08-31T18:32:57.788548Z", "contenu": {"ts": 1788201177.788548, "uid": "Cq0Tw03Oplz2OHLQwc", "id.orig_h": "192.168.90.6", "id.orig_p": 37310, "id.resp_h": "192.168.95.15", "id.resp_p": 502, "proto": "tcp", "duration": 0.0012180805206298828, "orig_bytes": 0, "resp_bytes": 0, "conn_state": "SF", "local_orig": true, "local_resp": true, "missed_bytes": 0, "history": "ShAFaf", "orig_pkts": 4, "orig_ip_bytes": 216, "resp_pkts": 3, "resp_ip_bytes": 164, "ip_proto": 6}}
{"id_journal": "S000-0007", "zone": "TO", "matrice": "ics", "source": "zeek", "horodatage": "2026-08-31T18:33:03.614345Z", "contenu": {"ts": 1788201183.614345, "uid": "CyfTL91PCRlC93PXxe", "id.orig_h": "192.168.90.6", "id.orig_p": 38592, "id.resp_h": "192.168.95.45", "id.resp_p": 502, "proto": "tcp", "duration": 1.0013580322265625e-05, "orig_bytes": 0, "resp_bytes": 0, "conn_state": "REJ", "local_orig": true, "local_resp": true, "missed_bytes": 0, "history": "Sr", "orig_pkts": 1, "orig_ip_bytes": 60, "resp_pkts": 1, "resp_ip_bytes": 40, "ip_proto": 6}}
\end{lstlisting}

\noindent Et la réponse associée, dans \texttt{verite\_terrain.jsonl} (utilisée à l'évaluation) :

\begin{lstlisting}[style=json]
{ "echantillon":"S000", "fichier":"S000.jsonl", "zone":"TO", "source":"zeek", "matrice":"ics",
  "tactique_attendue":{"id":"TA0102","nom":"Discovery"},
  "technique_attendue":{"id":"T0846","nom":"Remote System Discovery"},
  "sous_technique_attendue":{"id":"aucune","nom":"aucune"} }
\end{lstlisting}


%============================== BIBLIOGRAPHIE ================================
\cleardoublepage
\addcontentsline{toc}{chapter}{\bibname}
% Style de l'UQAC s'il est présent, repli sinon.
\IfFileExists{IEEEtranSN_francais.bst}
  {\bibliographystyle{IEEEtranSN_francais}}
  {\bibliographystyle{unsrtnat}}
\bibliography{revue_litterature,references}

\end{document}
